% Results and Analysis, organised by claim (C1-C6 from page 2).
\section{Results and Analysis}\label{sec:results}

All numbers come from the saved study \texttt{20260929-policy-study-v2}, which an
independent script recomputes with zero mismatches over 99 runs. Within a stream every
arm replays a byte-identical score trace, so differences come from decision rules alone.
Durations and delays are in sample-index units.

\subsection{The ratchet and its cost (C1, C2; H1)}\label{sec:results-ratchet}
\claim{1}{\Sproved\,\Smeasured}{We prove the ratchet and observe it in the saved
traces.}
The normal-only rolling threshold records 0 increases in 143 (valve1) and 140 (valve2)
replay steps. It falls from $1.309$ to $0.258$ on valve1 and from $1.126$ to $0.406$ on
valve2, below the median replay score (Figure~\ref{fig:ratchet}), exactly as
Proposition~\ref{prop:ratchet} predicts.

\claim{2}{\Smeasured}{We measure what the ratchet costs.}
Rolling (normal-only) detects the same single event as the fixed threshold, but with
$16.2\times$ the alert windows on valve1 (97 vs 6) and $3.9\times$ on valve2 (70 vs 18)
(Table~\ref{tab:baseline}). The extra alerts come from the falling threshold, not from
extra anomalies.

\begin{figure}[tbp]
  \centering
  \resizebox{0.9\linewidth}{!}{\providecolor{scblack}{RGB}{0,0,0}%
\providecolor{scblue}{RGB}{31,78,150}%
\providecolor{scred}{RGB}{170,40,30}%
\providecolor{scgreen}{RGB}{20,110,60}%
\providecolor{scgray}{RGB}{120,120,120}%
\providecolor{scevent}{RGB}{60,60,60}%
\begin{tikzpicture}[>=stealth,line join=round,line cap=round]
  \begin{scope}[xshift=0.00cm]
  \fill[scevent,opacity=0.14] (0.000,0) rectangle (3.994,4.300);
  \draw[scgray,solid,line width=0.5pt] (0.000,0.848) -- (0.040,0.765) -- (0.080,1.440) -- (0.120,1.425) -- (0.161,0.820) -- (0.201,0.883) -- (0.241,0.554) -- (0.281,0.497) -- (0.321,0.000) -- (0.361,1.159) -- (0.401,1.355) -- (0.442,1.179) -- (0.482,1.299) -- (0.522,1.356) -- (0.562,1.068) -- (0.602,0.000) -- (0.642,0.482) -- (0.682,0.915) -- (0.723,0.991) -- (0.763,0.926) -- (0.803,0.821) -- (0.843,1.407) -- (0.883,1.062) -- (0.923,0.870) -- (0.963,0.887) -- (1.004,0.094) -- (1.044,0.771) -- (1.084,0.894) -- (1.124,0.275) -- (1.164,0.554) -- (1.204,1.098) -- (1.244,0.000) -- (1.285,0.140) -- (1.325,0.954) -- (1.365,1.237) -- (1.405,0.955) -- (1.445,1.132) -- (1.485,0.928) -- (1.525,0.000) -- (1.565,0.659) -- (1.606,4.099) -- (1.646,0.746) -- (1.686,1.094) -- (1.726,1.048) -- (1.766,0.996) -- (1.806,0.869) -- (1.846,0.994) -- (1.887,1.458) -- (1.927,4.106) -- (1.967,0.637) -- (2.007,1.186) -- (2.047,1.118) -- (2.087,0.949) -- (2.127,0.000) -- (2.168,0.180) -- (2.208,0.647) -- (2.248,1.011) -- (2.288,1.216) -- (2.328,1.161) -- (2.368,0.925) -- (2.408,0.920) -- (2.449,1.271) -- (2.489,1.138) -- (2.529,0.830) -- (2.569,0.983) -- (2.609,0.000) -- (2.649,0.768) -- (2.689,1.203) -- (2.730,0.985) -- (2.770,0.880) -- (2.810,1.358) -- (2.850,0.904) -- (2.890,0.691) -- (2.930,0.925) -- (2.970,0.174) -- (3.011,0.859) -- (3.051,0.856) -- (3.091,0.780) -- (3.131,1.020) -- (3.171,1.278) -- (3.211,0.781) -- (3.251,1.260) -- (3.292,1.191) -- (3.332,1.107) -- (3.372,0.322) -- (3.412,0.000) -- (3.452,0.927) -- (3.492,1.057) -- (3.532,0.000) -- (3.573,0.926) -- (3.613,0.965) -- (3.653,1.187) -- (3.693,1.000) -- (3.733,0.973) -- (3.773,1.106) -- (3.813,0.714) -- (3.854,0.348) -- (3.894,1.093) -- (3.934,0.567) -- (3.974,1.338) -- (4.014,1.104) -- (4.054,1.157) -- (4.094,1.073) -- (4.135,0.977) -- (4.175,1.261) -- (4.215,0.469) -- (4.255,0.820) -- (4.295,0.269) -- (4.335,0.991) -- (4.375,0.253) -- (4.415,0.514) -- (4.456,1.133) -- (4.496,0.000) -- (4.536,1.042) -- (4.576,0.976) -- (4.616,1.114) -- (4.656,0.286) -- (4.696,0.349) -- (4.737,0.322) -- (4.777,0.432) -- (4.817,0.963) -- (4.857,0.314) -- (4.897,1.232) -- (4.937,0.805) -- (4.977,1.088) -- (5.018,1.137) -- (5.058,0.749) -- (5.098,0.577) -- (5.138,0.496) -- (5.178,0.509) -- (5.218,1.313) -- (5.258,1.111) -- (5.299,0.454) -- (5.339,0.000) -- (5.379,0.305) -- (5.419,0.000) -- (5.459,0.877) -- (5.499,0.638) -- (5.539,1.271) -- (5.580,1.279) -- (5.620,1.147) -- (5.660,1.103) -- (5.700,1.025);
  \draw[scblack,solid,line width=0.8pt] (0,1.381) -- (5.700,1.381);
  \draw[scred,dashed,line width=0.8pt] (0.000,1.381) -- (0.040,0.848) -- (0.080,0.848) -- (0.120,0.848) -- (0.161,0.848) -- (0.201,0.848) -- (0.241,0.848) -- (0.281,0.848) -- (0.321,0.848) -- (0.361,0.848) -- (0.401,0.848) -- (0.442,0.848) -- (0.482,0.848) -- (0.522,0.848) -- (0.562,0.848) -- (0.602,0.848) -- (0.642,0.848) -- (0.682,0.848) -- (0.723,0.848) -- (0.763,0.848) -- (0.803,0.848) -- (0.843,0.848) -- (0.883,0.848) -- (0.923,0.848) -- (0.963,0.848) -- (1.004,0.848) -- (1.044,0.848) -- (1.084,0.821) -- (1.124,0.821) -- (1.164,0.821) -- (1.204,0.821) -- (1.244,0.821) -- (1.285,0.821) -- (1.325,0.821) -- (1.365,0.821) -- (1.405,0.821) -- (1.445,0.821) -- (1.485,0.821) -- (1.525,0.821) -- (1.565,0.821) -- (1.606,0.821) -- (1.646,0.821) -- (1.686,0.821) -- (1.726,0.821) -- (1.766,0.821) -- (1.806,0.821) -- (1.846,0.821) -- (1.887,0.821) -- (1.927,0.821) -- (1.967,0.821) -- (2.007,0.771) -- (2.047,0.771) -- (2.087,0.771) -- (2.127,0.771) -- (2.168,0.771) -- (2.208,0.746) -- (2.248,0.746) -- (2.288,0.746) -- (2.328,0.746) -- (2.368,0.746) -- (2.408,0.746) -- (2.449,0.746) -- (2.489,0.746) -- (2.529,0.746) -- (2.569,0.746) -- (2.609,0.746) -- (2.649,0.746) -- (2.689,0.746) -- (2.730,0.746) -- (2.770,0.746) -- (2.810,0.746) -- (2.850,0.746) -- (2.890,0.746) -- (2.930,0.746) -- (2.970,0.746) -- (3.011,0.746) -- (3.051,0.746) -- (3.091,0.746) -- (3.131,0.746) -- (3.171,0.746) -- (3.211,0.746) -- (3.251,0.746) -- (3.292,0.746) -- (3.332,0.746) -- (3.372,0.746) -- (3.412,0.746) -- (3.452,0.746) -- (3.492,0.746) -- (3.532,0.746) -- (3.573,0.691) -- (3.613,0.691) -- (3.653,0.691) -- (3.693,0.691) -- (3.733,0.691) -- (3.773,0.691) -- (3.813,0.691) -- (3.854,0.691) -- (3.894,0.691) -- (3.934,0.691) -- (3.974,0.691) -- (4.014,0.691) -- (4.054,0.691) -- (4.094,0.691) -- (4.135,0.691) -- (4.175,0.691) -- (4.215,0.691) -- (4.255,0.691) -- (4.295,0.691) -- (4.335,0.691) -- (4.375,0.691) -- (4.415,0.691) -- (4.456,0.567) -- (4.496,0.567) -- (4.536,0.567) -- (4.576,0.567) -- (4.616,0.567) -- (4.656,0.567) -- (4.696,0.567) -- (4.737,0.567) -- (4.777,0.567) -- (4.817,0.567) -- (4.857,0.567) -- (4.897,0.514) -- (4.937,0.514) -- (4.977,0.514) -- (5.018,0.514) -- (5.058,0.514) -- (5.098,0.514) -- (5.138,0.514) -- (5.178,0.514) -- (5.218,0.514) -- (5.258,0.514) -- (5.299,0.514) -- (5.339,0.514) -- (5.379,0.509) -- (5.419,0.509) -- (5.459,0.509) -- (5.499,0.509) -- (5.539,0.509) -- (5.580,0.509) -- (5.620,0.509) -- (5.660,0.509) -- (5.700,0.509);
  \draw[scblue,dotted,line width=0.8pt] (0.000,1.381) -- (0.040,0.848) -- (0.080,0.848) -- (0.120,1.440) -- (0.161,1.440) -- (0.201,1.440) -- (0.241,1.440) -- (0.281,1.440) -- (0.321,1.440) -- (0.361,1.440) -- (0.401,1.440) -- (0.442,1.440) -- (0.482,1.440) -- (0.522,1.425) -- (0.562,1.356) -- (0.602,1.356) -- (0.642,1.356) -- (0.682,1.356) -- (0.723,1.356) -- (0.763,1.356) -- (0.803,1.356) -- (0.843,1.356) -- (0.883,1.407) -- (0.923,1.407) -- (0.963,1.407) -- (1.004,1.407) -- (1.044,1.407) -- (1.084,1.407) -- (1.124,1.407) -- (1.164,1.407) -- (1.204,1.407) -- (1.244,1.407) -- (1.285,1.098) -- (1.325,1.098) -- (1.365,1.098) -- (1.405,1.237) -- (1.445,1.237) -- (1.485,1.237) -- (1.525,1.237) -- (1.565,1.237) -- (1.606,1.237) -- (1.646,4.099) -- (1.686,4.099) -- (1.726,4.099) -- (1.766,4.099) -- (1.806,4.099) -- (1.846,4.099) -- (1.887,4.099) -- (1.927,4.099) -- (1.967,4.106) -- (2.007,4.106) -- (2.047,4.106) -- (2.087,4.106) -- (2.127,4.106) -- (2.168,4.106) -- (2.208,4.106) -- (2.248,4.106) -- (2.288,4.106) -- (2.328,4.106) -- (2.368,1.216) -- (2.408,1.216) -- (2.449,1.216) -- (2.489,1.271) -- (2.529,1.271) -- (2.569,1.271) -- (2.609,1.271) -- (2.649,1.271) -- (2.689,1.271) -- (2.730,1.271) -- (2.770,1.271) -- (2.810,1.271) -- (2.850,1.358) -- (2.890,1.358) -- (2.930,1.358) -- (2.970,1.358) -- (3.011,1.358) -- (3.051,1.358) -- (3.091,1.358) -- (3.131,1.358) -- (3.171,1.358) -- (3.211,1.358) -- (3.251,1.278) -- (3.292,1.278) -- (3.332,1.278) -- (3.372,1.278) -- (3.412,1.278) -- (3.452,1.278) -- (3.492,1.278) -- (3.532,1.278) -- (3.573,1.278) -- (3.613,1.260) -- (3.653,1.260) -- (3.693,1.191) -- (3.733,1.187) -- (3.773,1.187) -- (3.813,1.187) -- (3.854,1.187) -- (3.894,1.187) -- (3.934,1.187) -- (3.974,1.187) -- (4.014,1.338) -- (4.054,1.338) -- (4.094,1.338) -- (4.135,1.338) -- (4.175,1.338) -- (4.215,1.338) -- (4.255,1.338) -- (4.295,1.338) -- (4.335,1.338) -- (4.375,1.338) -- (4.415,1.261) -- (4.456,1.261) -- (4.496,1.261) -- (4.536,1.261) -- (4.576,1.261) -- (4.616,1.133) -- (4.656,1.133) -- (4.696,1.133) -- (4.737,1.133) -- (4.777,1.133) -- (4.817,1.133) -- (4.857,1.133) -- (4.897,1.114) -- (4.937,1.232) -- (4.977,1.232) -- (5.018,1.232) -- (5.058,1.232) -- (5.098,1.232) -- (5.138,1.232) -- (5.178,1.232) -- (5.218,1.232) -- (5.258,1.313) -- (5.299,1.313) -- (5.339,1.313) -- (5.379,1.313) -- (5.419,1.313) -- (5.459,1.313) -- (5.499,1.313) -- (5.539,1.313) -- (5.580,1.313) -- (5.620,1.313) -- (5.660,1.279) -- (5.700,1.279);
  \node[font=\scriptsize,scred,anchor=north west,fill=white,fill opacity=0.9,text opacity=1,inner sep=1pt,align=left] at (0.12,4.220) {start 1.309\\end 0.258};
  \draw[scblack] (0,0) rectangle (5.700,4.300);
  \draw[scblack] (0.000,0) -- (0.000,-0.08);
  \node[font=\scriptsize,below] at (0.000,-0.08) {577};
  \draw[scblack] (2.850,0) -- (2.850,-0.08);
  \node[font=\scriptsize,below] at (2.850,-0.08) {861};
  \draw[scblack] (5.700,0) -- (5.700,-0.08);
  \node[font=\scriptsize,below] at (5.700,-0.08) {1145};
  \draw[scblack] (0,0.000) -- (-0.08,0.000);
  \node[font=\scriptsize,left] at (-0.10,0.000) {$10^{-1}$};
  \draw[scblack] (0,1.237) -- (-0.08,1.237);
  \node[font=\scriptsize,left] at (-0.10,1.237) {$10^{0}$};
  \draw[scblack] (0,2.473) -- (-0.08,2.473);
  \node[font=\scriptsize,left] at (-0.10,2.473) {$10^{1}$};
  \draw[scblack] (0,3.710) -- (-0.08,3.710);
  \node[font=\scriptsize,left] at (-0.10,3.710) {$10^{2}$};
  \node[font=\scriptsize,below=0.42cm] at (2.850,0) {sample index (window end)};
  \node[font=\scriptsize,rotate=90] at (-0.920,2.150) {replay score};
  \node[font={\scriptsize\bfseries},above] at (2.850,4.340) {valve1 (score, log$_{10}$)};
  \end{scope}
  \begin{scope}[xshift=7.20cm]
  \fill[scevent,opacity=0.14] (0.000,0) rectangle (4.008,4.300);
  \draw[scgray,solid,line width=0.5pt] (0.000,1.239) -- (0.041,0.279) -- (0.082,0.156) -- (0.123,4.122) -- (0.164,1.366) -- (0.205,1.245) -- (0.246,1.404) -- (0.287,0.981) -- (0.328,1.277) -- (0.369,0.000) -- (0.410,1.164) -- (0.451,0.476) -- (0.492,0.661) -- (0.533,1.250) -- (0.574,1.072) -- (0.615,1.190) -- (0.656,0.972) -- (0.697,1.246) -- (0.738,0.516) -- (0.779,0.952) -- (0.820,0.000) -- (0.861,1.204) -- (0.902,0.000) -- (0.943,0.601) -- (0.984,0.302) -- (1.025,1.409) -- (1.066,1.210) -- (1.107,1.369) -- (1.148,0.996) -- (1.189,1.015) -- (1.230,1.104) -- (1.271,0.720) -- (1.312,0.266) -- (1.353,0.000) -- (1.394,0.488) -- (1.435,1.153) -- (1.476,1.093) -- (1.517,4.063) -- (1.558,1.453) -- (1.599,1.075) -- (1.640,0.000) -- (1.681,0.911) -- (1.722,0.786) -- (1.763,1.128) -- (1.804,0.884) -- (1.845,0.667) -- (1.886,1.217) -- (1.927,0.701) -- (1.968,0.988) -- (2.009,0.908) -- (2.050,1.034) -- (2.091,1.404) -- (2.132,1.347) -- (2.173,1.168) -- (2.214,0.971) -- (2.255,1.105) -- (2.296,0.000) -- (2.337,0.637) -- (2.378,0.923) -- (2.419,0.821) -- (2.460,0.887) -- (2.501,0.416) -- (2.542,0.654) -- (2.583,1.084) -- (2.624,1.043) -- (2.665,1.102) -- (2.706,1.075) -- (2.747,0.912) -- (2.788,1.335) -- (2.829,0.496) -- (2.871,0.000) -- (2.912,1.025) -- (2.953,1.184) -- (2.994,0.836) -- (3.035,0.815) -- (3.076,1.135) -- (3.117,1.256) -- (3.158,1.070) -- (3.199,0.840) -- (3.240,1.269) -- (3.281,1.210) -- (3.322,1.286) -- (3.363,0.909) -- (3.404,1.283) -- (3.445,1.221) -- (3.486,1.023) -- (3.527,1.081) -- (3.568,0.400) -- (3.609,0.760) -- (3.650,1.296) -- (3.691,1.205) -- (3.732,0.748) -- (3.773,1.185) -- (3.814,1.398) -- (3.855,1.320) -- (3.896,1.186) -- (3.937,1.446) -- (3.978,0.766) -- (4.019,0.926) -- (4.060,0.000) -- (4.101,0.676) -- (4.142,0.000) -- (4.183,0.412) -- (4.224,0.204) -- (4.265,0.706) -- (4.306,0.904) -- (4.347,1.080) -- (4.388,0.899) -- (4.429,1.279) -- (4.470,1.259) -- (4.511,0.254) -- (4.552,1.060) -- (4.593,1.076) -- (4.634,0.135) -- (4.675,0.792) -- (4.716,1.035) -- (4.757,1.246) -- (4.798,1.489) -- (4.839,4.146) -- (4.880,1.168) -- (4.921,1.355) -- (4.962,0.753) -- (5.003,1.222) -- (5.044,1.118) -- (5.085,0.137) -- (5.126,0.933) -- (5.167,0.274) -- (5.208,1.465) -- (5.249,1.301) -- (5.290,0.989) -- (5.331,0.745) -- (5.372,0.530) -- (5.413,1.027) -- (5.454,0.728) -- (5.495,0.000) -- (5.536,1.227) -- (5.577,1.270) -- (5.618,1.233) -- (5.659,1.182) -- (5.700,0.734);
  \draw[scblack,solid,line width=0.8pt] (0,1.300) -- (5.700,1.300);
  \draw[scred,dashed,line width=0.8pt] (0.000,1.300) -- (0.041,1.239) -- (0.082,1.239) -- (0.123,1.239) -- (0.164,1.239) -- (0.205,1.239) -- (0.246,1.239) -- (0.287,1.239) -- (0.328,1.239) -- (0.369,1.239) -- (0.410,1.239) -- (0.451,1.239) -- (0.492,1.239) -- (0.533,1.239) -- (0.574,1.239) -- (0.615,1.239) -- (0.656,1.239) -- (0.697,1.190) -- (0.738,1.190) -- (0.779,1.190) -- (0.820,1.190) -- (0.861,1.190) -- (0.902,1.190) -- (0.943,1.190) -- (0.984,1.190) -- (1.025,1.190) -- (1.066,1.190) -- (1.107,1.190) -- (1.148,1.190) -- (1.189,1.190) -- (1.230,1.190) -- (1.271,1.104) -- (1.312,1.104) -- (1.353,1.104) -- (1.394,1.104) -- (1.435,1.104) -- (1.476,1.104) -- (1.517,1.104) -- (1.558,1.104) -- (1.599,1.104) -- (1.640,1.104) -- (1.681,1.104) -- (1.722,1.104) -- (1.763,1.104) -- (1.804,1.104) -- (1.845,1.093) -- (1.886,1.093) -- (1.927,1.093) -- (1.968,1.093) -- (2.009,1.093) -- (2.050,1.093) -- (2.091,1.075) -- (2.132,1.075) -- (2.173,1.075) -- (2.214,1.075) -- (2.255,1.034) -- (2.296,1.034) -- (2.337,1.034) -- (2.378,1.034) -- (2.419,1.034) -- (2.460,1.034) -- (2.501,1.034) -- (2.542,1.034) -- (2.583,1.034) -- (2.624,1.034) -- (2.665,1.034) -- (2.706,1.034) -- (2.747,1.034) -- (2.788,1.034) -- (2.829,1.034) -- (2.871,0.971) -- (2.912,0.923) -- (2.953,0.923) -- (2.994,0.923) -- (3.035,0.923) -- (3.076,0.923) -- (3.117,0.923) -- (3.158,0.923) -- (3.199,0.923) -- (3.240,0.912) -- (3.281,0.912) -- (3.322,0.912) -- (3.363,0.912) -- (3.404,0.912) -- (3.445,0.912) -- (3.486,0.912) -- (3.527,0.912) -- (3.568,0.912) -- (3.609,0.912) -- (3.650,0.912) -- (3.691,0.912) -- (3.732,0.912) -- (3.773,0.912) -- (3.814,0.912) -- (3.855,0.912) -- (3.896,0.912) -- (3.937,0.912) -- (3.978,0.912) -- (4.019,0.909) -- (4.060,0.909) -- (4.101,0.909) -- (4.142,0.909) -- (4.183,0.909) -- (4.224,0.909) -- (4.265,0.909) -- (4.306,0.766) -- (4.347,0.766) -- (4.388,0.766) -- (4.429,0.766) -- (4.470,0.766) -- (4.511,0.766) -- (4.552,0.766) -- (4.593,0.766) -- (4.634,0.766) -- (4.675,0.766) -- (4.716,0.766) -- (4.757,0.766) -- (4.798,0.766) -- (4.839,0.766) -- (4.880,0.766) -- (4.921,0.766) -- (4.962,0.766) -- (5.003,0.766) -- (5.044,0.766) -- (5.085,0.766) -- (5.126,0.753) -- (5.167,0.753) -- (5.208,0.753) -- (5.249,0.753) -- (5.290,0.753) -- (5.331,0.753) -- (5.372,0.753) -- (5.413,0.753) -- (5.454,0.753) -- (5.495,0.753) -- (5.536,0.753) -- (5.577,0.753) -- (5.618,0.753) -- (5.659,0.753) -- (5.700,0.753);
  \draw[scblue,dotted,line width=0.8pt] (0.000,1.300) -- (0.041,1.239) -- (0.082,1.239) -- (0.123,1.239) -- (0.164,4.122) -- (0.205,4.122) -- (0.246,4.122) -- (0.287,4.122) -- (0.328,4.122) -- (0.369,4.122) -- (0.410,4.122) -- (0.451,4.122) -- (0.492,4.122) -- (0.533,4.122) -- (0.574,1.404) -- (0.615,1.404) -- (0.656,1.404) -- (0.697,1.277) -- (0.738,1.277) -- (0.779,1.250) -- (0.820,1.250) -- (0.861,1.250) -- (0.902,1.250) -- (0.943,1.250) -- (0.984,1.246) -- (1.025,1.246) -- (1.066,1.409) -- (1.107,1.409) -- (1.148,1.409) -- (1.189,1.409) -- (1.230,1.409) -- (1.271,1.409) -- (1.312,1.409) -- (1.353,1.409) -- (1.394,1.409) -- (1.435,1.409) -- (1.476,1.369) -- (1.517,1.369) -- (1.558,4.063) -- (1.599,4.063) -- (1.640,4.063) -- (1.681,4.063) -- (1.722,4.063) -- (1.763,4.063) -- (1.804,4.063) -- (1.845,4.063) -- (1.886,4.063) -- (1.927,4.063) -- (1.968,1.453) -- (2.009,1.217) -- (2.050,1.217) -- (2.091,1.217) -- (2.132,1.404) -- (2.173,1.404) -- (2.214,1.404) -- (2.255,1.404) -- (2.296,1.404) -- (2.337,1.404) -- (2.378,1.404) -- (2.419,1.404) -- (2.460,1.404) -- (2.501,1.404) -- (2.542,1.347) -- (2.583,1.168) -- (2.624,1.105) -- (2.665,1.105) -- (2.706,1.102) -- (2.747,1.102) -- (2.788,1.102) -- (2.829,1.335) -- (2.871,1.335) -- (2.912,1.335) -- (2.953,1.335) -- (2.994,1.335) -- (3.035,1.335) -- (3.076,1.335) -- (3.117,1.335) -- (3.158,1.335) -- (3.199,1.335) -- (3.240,1.256) -- (3.281,1.269) -- (3.322,1.269) -- (3.363,1.286) -- (3.404,1.286) -- (3.445,1.286) -- (3.486,1.286) -- (3.527,1.286) -- (3.568,1.286) -- (3.609,1.286) -- (3.650,1.286) -- (3.691,1.296) -- (3.732,1.296) -- (3.773,1.296) -- (3.814,1.296) -- (3.855,1.398) -- (3.896,1.398) -- (3.937,1.398) -- (3.978,1.446) -- (4.019,1.446) -- (4.060,1.446) -- (4.101,1.446) -- (4.142,1.446) -- (4.183,1.446) -- (4.224,1.446) -- (4.265,1.446) -- (4.306,1.446) -- (4.347,1.446) -- (4.388,1.080) -- (4.429,1.080) -- (4.470,1.279) -- (4.511,1.279) -- (4.552,1.279) -- (4.593,1.279) -- (4.634,1.279) -- (4.675,1.279) -- (4.716,1.279) -- (4.757,1.279) -- (4.798,1.279) -- (4.839,1.489) -- (4.880,4.146) -- (4.921,4.146) -- (4.962,4.146) -- (5.003,4.146) -- (5.044,4.146) -- (5.085,4.146) -- (5.126,4.146) -- (5.167,4.146) -- (5.208,4.146) -- (5.249,4.146) -- (5.290,1.465) -- (5.331,1.465) -- (5.372,1.465) -- (5.413,1.465) -- (5.454,1.465) -- (5.495,1.465) -- (5.536,1.465) -- (5.577,1.465) -- (5.618,1.465) -- (5.659,1.301) -- (5.700,1.270);
  \node[font=\scriptsize,scred,anchor=north west,fill=white,fill opacity=0.9,text opacity=1,inner sep=1pt,align=left] at (0.12,4.220) {start 1.126\\end 0.406};
  \draw[scblack] (0,0) rectangle (5.700,4.300);
  \draw[scblack] (0.000,0) -- (0.000,-0.08);
  \node[font=\scriptsize,below] at (0.000,-0.08) {565};
  \draw[scblack] (2.850,0) -- (2.850,-0.08);
  \node[font=\scriptsize,below] at (2.850,-0.08) {843};
  \draw[scblack] (5.700,0) -- (5.700,-0.08);
  \node[font=\scriptsize,below] at (5.700,-0.08) {1121};
  \draw[scblack] (0,0.000) -- (-0.08,0.000);
  \node[font=\scriptsize,left] at (-0.10,0.000) {$10^{-1}$};
  \draw[scblack] (0,1.237) -- (-0.08,1.237);
  \node[font=\scriptsize,left] at (-0.10,1.237) {$10^{0}$};
  \draw[scblack] (0,2.473) -- (-0.08,2.473);
  \node[font=\scriptsize,left] at (-0.10,2.473) {$10^{1}$};
  \draw[scblack] (0,3.710) -- (-0.08,3.710);
  \node[font=\scriptsize,left] at (-0.10,3.710) {$10^{2}$};
  \node[font=\scriptsize,below=0.42cm] at (2.850,0) {sample index (window end)};
  \node[font=\scriptsize,rotate=90] at (-0.920,2.150) {replay score};
  \node[font={\scriptsize\bfseries},above] at (2.850,4.340) {valve2 (score, log$_{10}$)};
  \end{scope}
  \draw[scgray,solid,line width=0.9pt] (0.000,-1.150) -- (0.500,-1.150);
  \node[font=\scriptsize,right] at (0.550,-1.150) {replay score};
  \draw[scblack,solid,line width=0.9pt] (2.600,-1.150) -- (3.100,-1.150);
  \node[font=\scriptsize,right] at (3.150,-1.150) {fixed $\theta_0$};
  \draw[scred,dashed,line width=0.9pt] (4.700,-1.150) -- (5.200,-1.150);
  \node[font=\scriptsize,right] at (5.250,-1.150) {rolling (normal-only)};
  \draw[scblue,dotted,line width=0.9pt] (8.600,-1.150) -- (9.100,-1.150);
  \node[font=\scriptsize,right] at (9.150,-1.150) {rolling (all)};
  \fill[scevent,opacity=0.14] (11.500,-1.270) rectangle (12.000,-1.030);
  \node[font=\scriptsize,right] at (12.050,-1.150) {labelled event span};
\end{tikzpicture}%
}
  \caption{Replay score (grey, $\log_{10}$) on valve1 (left) and valve2 (right), with
  the fixed threshold $\theta_0$ (black), the normal-only rolling threshold (red dashed)
  and the admit-all rolling threshold (blue dotted); the labelled event is shaded. Start
  and end threshold values are printed inside each panel. The normal-only threshold only
  ever falls.}
  \label{fig:ratchet}
\end{figure}

\subsection{The scorer, not the policy, limits detection (C3)}\label{sec:results-scorer}
\claim{3}{\Smeasured}{We measure that no alert rule can recover what the scorer does
not separate.}
Inside the labelled event only $6/100$ (valve1) and $13/98$ (valve2) replay windows
exceed $\theta_0$, and the longest run above it is 2 windows on both streams
(Figure~\ref{fig:scorer-limit}). Persistence rules therefore fall short of their own
requirement: \texttt{k\_consecutive} ($k=3$) misses both events and \texttt{m\_of\_n}
(3/5) misses valve1. Each event also covers about 70\% of the replay decisions, so a
single alert anywhere inside it scores full recall. On SKAB, recall of $1/1$ carries
little information; workload, false episodes and the in-event alerted fraction carry it.

\begin{figure}[tbp]
  \centering
  \resizebox{0.88\linewidth}{!}{\definecolor{scblack}{RGB}{ 25, 32, 56}%
\definecolor{scblue}{RGB}{  0,114,178}%
\definecolor{scred}{RGB}{ 213, 94,  0}%
\definecolor{scgreen}{RGB}{  0,158,115}%
\definecolor{scgray}{RGB}{ 88,108,140}%
\definecolor{scevent}{RGB}{230,159,  0}%
\begin{tikzpicture}[>=stealth,line join=round,line cap=round]
  \begin{scope}[xshift=0.00cm]
  \draw[scblack,densely dotted,line width=0.9pt] (0,1.200) -- (5.900,1.200);
  \draw[scblue,dotted,line width=0.8pt] (0.069,0.025) -- (0.206,0.052) -- (0.343,0.062) -- (0.480,0.147) -- (0.617,0.151) -- (0.755,0.156) -- (0.892,0.162) -- (1.029,0.165) -- (1.166,0.167) -- (1.303,0.176) -- (1.441,0.205) -- (1.578,0.214) -- (1.715,0.220) -- (1.852,0.231) -- (1.990,0.237) -- (2.127,0.239) -- (2.264,0.268) -- (2.401,0.301) -- (2.538,0.370) -- (2.676,0.411) -- (2.813,0.422) -- (2.950,0.470) -- (3.087,0.551) -- (3.224,0.564) -- (3.362,0.565) -- (3.499,0.580) -- (3.636,0.619) -- (3.773,0.638) -- (3.910,0.677) -- (4.048,0.696) -- (4.185,0.716) -- (4.322,0.717) -- (4.459,0.726) -- (4.597,0.730) -- (4.734,0.756) -- (4.871,0.761) -- (5.008,0.775) -- (5.145,0.790) -- (5.283,0.908) -- (5.420,0.960) -- (5.557,0.978) -- (5.694,0.991) -- (5.831,1.057);
  \draw[scred,solid,line width=1.0pt] (0.030,0.003) -- (0.088,0.015) -- (0.148,0.026) -- (0.207,0.028) -- (0.266,0.032) -- (0.325,0.051) -- (0.384,0.068) -- (0.443,0.080) -- (0.502,0.109) -- (0.560,0.119) -- (0.620,0.127) -- (0.679,0.128) -- (0.738,0.153) -- (0.797,0.167) -- (0.856,0.175) -- (0.915,0.225) -- (0.974,0.231) -- (1.032,0.257) -- (1.092,0.257) -- (1.151,0.264) -- (1.210,0.300) -- (1.268,0.306) -- (1.328,0.313) -- (1.387,0.332) -- (1.446,0.346) -- (1.505,0.368) -- (1.564,0.381) -- (1.623,0.383) -- (1.681,0.386) -- (1.740,0.392) -- (1.800,0.393) -- (1.859,0.422) -- (1.918,0.423) -- (1.977,0.430) -- (2.035,0.445) -- (2.095,0.452) -- (2.154,0.454) -- (2.213,0.463) -- (2.272,0.463) -- (2.331,0.472) -- (2.390,0.475) -- (2.449,0.479) -- (2.508,0.484) -- (2.567,0.493) -- (2.626,0.504) -- (2.685,0.508) -- (2.744,0.513) -- (2.803,0.513) -- (2.861,0.514) -- (2.921,0.514) -- (2.980,0.515) -- (3.039,0.516) -- (3.098,0.537) -- (3.157,0.542) -- (3.216,0.543) -- (3.275,0.553) -- (3.333,0.562) -- (3.393,0.572) -- (3.451,0.574) -- (3.510,0.580) -- (3.570,0.584) -- (3.629,0.585) -- (3.688,0.590) -- (3.747,0.602) -- (3.806,0.613) -- (3.865,0.645) -- (3.924,0.657) -- (3.983,0.663) -- (4.042,0.670) -- (4.101,0.702) -- (4.160,0.702) -- (4.218,0.709) -- (4.277,0.720) -- (4.337,0.720) -- (4.396,0.735) -- (4.455,0.755) -- (4.514,0.764) -- (4.573,0.794) -- (4.632,0.797) -- (4.691,0.824) -- (4.750,0.834) -- (4.808,0.836) -- (4.867,0.842) -- (4.926,0.861) -- (4.986,0.883) -- (5.045,0.918) -- (5.104,0.958) -- (5.163,0.977) -- (5.222,0.990) -- (5.281,1.030) -- (5.340,1.107) -- (5.399,1.144) -- (5.458,1.145) -- (5.517,1.150) -- (5.575,1.260) -- (5.635,1.301) -- (5.694,1.339) -- (5.753,1.385) -- (5.812,3.600) -- (5.871,3.600);
  \node[font=\scriptsize,scred,align=left,anchor=north west] at (0.05,3.550) {exceed $\theta_0$: 6/100\\max run: 2};
  \draw[scblack] (0,0) rectangle (5.900,3.600);
  \draw[scblack] (0.000,0) -- (0.000,-0.08);
  \node[font=\scriptsize,below] at (0.000,-0.08) {0};
  \draw[scblack] (2.950,0) -- (2.950,-0.08);
  \node[font=\scriptsize,below] at (2.950,-0.08) {0.5};
  \draw[scblack] (5.900,0) -- (5.900,-0.08);
  \node[font=\scriptsize,below] at (5.900,-0.08) {1};
  \draw[scblack] (0,0.000) -- (-0.08,0.000);
  \node[font=\scriptsize,left] at (-0.10,0.000) {0};
  \draw[scblack] (0,1.200) -- (-0.08,1.200);
  \node[font=\scriptsize,left] at (-0.10,1.200) {1 ($\theta_0$)};
  \draw[scblack] (0,2.400) -- (-0.08,2.400);
  \node[font=\scriptsize,left] at (-0.10,2.400) {2};
  \draw[scblack] (0,3.600) -- (-0.08,3.600);
  \node[font=\scriptsize,left] at (-0.10,3.600) {3};
  \node[font=\scriptsize,below=0.42cm] at (2.950,0) {within-set quantile};
  \node[font=\scriptsize,rotate=90] at (-1.250,1.800) {score $/\ \theta_0$};
  \node[font={\scriptsize\bfseries},above] at (2.950,3.640) {valve1 (sorted score $/\ \theta_0$)};
  \end{scope}
  \begin{scope}[xshift=7.60cm]
  \draw[scblack,densely dotted,line width=0.9pt] (0,1.200) -- (5.900,1.200);
  \draw[scblue,dotted,line width=0.8pt] (0.070,0.035) -- (0.211,0.047) -- (0.351,0.096) -- (0.492,0.137) -- (0.632,0.138) -- (0.773,0.156) -- (0.913,0.171) -- (1.054,0.178) -- (1.194,0.230) -- (1.335,0.286) -- (1.475,0.375) -- (1.615,0.397) -- (1.756,0.413) -- (1.896,0.418) -- (2.037,0.427) -- (2.177,0.433) -- (2.318,0.466) -- (2.458,0.568) -- (2.599,0.574) -- (2.739,0.597) -- (2.880,0.606) -- (3.020,0.672) -- (3.161,0.722) -- (3.301,0.732) -- (3.442,0.767) -- (3.582,0.790) -- (3.723,0.796) -- (3.863,0.856) -- (4.004,0.938) -- (4.144,0.964) -- (4.285,1.037) -- (4.425,1.048) -- (4.565,1.058) -- (4.706,1.084) -- (4.846,1.110) -- (4.987,1.134) -- (5.127,1.154) -- (5.268,1.202) -- (5.408,1.329) -- (5.549,1.630) -- (5.689,1.707) -- (5.830,3.600);
  \draw[scred,solid,line width=1.0pt] (0.030,0.011) -- (0.090,0.014) -- (0.151,0.078) -- (0.211,0.078) -- (0.271,0.097) -- (0.331,0.098) -- (0.391,0.104) -- (0.452,0.143) -- (0.512,0.175) -- (0.572,0.179) -- (0.632,0.187) -- (0.692,0.225) -- (0.753,0.231) -- (0.813,0.258) -- (0.873,0.264) -- (0.933,0.269) -- (0.993,0.278) -- (1.054,0.327) -- (1.114,0.349) -- (1.174,0.361) -- (1.234,0.365) -- (1.294,0.369) -- (1.355,0.394) -- (1.415,0.408) -- (1.475,0.429) -- (1.535,0.439) -- (1.595,0.443) -- (1.656,0.461) -- (1.716,0.486) -- (1.776,0.492) -- (1.836,0.505) -- (1.896,0.510) -- (1.957,0.553) -- (2.017,0.556) -- (2.077,0.579) -- (2.137,0.579) -- (2.197,0.581) -- (2.258,0.583) -- (2.318,0.595) -- (2.378,0.628) -- (2.438,0.650) -- (2.498,0.652) -- (2.559,0.663) -- (2.619,0.672) -- (2.679,0.681) -- (2.739,0.705) -- (2.799,0.717) -- (2.860,0.718) -- (2.920,0.731) -- (2.980,0.743) -- (3.040,0.781) -- (3.101,0.784) -- (3.161,0.789) -- (3.221,0.789) -- (3.281,0.798) -- (3.341,0.803) -- (3.402,0.817) -- (3.462,0.830) -- (3.522,0.833) -- (3.582,0.835) -- (3.642,0.871) -- (3.703,0.883) -- (3.763,0.913) -- (3.823,0.931) -- (3.883,0.939) -- (3.943,0.967) -- (4.004,0.968) -- (4.064,0.971) -- (4.124,0.978) -- (4.184,1.004) -- (4.244,1.006) -- (4.305,1.014) -- (4.365,1.015) -- (4.425,1.028) -- (4.485,1.036) -- (4.545,1.071) -- (4.606,1.083) -- (4.666,1.085) -- (4.726,1.093) -- (4.786,1.105) -- (4.846,1.131) -- (4.907,1.150) -- (4.967,1.163) -- (5.027,1.168) -- (5.087,1.191) -- (5.147,1.244) -- (5.208,1.279) -- (5.268,1.309) -- (5.328,1.357) -- (5.388,1.365) -- (5.448,1.441) -- (5.509,1.456) -- (5.569,1.457) -- (5.629,1.469) -- (5.689,1.574) -- (5.749,1.596) -- (5.810,3.600) -- (5.870,3.600);
  \node[font=\scriptsize,scred,align=left,anchor=north west] at (0.05,3.550) {exceed $\theta_0$: 13/98\\max run: 2};
  \draw[scblack] (0,0) rectangle (5.900,3.600);
  \draw[scblack] (0.000,0) -- (0.000,-0.08);
  \node[font=\scriptsize,below] at (0.000,-0.08) {0};
  \draw[scblack] (2.950,0) -- (2.950,-0.08);
  \node[font=\scriptsize,below] at (2.950,-0.08) {0.5};
  \draw[scblack] (5.900,0) -- (5.900,-0.08);
  \node[font=\scriptsize,below] at (5.900,-0.08) {1};
  \draw[scblack] (0,0.000) -- (-0.08,0.000);
  \node[font=\scriptsize,left] at (-0.10,0.000) {0};
  \draw[scblack] (0,1.200) -- (-0.08,1.200);
  \node[font=\scriptsize,left] at (-0.10,1.200) {1 ($\theta_0$)};
  \draw[scblack] (0,2.400) -- (-0.08,2.400);
  \node[font=\scriptsize,left] at (-0.10,2.400) {2};
  \draw[scblack] (0,3.600) -- (-0.08,3.600);
  \node[font=\scriptsize,left] at (-0.10,3.600) {3};
  \node[font=\scriptsize,below=0.42cm] at (2.950,0) {within-set quantile};
  \node[font=\scriptsize,rotate=90] at (-1.250,1.800) {score $/\ \theta_0$};
  \node[font={\scriptsize\bfseries},above] at (2.950,3.640) {valve2 (sorted score $/\ \theta_0$)};
  \end{scope}
  \draw[scred,solid,line width=1.0pt] (0.000,-1.050) -- (0.500,-1.050);
  \node[font=\scriptsize,right] at (0.550,-1.050) {in-event windows};
  \draw[scblue,dotted,line width=0.8pt] (3.000,-1.050) -- (3.500,-1.050);
  \node[font=\scriptsize,right] at (3.550,-1.050) {out-of-event windows};
  \draw[scblack,densely dotted,line width=0.9pt] (6.800,-1.050) -- (7.300,-1.050);
  \node[font=\scriptsize,right] at (7.350,-1.050) {$\theta_0$ (ratio $=1$)};
\end{tikzpicture}%
}
  \caption{Scores sorted within each set, as a ratio to $\theta_0$: in-event windows (red
  solid), out-of-event windows (blue dotted), $\theta_0$ at ratio 1 (black dotted). Few
  in-event windows cross the line. The $y$ axis is clipped at $3\theta_0$; 2 (valve1) and
  3 (valve2) windows lie above the clip.}
  \label{fig:scorer-limit}
\end{figure}

\subsection{No harm on SKAB, and no evidence of benefit (C5; H2)}\label{sec:results-baseline}
\claim{5}{\Smeasured}{We measure that the controller is identical to the fixed
threshold on both SKAB streams.}
On valve1 both give 6 alert windows, 4 episodes, 0 false episodes and recall $1/1$; on
valve2 both give 18 windows, 12 episodes, 3 false episodes and recall $1/1$. The
controller makes 0 deferrals and 0 recalibrations, because its guard is never reached.
H2 therefore holds as a no-harm check; it says nothing about benefit. Fixed and
hysteresis detect each event at low alert volume. The Sun et al.\ arm decides on only
$54/143$ (valve1) and $57/140$ (valve2) windows and reaches its low workload by
abstaining; it is the only arm that defers on SKAB.

\begin{table}[tbp]
  \centering
  \caption{All 11 arms on the two SKAB development streams, one common score trace per
  stream. Recall and coverage are shown as num/den. Calibration threshold $\theta_0$:
  valve1 $1.309$ (43 windows), valve2 $1.126$ (40 windows). ``In-ev.\ alert'' is the
  in-event alerted fraction; ``--'' marks an undefined value (mean delay when an arm
  raises no in-event alert).}
  \label{tab:baseline}
  \scriptsize
  \setlength{\tabcolsep}{3.2pt}
  \begin{tabular}{llrrcrrrrcrr}
\toprule
Stream & Arm & Alert & Epis. & Recall & False & Mean & Non-ev. & In-ev. & Cover. & Def. & Recal. \\
& & win. & & $n/d$ & epis. & delay & dur. & alert & $n/d$ & & \\
\midrule
valve1 & fixed & 6 & 4 & 1/1 & 0 & 11.0 & 0 & 0.060 & 143/143 & 0 & 0 \\
 & rolling (n.o.) & 97 & 27 & 1/1 & 8 & 11.0 & 109 & 0.693 & 143/143 & 0 & 0 \\
 & rolling (all) & 10 & 10 & 1/1 & 2 & 11.0 & 10 & 0.075 & 143/143 & 0 & 0 \\
 & k-consec. & 0 & 0 & 0/1 & 0 & -- & 0 & 0.000 & 143/143 & 0 & 0 \\
 & m-of-n & 0 & 0 & 0/1 & 0 & -- & 0 & 0.000 & 143/143 & 0 & 0 \\
 & hysteresis & 6 & 4 & 1/1 & 0 & 11.0 & 0 & 0.060 & 143/143 & 0 & 0 \\
 & anchored & 6 & 4 & 1/1 & 0 & 11.0 & 0 & 0.060 & 143/143 & 0 & 0 \\
 & no-cap & 6 & 4 & 1/1 & 0 & 11.0 & 0 & 0.060 & 143/143 & 0 & 0 \\
 & no-stability & 6 & 4 & 1/1 & 0 & 11.0 & 0 & 0.060 & 143/143 & 0 & 0 \\
 & no-defer & 6 & 4 & 1/1 & 0 & 11.0 & 0 & 0.060 & 143/143 & 0 & 0 \\
 & Sun & 3 & 3 & 1/1 & 0 & 11.0 & 0 & 0.030 & 54/143 & 89 & 0 \\
\midrule
valve2 & fixed & 18 & 12 & 1/1 & 3 & 15.0 & 20 & 0.132 & 140/140 & 0 & 0 \\
 & rolling (n.o.) & 70 & 29 & 1/1 & 9 & 15.0 & 104 & 0.447 & 140/140 & 0 & 0 \\
 & rolling (all) & 13 & 12 & 1/1 & 2 & 15.0 & 12 & 0.102 & 140/140 & 0 & 0 \\
 & k-consec. & 0 & 0 & 0/1 & 0 & -- & 0 & 0.000 & 140/140 & 0 & 0 \\
 & m-of-n & 6 & 3 & 1/1 & 1 & 27.0 & 9 & 0.038 & 140/140 & 0 & 0 \\
 & hysteresis & 21 & 9 & 1/1 & 3 & 15.0 & 20 & 0.162 & 140/140 & 0 & 0 \\
 & anchored & 18 & 12 & 1/1 & 3 & 15.0 & 20 & 0.132 & 140/140 & 0 & 0 \\
 & no-cap & 18 & 12 & 1/1 & 3 & 15.0 & 20 & 0.132 & 140/140 & 0 & 0 \\
 & no-stability & 18 & 12 & 1/1 & 3 & 15.0 & 20 & 0.132 & 140/140 & 0 & 0 \\
 & no-defer & 18 & 12 & 1/1 & 3 & 15.0 & 20 & 0.132 & 140/140 & 0 & 0 \\
 & Sun & 2 & 2 & 1/1 & 1 & 15.0 & 4 & 0.010 & 57/140 & 83 & 0 \\
\bottomrule
\end{tabular}%

\end{table}

\subsection{Mechanism and declared failures under synthetic stress (C6; H3, H4)}\label{sec:results-synthetic}
\claim{6}{\Smeasured}{We measure where the mechanism works and where it fails.}
Table~\ref{tab:synth} and Figure~\ref{fig:synthetic-map} give the synthetic family.
\begin{itemize}[leftmargin=1.4em, itemsep=1pt, topsep=1pt]
  \item \emph{Works on a benign step (Y1).} Non-event alert duration is $16$, against
    $597$ for fixed and $917$ for normal-only rolling, at full coverage $260/260$, with
    one recalibration and no deferral.
  \item \emph{Still catches later faults (Y3, Y7).} After adapting, it alerts the spiky
    fault with recall $1/1$ and in-event alerted fraction $0.29$ on both.
  \item \emph{Fails on a benign ramp (Y2).} It recalibrates once during the ramp and then
    locks into periodic single-window alerts: $204$ over $50$ episodes at coverage
    $257/260$ (3 deferrals), against $577$ for fixed.
  \item \emph{Absorbs a fault shaped like a step (Y4), as declared.} Its in-event alerted
    fraction is $0.020$ against $0.746$ for fixed. No label-free rule can separate a
    benign shift from a fault of the same shape (H4).
\end{itemize}

\begin{table}[tbp]
  \centering
  \caption{Synthetic cases Y1--Y7, five arms (fixed, admit-all rolling, anchored,
  no-cap, Sun). ``Non-ev.\ dur.'' is non-event alert duration; ``In-ev.\ alert'' is the
  in-event alerted fraction; ``Def./Rec.'' is deferrals/recalibrations; ``--'' marks a
  value undefined because the case has no labelled event.}
  \label{tab:synth}
  \footnotesize
  \setlength{\tabcolsep}{4pt}
  \begin{tabular}{llrrcrr}
\toprule
Scen. & Arm & Non-ev. & In-ev. & Recall & Cover. & Def./Rec. \\
& & dur. & alert & $n/d$ & $n/d$ & \\
\midrule
Y1 & fixed & 597 & -- & -- & 260/260 & 0/0 \\
 & rolling (all) & 12 & -- & -- & 260/260 & 0/0 \\
 & anchored & 16 & -- & -- & 260/260 & 0/1 \\
 & no-cap & 16 & -- & -- & 260/260 & 0/1 \\
 & Sun & 8 & -- & -- & 22/260 & 238/0 \\
\midrule
Y2 & fixed & 577 & -- & -- & 260/260 & 0/0 \\
 & rolling (all) & 28 & -- & -- & 260/260 & 0/0 \\
 & anchored & 204 & -- & -- & 257/260 & 3/1 \\
 & no-cap & 204 & -- & -- & 257/260 & 3/1 \\
 & Sun & 24 & -- & -- & 23/260 & 237/0 \\
\midrule
Y3 & fixed & 521 & 1.000 & 1/1 & 260/260 & 0/0 \\
 & rolling (all) & 12 & 0.040 & 1/1 & 260/260 & 0/0 \\
 & anchored & 19 & 0.290 & 1/1 & 250/260 & 10/1 \\
 & no-cap & 16 & 0.160 & 1/1 & 260/260 & 0/3 \\
 & Sun & 8 & 0.040 & 1/1 & 22/260 & 238/0 \\
\midrule
Y4 & fixed & 0 & 0.746 & 1/1 & 260/260 & 0/0 \\
 & rolling (all) & 4 & 0.010 & 1/1 & 260/260 & 0/0 \\
 & anchored & 0 & 0.020 & 1/1 & 260/260 & 0/1 \\
 & no-cap & 0 & 0.020 & 1/1 & 260/260 & 0/1 \\
 & Sun & 0 & 0.010 & 1/1 & 22/260 & 238/0 \\
\midrule
Y5 & fixed & 3 & 0.490 & 1/1 & 260/260 & 0/0 \\
 & rolling (all) & 4 & 0.040 & 1/1 & 260/260 & 0/0 \\
 & anchored & 3 & 0.290 & 1/1 & 250/260 & 10/0 \\
 & no-cap & 0 & 0.160 & 1/1 & 260/260 & 0/2 \\
 & Sun & 0 & 0.040 & 1/1 & 90/260 & 170/0 \\
\midrule
Y6 & fixed & 797 & -- & -- & 260/260 & 0/0 \\
 & rolling (all) & 12 & -- & -- & 260/260 & 0/0 \\
 & anchored & 757 & -- & -- & 250/260 & 10/0 \\
 & no-cap & 16 & -- & -- & 259/260 & 1/1 \\
 & Sun & 8 & -- & -- & 22/260 & 238/0 \\
\midrule
Y7 & fixed & 227 & 0.490 & 1/1 & 260/260 & 0/0 \\
 & rolling (all) & 12 & 0.040 & 1/1 & 260/260 & 0/0 \\
 & anchored & 19 & 0.290 & 1/1 & 250/260 & 10/2 \\
 & no-cap & 16 & 0.160 & 1/1 & 260/260 & 0/4 \\
 & Sun & 8 & 0.040 & 1/1 & 60/260 & 200/0 \\
\bottomrule
\end{tabular}%

\end{table}

\begin{figure}[tbp]
  \centering
  \resizebox{0.78\linewidth}{!}{\definecolor{scblack}{RGB}{ 25, 32, 56}%
\definecolor{scblue}{RGB}{  0,114,178}%
\definecolor{scred}{RGB}{ 213, 94,  0}%
\definecolor{scgreen}{RGB}{  0,158,115}%
\definecolor{scgray}{RGB}{ 88,108,140}%
\definecolor{scevent}{RGB}{230,159,  0}%
\begin{tikzpicture}[>=stealth,line join=round,line cap=round]
  \begin{scope}[xshift=0.00cm,yshift=0.00cm]
  \draw[scgray,solid,line width=0.4pt] (0.000,0.000) -- (0.014,0.635) -- (0.028,0.635) -- (0.042,0.000) -- (0.056,0.000) -- (0.069,0.635) -- (0.083,0.635) -- (0.097,0.000) -- (0.111,0.000) -- (0.125,0.635) -- (0.139,0.635) -- (0.153,0.000) -- (0.167,0.000) -- (0.181,0.635) -- (0.195,0.635) -- (0.208,0.000) -- (0.222,0.000) -- (0.236,0.635) -- (0.250,0.635) -- (0.264,0.000) -- (0.278,0.000) -- (0.292,0.635) -- (0.306,0.635) -- (0.320,0.000) -- (0.334,0.000) -- (0.347,0.635) -- (0.361,0.635) -- (0.375,0.000) -- (0.389,0.000) -- (0.403,0.635) -- (0.417,0.635) -- (0.431,0.000) -- (0.445,0.000) -- (0.459,0.635) -- (0.473,0.635) -- (0.486,0.000) -- (0.500,0.000) -- (0.514,0.635) -- (0.528,0.635) -- (0.542,0.000) -- (0.556,0.000) -- (0.570,0.635) -- (0.584,0.635) -- (0.598,0.000) -- (0.612,0.000) -- (0.625,0.635) -- (0.639,0.635) -- (0.653,0.000) -- (0.667,0.000) -- (0.681,0.635) -- (0.695,0.635) -- (0.709,0.000) -- (0.723,0.000) -- (0.737,0.635) -- (0.751,0.635) -- (0.764,0.000) -- (0.778,0.000) -- (0.792,0.635) -- (0.806,0.635) -- (0.820,0.000) -- (0.834,1.270) -- (0.848,1.905) -- (0.862,0.635) -- (0.876,1.270) -- (0.890,1.270) -- (0.903,1.905) -- (0.917,0.635) -- (0.931,1.270) -- (0.945,1.270) -- (0.959,1.905) -- (0.973,0.635) -- (0.987,1.270) -- (1.001,1.270) -- (1.015,1.905) -- (1.029,0.635) -- (1.042,1.270) -- (1.056,1.270) -- (1.070,1.905) -- (1.084,0.635) -- (1.098,1.270) -- (1.112,1.270) -- (1.126,1.905) -- (1.140,0.635) -- (1.154,1.270) -- (1.168,1.270) -- (1.181,1.905) -- (1.195,0.635) -- (1.209,1.270) -- (1.223,1.270) -- (1.237,1.905) -- (1.251,0.635) -- (1.265,1.270) -- (1.279,1.270) -- (1.293,1.905) -- (1.307,0.635) -- (1.320,1.270) -- (1.334,1.270) -- (1.348,1.905) -- (1.362,0.635) -- (1.376,1.270) -- (1.390,1.270) -- (1.404,1.905) -- (1.418,0.635) -- (1.432,1.270) -- (1.446,1.270) -- (1.459,1.905) -- (1.473,0.635) -- (1.487,1.270) -- (1.501,1.270) -- (1.515,1.905) -- (1.529,0.635) -- (1.543,1.270) -- (1.557,1.270) -- (1.571,1.905) -- (1.585,0.635) -- (1.598,1.270) -- (1.612,1.270) -- (1.626,1.905) -- (1.640,0.635) -- (1.654,1.270) -- (1.668,1.270) -- (1.682,1.905) -- (1.696,0.635) -- (1.710,1.270) -- (1.724,1.270) -- (1.737,1.905) -- (1.751,0.635) -- (1.765,1.270) -- (1.779,1.270) -- (1.793,1.905) -- (1.807,0.635) -- (1.821,1.270) -- (1.835,1.270) -- (1.849,1.905) -- (1.863,0.635) -- (1.876,1.270) -- (1.890,1.270) -- (1.904,1.905) -- (1.918,0.635) -- (1.932,1.270) -- (1.946,1.270) -- (1.960,1.905) -- (1.974,0.635) -- (1.988,1.270) -- (2.002,1.270) -- (2.015,1.905) -- (2.029,0.635) -- (2.043,1.270) -- (2.057,1.270) -- (2.071,1.905) -- (2.085,0.635) -- (2.099,1.270) -- (2.113,1.270) -- (2.127,1.905) -- (2.141,0.635) -- (2.154,1.270) -- (2.168,1.270) -- (2.182,1.905) -- (2.196,0.635) -- (2.210,1.270) -- (2.224,1.270) -- (2.238,1.905) -- (2.252,0.635) -- (2.266,1.270) -- (2.280,1.270) -- (2.293,1.905) -- (2.307,0.635) -- (2.321,1.270) -- (2.335,1.270) -- (2.349,1.905) -- (2.363,0.635) -- (2.377,1.270) -- (2.391,1.270) -- (2.405,1.905) -- (2.419,0.635) -- (2.432,1.270) -- (2.446,1.270) -- (2.460,1.905) -- (2.474,0.635) -- (2.488,1.270) -- (2.502,1.270) -- (2.516,1.905) -- (2.530,0.635) -- (2.544,1.270) -- (2.558,1.270) -- (2.571,1.905) -- (2.585,0.635) -- (2.599,1.270) -- (2.613,1.270) -- (2.627,1.905) -- (2.641,0.635) -- (2.655,1.270) -- (2.669,1.270) -- (2.683,1.905) -- (2.697,0.635) -- (2.710,1.270) -- (2.724,1.270) -- (2.738,1.905) -- (2.752,0.635) -- (2.766,1.270) -- (2.780,1.270) -- (2.794,1.905) -- (2.808,0.635) -- (2.822,1.270) -- (2.836,1.270) -- (2.849,1.905) -- (2.863,0.635) -- (2.877,1.270) -- (2.891,1.270) -- (2.905,1.905) -- (2.919,0.635) -- (2.933,1.270) -- (2.947,1.270) -- (2.961,1.905) -- (2.975,0.635) -- (2.988,1.270) -- (3.002,1.270) -- (3.016,1.905) -- (3.030,0.635) -- (3.044,1.270) -- (3.058,1.270) -- (3.072,1.905) -- (3.086,0.635) -- (3.100,1.270) -- (3.114,1.270) -- (3.127,1.905) -- (3.141,0.635) -- (3.155,1.270) -- (3.169,1.270) -- (3.183,1.905) -- (3.197,0.635) -- (3.211,1.270) -- (3.225,1.270) -- (3.239,1.905) -- (3.253,0.635) -- (3.266,1.270) -- (3.280,1.270) -- (3.294,1.905) -- (3.308,0.635) -- (3.322,1.270) -- (3.336,1.270) -- (3.350,1.905) -- (3.364,0.635) -- (3.378,1.270) -- (3.392,1.270) -- (3.405,1.905) -- (3.419,0.635) -- (3.433,1.270) -- (3.447,1.270) -- (3.461,1.905) -- (3.475,0.635) -- (3.489,1.270) -- (3.503,1.270) -- (3.517,1.905) -- (3.531,0.635) -- (3.544,1.270) -- (3.558,1.270) -- (3.572,1.905) -- (3.586,0.635) -- (3.600,1.270);
  \draw[scblack,solid,line width=0.7pt] (0,0.635) -- (3.600,0.635);
  \draw[scblue,dotted,line width=0.7pt] (0.000,0.635) -- (0.014,0.000) -- (0.028,0.635) -- (0.042,0.635) -- (0.056,0.635) -- (0.069,0.635) -- (0.083,0.635) -- (0.097,0.635) -- (0.111,0.635) -- (0.125,0.635) -- (0.139,0.635) -- (0.153,0.635) -- (0.167,0.635) -- (0.181,0.635) -- (0.195,0.635) -- (0.208,0.635) -- (0.222,0.635) -- (0.236,0.635) -- (0.250,0.635) -- (0.264,0.635) -- (0.278,0.635) -- (0.292,0.635) -- (0.306,0.635) -- (0.320,0.635) -- (0.334,0.635) -- (0.347,0.635) -- (0.361,0.635) -- (0.375,0.635) -- (0.389,0.635) -- (0.403,0.635) -- (0.417,0.635) -- (0.431,0.635) -- (0.445,0.635) -- (0.459,0.635) -- (0.473,0.635) -- (0.486,0.635) -- (0.500,0.635) -- (0.514,0.635) -- (0.528,0.635) -- (0.542,0.635) -- (0.556,0.635) -- (0.570,0.635) -- (0.584,0.635) -- (0.598,0.635) -- (0.612,0.635) -- (0.625,0.635) -- (0.639,0.635) -- (0.653,0.635) -- (0.667,0.635) -- (0.681,0.635) -- (0.695,0.635) -- (0.709,0.635) -- (0.723,0.635) -- (0.737,0.635) -- (0.751,0.635) -- (0.764,0.635) -- (0.778,0.635) -- (0.792,0.635) -- (0.806,0.635) -- (0.820,0.635) -- (0.834,0.635) -- (0.848,1.270) -- (0.862,1.905) -- (0.876,1.905) -- (0.890,1.905) -- (0.903,1.905) -- (0.917,1.905) -- (0.931,1.905) -- (0.945,1.905) -- (0.959,1.905) -- (0.973,1.905) -- (0.987,1.905) -- (1.001,1.905) -- (1.015,1.905) -- (1.029,1.905) -- (1.042,1.905) -- (1.056,1.905) -- (1.070,1.905) -- (1.084,1.905) -- (1.098,1.905) -- (1.112,1.905) -- (1.126,1.905) -- (1.140,1.905) -- (1.154,1.905) -- (1.168,1.905) -- (1.181,1.905) -- (1.195,1.905) -- (1.209,1.905) -- (1.223,1.905) -- (1.237,1.905) -- (1.251,1.905) -- (1.265,1.905) -- (1.279,1.905) -- (1.293,1.905) -- (1.307,1.905) -- (1.320,1.905) -- (1.334,1.905) -- (1.348,1.905) -- (1.362,1.905) -- (1.376,1.905) -- (1.390,1.905) -- (1.404,1.905) -- (1.418,1.905) -- (1.432,1.905) -- (1.446,1.905) -- (1.459,1.905) -- (1.473,1.905) -- (1.487,1.905) -- (1.501,1.905) -- (1.515,1.905) -- (1.529,1.905) -- (1.543,1.905) -- (1.557,1.905) -- (1.571,1.905) -- (1.585,1.905) -- (1.598,1.905) -- (1.612,1.905) -- (1.626,1.905) -- (1.640,1.905) -- (1.654,1.905) -- (1.668,1.905) -- (1.682,1.905) -- (1.696,1.905) -- (1.710,1.905) -- (1.724,1.905) -- (1.737,1.905) -- (1.751,1.905) -- (1.765,1.905) -- (1.779,1.905) -- (1.793,1.905) -- (1.807,1.905) -- (1.821,1.905) -- (1.835,1.905) -- (1.849,1.905) -- (1.863,1.905) -- (1.876,1.905) -- (1.890,1.905) -- (1.904,1.905) -- (1.918,1.905) -- (1.932,1.905) -- (1.946,1.905) -- (1.960,1.905) -- (1.974,1.905) -- (1.988,1.905) -- (2.002,1.905) -- (2.015,1.905) -- (2.029,1.905) -- (2.043,1.905) -- (2.057,1.905) -- (2.071,1.905) -- (2.085,1.905) -- (2.099,1.905) -- (2.113,1.905) -- (2.127,1.905) -- (2.141,1.905) -- (2.154,1.905) -- (2.168,1.905) -- (2.182,1.905) -- (2.196,1.905) -- (2.210,1.905) -- (2.224,1.905) -- (2.238,1.905) -- (2.252,1.905) -- (2.266,1.905) -- (2.280,1.905) -- (2.293,1.905) -- (2.307,1.905) -- (2.321,1.905) -- (2.335,1.905) -- (2.349,1.905) -- (2.363,1.905) -- (2.377,1.905) -- (2.391,1.905) -- (2.405,1.905) -- (2.419,1.905) -- (2.432,1.905) -- (2.446,1.905) -- (2.460,1.905) -- (2.474,1.905) -- (2.488,1.905) -- (2.502,1.905) -- (2.516,1.905) -- (2.530,1.905) -- (2.544,1.905) -- (2.558,1.905) -- (2.571,1.905) -- (2.585,1.905) -- (2.599,1.905) -- (2.613,1.905) -- (2.627,1.905) -- (2.641,1.905) -- (2.655,1.905) -- (2.669,1.905) -- (2.683,1.905) -- (2.697,1.905) -- (2.710,1.905) -- (2.724,1.905) -- (2.738,1.905) -- (2.752,1.905) -- (2.766,1.905) -- (2.780,1.905) -- (2.794,1.905) -- (2.808,1.905) -- (2.822,1.905) -- (2.836,1.905) -- (2.849,1.905) -- (2.863,1.905) -- (2.877,1.905) -- (2.891,1.905) -- (2.905,1.905) -- (2.919,1.905) -- (2.933,1.905) -- (2.947,1.905) -- (2.961,1.905) -- (2.975,1.905) -- (2.988,1.905) -- (3.002,1.905) -- (3.016,1.905) -- (3.030,1.905) -- (3.044,1.905) -- (3.058,1.905) -- (3.072,1.905) -- (3.086,1.905) -- (3.100,1.905) -- (3.114,1.905) -- (3.127,1.905) -- (3.141,1.905) -- (3.155,1.905) -- (3.169,1.905) -- (3.183,1.905) -- (3.197,1.905) -- (3.211,1.905) -- (3.225,1.905) -- (3.239,1.905) -- (3.253,1.905) -- (3.266,1.905) -- (3.280,1.905) -- (3.294,1.905) -- (3.308,1.905) -- (3.322,1.905) -- (3.336,1.905) -- (3.350,1.905) -- (3.364,1.905) -- (3.378,1.905) -- (3.392,1.905) -- (3.405,1.905) -- (3.419,1.905) -- (3.433,1.905) -- (3.447,1.905) -- (3.461,1.905) -- (3.475,1.905) -- (3.489,1.905) -- (3.503,1.905) -- (3.517,1.905) -- (3.531,1.905) -- (3.544,1.905) -- (3.558,1.905) -- (3.572,1.905) -- (3.586,1.905) -- (3.600,1.905);
  \draw[scgreen,densely dashed,line width=1.0pt] (0.000,0.635) -- (0.014,0.635) -- (0.028,0.635) -- (0.042,0.635) -- (0.056,0.635) -- (0.069,0.635) -- (0.083,0.635) -- (0.097,0.635) -- (0.111,0.635) -- (0.125,0.635) -- (0.139,0.635) -- (0.153,0.635) -- (0.167,0.635) -- (0.181,0.635) -- (0.195,0.635) -- (0.208,0.635) -- (0.222,0.635) -- (0.236,0.635) -- (0.250,0.635) -- (0.264,0.635) -- (0.278,0.635) -- (0.292,0.635) -- (0.306,0.635) -- (0.320,0.635) -- (0.334,0.635) -- (0.347,0.635) -- (0.361,0.635) -- (0.375,0.635) -- (0.389,0.635) -- (0.403,0.635) -- (0.417,0.635) -- (0.431,0.635) -- (0.445,0.635) -- (0.459,0.635) -- (0.473,0.635) -- (0.486,0.635) -- (0.500,0.635) -- (0.514,0.635) -- (0.528,0.635) -- (0.542,0.635) -- (0.556,0.635) -- (0.570,0.635) -- (0.584,0.635) -- (0.598,0.635) -- (0.612,0.635) -- (0.625,0.635) -- (0.639,0.635) -- (0.653,0.635) -- (0.667,0.635) -- (0.681,0.635) -- (0.695,0.635) -- (0.709,0.635) -- (0.723,0.635) -- (0.737,0.635) -- (0.751,0.635) -- (0.764,0.635) -- (0.778,0.635) -- (0.792,0.635) -- (0.806,0.635) -- (0.820,0.635) -- (0.834,0.635) -- (0.848,0.635) -- (0.862,0.635) -- (0.876,0.635) -- (0.890,0.635) -- (0.903,1.905) -- (0.917,1.905) -- (0.931,1.905) -- (0.945,1.905) -- (0.959,1.905) -- (0.973,1.905) -- (0.987,1.905) -- (1.001,1.905) -- (1.015,1.905) -- (1.029,1.905) -- (1.042,1.905) -- (1.056,1.905) -- (1.070,1.905) -- (1.084,1.905) -- (1.098,1.905) -- (1.112,1.905) -- (1.126,1.905) -- (1.140,1.905) -- (1.154,1.905) -- (1.168,1.905) -- (1.181,1.905) -- (1.195,1.905) -- (1.209,1.905) -- (1.223,1.905) -- (1.237,1.905) -- (1.251,1.905) -- (1.265,1.905) -- (1.279,1.905) -- (1.293,1.905) -- (1.307,1.905) -- (1.320,1.905) -- (1.334,1.905) -- (1.348,1.905) -- (1.362,1.905) -- (1.376,1.905) -- (1.390,1.905) -- (1.404,1.905) -- (1.418,1.905) -- (1.432,1.905) -- (1.446,1.905) -- (1.459,1.905) -- (1.473,1.905) -- (1.487,1.905) -- (1.501,1.905) -- (1.515,1.905) -- (1.529,1.905) -- (1.543,1.905) -- (1.557,1.905) -- (1.571,1.905) -- (1.585,1.905) -- (1.598,1.905) -- (1.612,1.905) -- (1.626,1.905) -- (1.640,1.905) -- (1.654,1.905) -- (1.668,1.905) -- (1.682,1.905) -- (1.696,1.905) -- (1.710,1.905) -- (1.724,1.905) -- (1.737,1.905) -- (1.751,1.905) -- (1.765,1.905) -- (1.779,1.905) -- (1.793,1.905) -- (1.807,1.905) -- (1.821,1.905) -- (1.835,1.905) -- (1.849,1.905) -- (1.863,1.905) -- (1.876,1.905) -- (1.890,1.905) -- (1.904,1.905) -- (1.918,1.905) -- (1.932,1.905) -- (1.946,1.905) -- (1.960,1.905) -- (1.974,1.905) -- (1.988,1.905) -- (2.002,1.905) -- (2.015,1.905) -- (2.029,1.905) -- (2.043,1.905) -- (2.057,1.905) -- (2.071,1.905) -- (2.085,1.905) -- (2.099,1.905) -- (2.113,1.905) -- (2.127,1.905) -- (2.141,1.905) -- (2.154,1.905) -- (2.168,1.905) -- (2.182,1.905) -- (2.196,1.905) -- (2.210,1.905) -- (2.224,1.905) -- (2.238,1.905) -- (2.252,1.905) -- (2.266,1.905) -- (2.280,1.905) -- (2.293,1.905) -- (2.307,1.905) -- (2.321,1.905) -- (2.335,1.905) -- (2.349,1.905) -- (2.363,1.905) -- (2.377,1.905) -- (2.391,1.905) -- (2.405,1.905) -- (2.419,1.905) -- (2.432,1.905) -- (2.446,1.905) -- (2.460,1.905) -- (2.474,1.905) -- (2.488,1.905) -- (2.502,1.905) -- (2.516,1.905) -- (2.530,1.905) -- (2.544,1.905) -- (2.558,1.905) -- (2.571,1.905) -- (2.585,1.905) -- (2.599,1.905) -- (2.613,1.905) -- (2.627,1.905) -- (2.641,1.905) -- (2.655,1.905) -- (2.669,1.905) -- (2.683,1.905) -- (2.697,1.905) -- (2.710,1.905) -- (2.724,1.905) -- (2.738,1.905) -- (2.752,1.905) -- (2.766,1.905) -- (2.780,1.905) -- (2.794,1.905) -- (2.808,1.905) -- (2.822,1.905) -- (2.836,1.905) -- (2.849,1.905) -- (2.863,1.905) -- (2.877,1.905) -- (2.891,1.905) -- (2.905,1.905) -- (2.919,1.905) -- (2.933,1.905) -- (2.947,1.905) -- (2.961,1.905) -- (2.975,1.905) -- (2.988,1.905) -- (3.002,1.905) -- (3.016,1.905) -- (3.030,1.905) -- (3.044,1.905) -- (3.058,1.905) -- (3.072,1.905) -- (3.086,1.905) -- (3.100,1.905) -- (3.114,1.905) -- (3.127,1.905) -- (3.141,1.905) -- (3.155,1.905) -- (3.169,1.905) -- (3.183,1.905) -- (3.197,1.905) -- (3.211,1.905) -- (3.225,1.905) -- (3.239,1.905) -- (3.253,1.905) -- (3.266,1.905) -- (3.280,1.905) -- (3.294,1.905) -- (3.308,1.905) -- (3.322,1.905) -- (3.336,1.905) -- (3.350,1.905) -- (3.364,1.905) -- (3.378,1.905) -- (3.392,1.905) -- (3.405,1.905) -- (3.419,1.905) -- (3.433,1.905) -- (3.447,1.905) -- (3.461,1.905) -- (3.475,1.905) -- (3.489,1.905) -- (3.503,1.905) -- (3.517,1.905) -- (3.531,1.905) -- (3.544,1.905) -- (3.558,1.905) -- (3.572,1.905) -- (3.586,1.905) -- (3.600,1.905);
  \node[scgreen,font=\scriptsize] at (0.903,2.020) {$\bullet$};
  \draw[scblack] (0,0) rectangle (3.600,2.000);
  \draw[scblack] (0.000,0) -- (0.000,-0.08);
  \node[font=\scriptsize,below] at (0.000,-0.08) {563};
  \draw[scblack] (3.600,0) -- (3.600,-0.08);
  \node[font=\scriptsize,below] at (3.600,-0.08) {1599};
  \draw[scblack] (0,0.000) -- (-0.08,0.000);
  \node[font=\scriptsize,left] at (-0.10,0.000) {0};
  \draw[scblack] (0,0.635) -- (-0.08,0.635);
  \node[font=\scriptsize,left] at (-0.10,0.635) {$\theta_0$};
  \draw[scblack] (0,2.019) -- (-0.08,2.019);
  \node[font=\scriptsize,left] at (-0.10,2.019) {3.4};
  \node[font=\scriptsize,rotate=90] at (-0.720,1.000) {score};
  \node[font={\scriptsize\bfseries},above] at (1.800,2.040) {Y1};
  \end{scope}
  \begin{scope}[xshift=4.70cm,yshift=0.00cm]
  \draw[scgray,solid,line width=0.4pt] (0.000,0.000) -- (0.014,0.635) -- (0.028,0.635) -- (0.042,0.000) -- (0.056,0.000) -- (0.069,0.635) -- (0.083,0.635) -- (0.097,0.000) -- (0.111,0.000) -- (0.125,0.635) -- (0.139,0.635) -- (0.153,0.000) -- (0.167,0.000) -- (0.181,0.635) -- (0.195,0.635) -- (0.208,0.000) -- (0.222,0.000) -- (0.236,0.635) -- (0.250,0.635) -- (0.264,0.000) -- (0.278,0.000) -- (0.292,0.635) -- (0.306,0.635) -- (0.320,0.000) -- (0.334,0.000) -- (0.347,0.635) -- (0.361,0.635) -- (0.375,0.000) -- (0.389,0.000) -- (0.403,0.635) -- (0.417,0.635) -- (0.431,0.000) -- (0.445,0.000) -- (0.459,0.635) -- (0.473,0.635) -- (0.486,0.000) -- (0.500,0.000) -- (0.514,0.635) -- (0.528,0.635) -- (0.542,0.000) -- (0.556,0.000) -- (0.570,0.635) -- (0.584,0.635) -- (0.598,0.000) -- (0.612,0.000) -- (0.625,0.635) -- (0.639,0.635) -- (0.653,0.000) -- (0.667,0.000) -- (0.681,0.635) -- (0.695,0.635) -- (0.709,0.000) -- (0.723,0.000) -- (0.737,0.635) -- (0.751,0.635) -- (0.764,0.000) -- (0.778,0.000) -- (0.792,0.635) -- (0.806,0.635) -- (0.820,0.000) -- (0.834,0.159) -- (0.848,0.794) -- (0.862,0.397) -- (0.876,0.317) -- (0.890,0.317) -- (0.903,1.111) -- (0.917,0.159) -- (0.931,0.556) -- (0.945,0.635) -- (0.959,1.270) -- (0.973,0.159) -- (0.987,0.794) -- (1.001,0.873) -- (1.015,1.587) -- (1.029,0.317) -- (1.042,1.111) -- (1.056,1.111) -- (1.070,1.825) -- (1.084,0.635) -- (1.098,1.270) -- (1.112,1.270) -- (1.126,1.905) -- (1.140,0.635) -- (1.154,1.270) -- (1.168,1.270) -- (1.181,1.905) -- (1.195,0.635) -- (1.209,1.270) -- (1.223,1.270) -- (1.237,1.905) -- (1.251,0.635) -- (1.265,1.270) -- (1.279,1.270) -- (1.293,1.905) -- (1.307,0.635) -- (1.320,1.270) -- (1.334,1.270) -- (1.348,1.905) -- (1.362,0.635) -- (1.376,1.270) -- (1.390,1.270) -- (1.404,1.905) -- (1.418,0.635) -- (1.432,1.270) -- (1.446,1.270) -- (1.459,1.905) -- (1.473,0.635) -- (1.487,1.270) -- (1.501,1.270) -- (1.515,1.905) -- (1.529,0.635) -- (1.543,1.270) -- (1.557,1.270) -- (1.571,1.905) -- (1.585,0.635) -- (1.598,1.270) -- (1.612,1.270) -- (1.626,1.905) -- (1.640,0.635) -- (1.654,1.270) -- (1.668,1.270) -- (1.682,1.905) -- (1.696,0.635) -- (1.710,1.270) -- (1.724,1.270) -- (1.737,1.905) -- (1.751,0.635) -- (1.765,1.270) -- (1.779,1.270) -- (1.793,1.905) -- (1.807,0.635) -- (1.821,1.270) -- (1.835,1.270) -- (1.849,1.905) -- (1.863,0.635) -- (1.876,1.270) -- (1.890,1.270) -- (1.904,1.905) -- (1.918,0.635) -- (1.932,1.270) -- (1.946,1.270) -- (1.960,1.905) -- (1.974,0.635) -- (1.988,1.270) -- (2.002,1.270) -- (2.015,1.905) -- (2.029,0.635) -- (2.043,1.270) -- (2.057,1.270) -- (2.071,1.905) -- (2.085,0.635) -- (2.099,1.270) -- (2.113,1.270) -- (2.127,1.905) -- (2.141,0.635) -- (2.154,1.270) -- (2.168,1.270) -- (2.182,1.905) -- (2.196,0.635) -- (2.210,1.270) -- (2.224,1.270) -- (2.238,1.905) -- (2.252,0.635) -- (2.266,1.270) -- (2.280,1.270) -- (2.293,1.905) -- (2.307,0.635) -- (2.321,1.270) -- (2.335,1.270) -- (2.349,1.905) -- (2.363,0.635) -- (2.377,1.270) -- (2.391,1.270) -- (2.405,1.905) -- (2.419,0.635) -- (2.432,1.270) -- (2.446,1.270) -- (2.460,1.905) -- (2.474,0.635) -- (2.488,1.270) -- (2.502,1.270) -- (2.516,1.905) -- (2.530,0.635) -- (2.544,1.270) -- (2.558,1.270) -- (2.571,1.905) -- (2.585,0.635) -- (2.599,1.270) -- (2.613,1.270) -- (2.627,1.905) -- (2.641,0.635) -- (2.655,1.270) -- (2.669,1.270) -- (2.683,1.905) -- (2.697,0.635) -- (2.710,1.270) -- (2.724,1.270) -- (2.738,1.905) -- (2.752,0.635) -- (2.766,1.270) -- (2.780,1.270) -- (2.794,1.905) -- (2.808,0.635) -- (2.822,1.270) -- (2.836,1.270) -- (2.849,1.905) -- (2.863,0.635) -- (2.877,1.270) -- (2.891,1.270) -- (2.905,1.905) -- (2.919,0.635) -- (2.933,1.270) -- (2.947,1.270) -- (2.961,1.905) -- (2.975,0.635) -- (2.988,1.270) -- (3.002,1.270) -- (3.016,1.905) -- (3.030,0.635) -- (3.044,1.270) -- (3.058,1.270) -- (3.072,1.905) -- (3.086,0.635) -- (3.100,1.270) -- (3.114,1.270) -- (3.127,1.905) -- (3.141,0.635) -- (3.155,1.270) -- (3.169,1.270) -- (3.183,1.905) -- (3.197,0.635) -- (3.211,1.270) -- (3.225,1.270) -- (3.239,1.905) -- (3.253,0.635) -- (3.266,1.270) -- (3.280,1.270) -- (3.294,1.905) -- (3.308,0.635) -- (3.322,1.270) -- (3.336,1.270) -- (3.350,1.905) -- (3.364,0.635) -- (3.378,1.270) -- (3.392,1.270) -- (3.405,1.905) -- (3.419,0.635) -- (3.433,1.270) -- (3.447,1.270) -- (3.461,1.905) -- (3.475,0.635) -- (3.489,1.270) -- (3.503,1.270) -- (3.517,1.905) -- (3.531,0.635) -- (3.544,1.270) -- (3.558,1.270) -- (3.572,1.905) -- (3.586,0.635) -- (3.600,1.270);
  \draw[scblack,solid,line width=0.7pt] (0,0.635) -- (3.600,0.635);
  \draw[scblue,dotted,line width=0.7pt] (0.000,0.635) -- (0.014,0.000) -- (0.028,0.635) -- (0.042,0.635) -- (0.056,0.635) -- (0.069,0.635) -- (0.083,0.635) -- (0.097,0.635) -- (0.111,0.635) -- (0.125,0.635) -- (0.139,0.635) -- (0.153,0.635) -- (0.167,0.635) -- (0.181,0.635) -- (0.195,0.635) -- (0.208,0.635) -- (0.222,0.635) -- (0.236,0.635) -- (0.250,0.635) -- (0.264,0.635) -- (0.278,0.635) -- (0.292,0.635) -- (0.306,0.635) -- (0.320,0.635) -- (0.334,0.635) -- (0.347,0.635) -- (0.361,0.635) -- (0.375,0.635) -- (0.389,0.635) -- (0.403,0.635) -- (0.417,0.635) -- (0.431,0.635) -- (0.445,0.635) -- (0.459,0.635) -- (0.473,0.635) -- (0.486,0.635) -- (0.500,0.635) -- (0.514,0.635) -- (0.528,0.635) -- (0.542,0.635) -- (0.556,0.635) -- (0.570,0.635) -- (0.584,0.635) -- (0.598,0.635) -- (0.612,0.635) -- (0.625,0.635) -- (0.639,0.635) -- (0.653,0.635) -- (0.667,0.635) -- (0.681,0.635) -- (0.695,0.635) -- (0.709,0.635) -- (0.723,0.635) -- (0.737,0.635) -- (0.751,0.635) -- (0.764,0.635) -- (0.778,0.635) -- (0.792,0.635) -- (0.806,0.635) -- (0.820,0.635) -- (0.834,0.635) -- (0.848,0.635) -- (0.862,0.794) -- (0.876,0.794) -- (0.890,0.794) -- (0.903,0.794) -- (0.917,1.111) -- (0.931,1.111) -- (0.945,1.111) -- (0.959,1.111) -- (0.973,1.270) -- (0.987,1.270) -- (1.001,1.270) -- (1.015,1.270) -- (1.029,1.587) -- (1.042,1.587) -- (1.056,1.587) -- (1.070,1.587) -- (1.084,1.825) -- (1.098,1.825) -- (1.112,1.825) -- (1.126,1.825) -- (1.140,1.905) -- (1.154,1.905) -- (1.168,1.905) -- (1.181,1.905) -- (1.195,1.905) -- (1.209,1.905) -- (1.223,1.905) -- (1.237,1.905) -- (1.251,1.905) -- (1.265,1.905) -- (1.279,1.905) -- (1.293,1.905) -- (1.307,1.905) -- (1.320,1.905) -- (1.334,1.905) -- (1.348,1.905) -- (1.362,1.905) -- (1.376,1.905) -- (1.390,1.905) -- (1.404,1.905) -- (1.418,1.905) -- (1.432,1.905) -- (1.446,1.905) -- (1.459,1.905) -- (1.473,1.905) -- (1.487,1.905) -- (1.501,1.905) -- (1.515,1.905) -- (1.529,1.905) -- (1.543,1.905) -- (1.557,1.905) -- (1.571,1.905) -- (1.585,1.905) -- (1.598,1.905) -- (1.612,1.905) -- (1.626,1.905) -- (1.640,1.905) -- (1.654,1.905) -- (1.668,1.905) -- (1.682,1.905) -- (1.696,1.905) -- (1.710,1.905) -- (1.724,1.905) -- (1.737,1.905) -- (1.751,1.905) -- (1.765,1.905) -- (1.779,1.905) -- (1.793,1.905) -- (1.807,1.905) -- (1.821,1.905) -- (1.835,1.905) -- (1.849,1.905) -- (1.863,1.905) -- (1.876,1.905) -- (1.890,1.905) -- (1.904,1.905) -- (1.918,1.905) -- (1.932,1.905) -- (1.946,1.905) -- (1.960,1.905) -- (1.974,1.905) -- (1.988,1.905) -- (2.002,1.905) -- (2.015,1.905) -- (2.029,1.905) -- (2.043,1.905) -- (2.057,1.905) -- (2.071,1.905) -- (2.085,1.905) -- (2.099,1.905) -- (2.113,1.905) -- (2.127,1.905) -- (2.141,1.905) -- (2.154,1.905) -- (2.168,1.905) -- (2.182,1.905) -- (2.196,1.905) -- (2.210,1.905) -- (2.224,1.905) -- (2.238,1.905) -- (2.252,1.905) -- (2.266,1.905) -- (2.280,1.905) -- (2.293,1.905) -- (2.307,1.905) -- (2.321,1.905) -- (2.335,1.905) -- (2.349,1.905) -- (2.363,1.905) -- (2.377,1.905) -- (2.391,1.905) -- (2.405,1.905) -- (2.419,1.905) -- (2.432,1.905) -- (2.446,1.905) -- (2.460,1.905) -- (2.474,1.905) -- (2.488,1.905) -- (2.502,1.905) -- (2.516,1.905) -- (2.530,1.905) -- (2.544,1.905) -- (2.558,1.905) -- (2.571,1.905) -- (2.585,1.905) -- (2.599,1.905) -- (2.613,1.905) -- (2.627,1.905) -- (2.641,1.905) -- (2.655,1.905) -- (2.669,1.905) -- (2.683,1.905) -- (2.697,1.905) -- (2.710,1.905) -- (2.724,1.905) -- (2.738,1.905) -- (2.752,1.905) -- (2.766,1.905) -- (2.780,1.905) -- (2.794,1.905) -- (2.808,1.905) -- (2.822,1.905) -- (2.836,1.905) -- (2.849,1.905) -- (2.863,1.905) -- (2.877,1.905) -- (2.891,1.905) -- (2.905,1.905) -- (2.919,1.905) -- (2.933,1.905) -- (2.947,1.905) -- (2.961,1.905) -- (2.975,1.905) -- (2.988,1.905) -- (3.002,1.905) -- (3.016,1.905) -- (3.030,1.905) -- (3.044,1.905) -- (3.058,1.905) -- (3.072,1.905) -- (3.086,1.905) -- (3.100,1.905) -- (3.114,1.905) -- (3.127,1.905) -- (3.141,1.905) -- (3.155,1.905) -- (3.169,1.905) -- (3.183,1.905) -- (3.197,1.905) -- (3.211,1.905) -- (3.225,1.905) -- (3.239,1.905) -- (3.253,1.905) -- (3.266,1.905) -- (3.280,1.905) -- (3.294,1.905) -- (3.308,1.905) -- (3.322,1.905) -- (3.336,1.905) -- (3.350,1.905) -- (3.364,1.905) -- (3.378,1.905) -- (3.392,1.905) -- (3.405,1.905) -- (3.419,1.905) -- (3.433,1.905) -- (3.447,1.905) -- (3.461,1.905) -- (3.475,1.905) -- (3.489,1.905) -- (3.503,1.905) -- (3.517,1.905) -- (3.531,1.905) -- (3.544,1.905) -- (3.558,1.905) -- (3.572,1.905) -- (3.586,1.905) -- (3.600,1.905);
  \draw[scgreen,densely dashed,line width=1.0pt] (0.000,0.635) -- (0.014,0.635) -- (0.028,0.635) -- (0.042,0.635) -- (0.056,0.635) -- (0.069,0.635) -- (0.083,0.635) -- (0.097,0.635) -- (0.111,0.635) -- (0.125,0.635) -- (0.139,0.635) -- (0.153,0.635) -- (0.167,0.635) -- (0.181,0.635) -- (0.195,0.635) -- (0.208,0.635) -- (0.222,0.635) -- (0.236,0.635) -- (0.250,0.635) -- (0.264,0.635) -- (0.278,0.635) -- (0.292,0.635) -- (0.306,0.635) -- (0.320,0.635) -- (0.334,0.635) -- (0.347,0.635) -- (0.361,0.635) -- (0.375,0.635) -- (0.389,0.635) -- (0.403,0.635) -- (0.417,0.635) -- (0.431,0.635) -- (0.445,0.635) -- (0.459,0.635) -- (0.473,0.635) -- (0.486,0.635) -- (0.500,0.635) -- (0.514,0.635) -- (0.528,0.635) -- (0.542,0.635) -- (0.556,0.635) -- (0.570,0.635) -- (0.584,0.635) -- (0.598,0.635) -- (0.612,0.635) -- (0.625,0.635) -- (0.639,0.635) -- (0.653,0.635) -- (0.667,0.635) -- (0.681,0.635) -- (0.695,0.635) -- (0.709,0.635) -- (0.723,0.635) -- (0.737,0.635) -- (0.751,0.635) -- (0.764,0.635) -- (0.778,0.635) -- (0.792,0.635) -- (0.806,0.635) -- (0.820,0.635) -- (0.834,0.635) -- (0.848,0.635) -- (0.862,0.635) -- (0.876,0.635) -- (0.890,0.635) -- (0.903,0.635) -- (0.917,0.635) -- (0.931,0.635) -- (0.945,0.635) -- (0.959,0.635) -- (0.973,0.635) -- (0.987,0.635) -- (1.001,0.635) -- (1.015,0.635) -- (1.029,0.635) -- (1.042,0.635) -- (1.056,1.587) -- (1.070,1.587) -- (1.084,1.587) -- (1.098,1.587) -- (1.112,1.587) -- (1.126,1.587) -- (1.140,1.587) -- (1.154,1.587) -- (1.168,1.587) -- (1.181,1.587) -- (1.195,1.587) -- (1.209,1.587) -- (1.223,1.587) -- (1.237,1.587) -- (1.251,1.587) -- (1.265,1.587) -- (1.279,1.587) -- (1.293,1.587) -- (1.307,1.587) -- (1.320,1.587) -- (1.334,1.587) -- (1.348,1.587) -- (1.362,1.587) -- (1.376,1.587) -- (1.390,1.587) -- (1.404,1.587) -- (1.418,1.587) -- (1.432,1.587) -- (1.446,1.587) -- (1.459,1.587) -- (1.473,1.587) -- (1.487,1.587) -- (1.501,1.587) -- (1.515,1.587) -- (1.529,1.587) -- (1.543,1.587) -- (1.557,1.587) -- (1.571,1.587) -- (1.585,1.587) -- (1.598,1.587) -- (1.612,1.587) -- (1.626,1.587) -- (1.640,1.587) -- (1.654,1.587) -- (1.668,1.587) -- (1.682,1.587) -- (1.696,1.587) -- (1.710,1.587) -- (1.724,1.587) -- (1.737,1.587) -- (1.751,1.587) -- (1.765,1.587) -- (1.779,1.587) -- (1.793,1.587) -- (1.807,1.587) -- (1.821,1.587) -- (1.835,1.587) -- (1.849,1.587) -- (1.863,1.587) -- (1.876,1.587) -- (1.890,1.587) -- (1.904,1.587) -- (1.918,1.587) -- (1.932,1.587) -- (1.946,1.587) -- (1.960,1.587) -- (1.974,1.587) -- (1.988,1.587) -- (2.002,1.587) -- (2.015,1.587) -- (2.029,1.587) -- (2.043,1.587) -- (2.057,1.587) -- (2.071,1.587) -- (2.085,1.587) -- (2.099,1.587) -- (2.113,1.587) -- (2.127,1.587) -- (2.141,1.587) -- (2.154,1.587) -- (2.168,1.587) -- (2.182,1.587) -- (2.196,1.587) -- (2.210,1.587) -- (2.224,1.587) -- (2.238,1.587) -- (2.252,1.587) -- (2.266,1.587) -- (2.280,1.587) -- (2.293,1.587) -- (2.307,1.587) -- (2.321,1.587) -- (2.335,1.587) -- (2.349,1.587) -- (2.363,1.587) -- (2.377,1.587) -- (2.391,1.587) -- (2.405,1.587) -- (2.419,1.587) -- (2.432,1.587) -- (2.446,1.587) -- (2.460,1.587) -- (2.474,1.587) -- (2.488,1.587) -- (2.502,1.587) -- (2.516,1.587) -- (2.530,1.587) -- (2.544,1.587) -- (2.558,1.587) -- (2.571,1.587) -- (2.585,1.587) -- (2.599,1.587) -- (2.613,1.587) -- (2.627,1.587) -- (2.641,1.587) -- (2.655,1.587) -- (2.669,1.587) -- (2.683,1.587) -- (2.697,1.587) -- (2.710,1.587) -- (2.724,1.587) -- (2.738,1.587) -- (2.752,1.587) -- (2.766,1.587) -- (2.780,1.587) -- (2.794,1.587) -- (2.808,1.587) -- (2.822,1.587) -- (2.836,1.587) -- (2.849,1.587) -- (2.863,1.587) -- (2.877,1.587) -- (2.891,1.587) -- (2.905,1.587) -- (2.919,1.587) -- (2.933,1.587) -- (2.947,1.587) -- (2.961,1.587) -- (2.975,1.587) -- (2.988,1.587) -- (3.002,1.587) -- (3.016,1.587) -- (3.030,1.587) -- (3.044,1.587) -- (3.058,1.587) -- (3.072,1.587) -- (3.086,1.587) -- (3.100,1.587) -- (3.114,1.587) -- (3.127,1.587) -- (3.141,1.587) -- (3.155,1.587) -- (3.169,1.587) -- (3.183,1.587) -- (3.197,1.587) -- (3.211,1.587) -- (3.225,1.587) -- (3.239,1.587) -- (3.253,1.587) -- (3.266,1.587) -- (3.280,1.587) -- (3.294,1.587) -- (3.308,1.587) -- (3.322,1.587) -- (3.336,1.587) -- (3.350,1.587) -- (3.364,1.587) -- (3.378,1.587) -- (3.392,1.587) -- (3.405,1.587) -- (3.419,1.587) -- (3.433,1.587) -- (3.447,1.587) -- (3.461,1.587) -- (3.475,1.587) -- (3.489,1.587) -- (3.503,1.587) -- (3.517,1.587) -- (3.531,1.587) -- (3.544,1.587) -- (3.558,1.587) -- (3.572,1.587) -- (3.586,1.587) -- (3.600,1.587);
  \draw[scred,line width=0.5pt] (1.015,0) -- (1.015,0.16);
  \draw[scred,line width=0.5pt] (1.029,0) -- (1.029,0.16);
  \draw[scred,line width=0.5pt] (1.042,0) -- (1.042,0.16);
  \node[scgreen,font=\scriptsize] at (1.056,2.020) {$\bullet$};
  \draw[scblack] (0,0) rectangle (3.600,2.000);
  \draw[scblack] (0.000,0) -- (0.000,-0.08);
  \node[font=\scriptsize,below] at (0.000,-0.08) {563};
  \draw[scblack] (3.600,0) -- (3.600,-0.08);
  \node[font=\scriptsize,below] at (3.600,-0.08) {1599};
  \draw[scblack] (0,0.000) -- (-0.08,0.000);
  \node[font=\scriptsize,left] at (-0.10,0.000) {0};
  \draw[scblack] (0,0.635) -- (-0.08,0.635);
  \node[font=\scriptsize,left] at (-0.10,0.635) {$\theta_0$};
  \draw[scblack] (0,2.019) -- (-0.08,2.019);
  \node[font=\scriptsize,left] at (-0.10,2.019) {3.4};
  \node[font={\scriptsize\bfseries},above] at (1.800,2.040) {Y2};
  \end{scope}
  \begin{scope}[xshift=9.40cm,yshift=0.00cm]
  \fill[scevent,opacity=0.14] (2.214,0) rectangle (2.561,2.000);
  \draw[scgray,solid,line width=0.4pt] (0.000,0.000) -- (0.014,0.238) -- (0.028,0.238) -- (0.042,0.000) -- (0.056,0.000) -- (0.069,0.238) -- (0.083,0.238) -- (0.097,0.000) -- (0.111,0.000) -- (0.125,0.238) -- (0.139,0.238) -- (0.153,0.000) -- (0.167,0.000) -- (0.181,0.238) -- (0.195,0.238) -- (0.208,0.000) -- (0.222,0.000) -- (0.236,0.238) -- (0.250,0.238) -- (0.264,0.000) -- (0.278,0.000) -- (0.292,0.238) -- (0.306,0.238) -- (0.320,0.000) -- (0.334,0.000) -- (0.347,0.238) -- (0.361,0.238) -- (0.375,0.000) -- (0.389,0.000) -- (0.403,0.238) -- (0.417,0.238) -- (0.431,0.000) -- (0.445,0.000) -- (0.459,0.238) -- (0.473,0.238) -- (0.486,0.000) -- (0.500,0.000) -- (0.514,0.238) -- (0.528,0.238) -- (0.542,0.000) -- (0.556,0.000) -- (0.570,0.238) -- (0.584,0.238) -- (0.598,0.000) -- (0.612,0.000) -- (0.625,0.238) -- (0.639,0.238) -- (0.653,0.000) -- (0.667,0.000) -- (0.681,0.238) -- (0.695,0.238) -- (0.709,0.000) -- (0.723,0.000) -- (0.737,0.238) -- (0.751,0.238) -- (0.764,0.000) -- (0.778,0.000) -- (0.792,0.238) -- (0.806,0.238) -- (0.820,0.000) -- (0.834,0.476) -- (0.848,0.714) -- (0.862,0.238) -- (0.876,0.476) -- (0.890,0.476) -- (0.903,0.714) -- (0.917,0.238) -- (0.931,0.476) -- (0.945,0.476) -- (0.959,0.714) -- (0.973,0.238) -- (0.987,0.476) -- (1.001,0.476) -- (1.015,0.714) -- (1.029,0.238) -- (1.042,0.476) -- (1.056,0.476) -- (1.070,0.714) -- (1.084,0.238) -- (1.098,0.476) -- (1.112,0.476) -- (1.126,0.714) -- (1.140,0.238) -- (1.154,0.476) -- (1.168,0.476) -- (1.181,0.714) -- (1.195,0.238) -- (1.209,0.476) -- (1.223,0.476) -- (1.237,0.714) -- (1.251,0.238) -- (1.265,0.476) -- (1.279,0.476) -- (1.293,0.714) -- (1.307,0.238) -- (1.320,0.476) -- (1.334,0.476) -- (1.348,0.714) -- (1.362,0.238) -- (1.376,0.476) -- (1.390,0.476) -- (1.404,0.714) -- (1.418,0.238) -- (1.432,0.476) -- (1.446,0.476) -- (1.459,0.714) -- (1.473,0.238) -- (1.487,0.476) -- (1.501,0.476) -- (1.515,0.714) -- (1.529,0.238) -- (1.543,0.476) -- (1.557,0.476) -- (1.571,0.714) -- (1.585,0.238) -- (1.598,0.476) -- (1.612,0.476) -- (1.626,0.714) -- (1.640,0.238) -- (1.654,0.476) -- (1.668,0.476) -- (1.682,0.714) -- (1.696,0.238) -- (1.710,0.476) -- (1.724,0.476) -- (1.737,0.714) -- (1.751,0.238) -- (1.765,0.476) -- (1.779,0.476) -- (1.793,0.714) -- (1.807,0.238) -- (1.821,0.476) -- (1.835,0.476) -- (1.849,0.714) -- (1.863,0.238) -- (1.876,0.476) -- (1.890,0.476) -- (1.904,0.714) -- (1.918,0.238) -- (1.932,0.476) -- (1.946,0.476) -- (1.960,0.714) -- (1.974,0.238) -- (1.988,0.476) -- (2.002,0.476) -- (2.015,0.714) -- (2.029,0.238) -- (2.043,0.476) -- (2.057,0.476) -- (2.071,0.714) -- (2.085,0.238) -- (2.099,0.476) -- (2.113,0.476) -- (2.127,0.714) -- (2.141,0.238) -- (2.154,0.476) -- (2.168,0.476) -- (2.182,0.714) -- (2.196,0.238) -- (2.210,0.476) -- (2.224,1.905) -- (2.238,0.714) -- (2.252,1.667) -- (2.266,0.476) -- (2.280,1.905) -- (2.293,0.714) -- (2.307,1.667) -- (2.321,0.476) -- (2.335,1.905) -- (2.349,0.714) -- (2.363,1.667) -- (2.377,0.476) -- (2.391,1.905) -- (2.405,0.714) -- (2.419,1.667) -- (2.432,0.476) -- (2.446,1.905) -- (2.460,0.714) -- (2.474,1.667) -- (2.488,0.476) -- (2.502,1.905) -- (2.516,0.714) -- (2.530,1.667) -- (2.544,0.476) -- (2.558,1.905) -- (2.571,0.714) -- (2.585,0.238) -- (2.599,0.476) -- (2.613,0.476) -- (2.627,0.714) -- (2.641,0.238) -- (2.655,0.476) -- (2.669,0.476) -- (2.683,0.714) -- (2.697,0.238) -- (2.710,0.476) -- (2.724,0.476) -- (2.738,0.714) -- (2.752,0.238) -- (2.766,0.476) -- (2.780,0.476) -- (2.794,0.714) -- (2.808,0.238) -- (2.822,0.476) -- (2.836,0.476) -- (2.849,0.714) -- (2.863,0.238) -- (2.877,0.476) -- (2.891,0.476) -- (2.905,0.714) -- (2.919,0.238) -- (2.933,0.476) -- (2.947,0.476) -- (2.961,0.714) -- (2.975,0.238) -- (2.988,0.476) -- (3.002,0.476) -- (3.016,0.714) -- (3.030,0.238) -- (3.044,0.476) -- (3.058,0.476) -- (3.072,0.714) -- (3.086,0.238) -- (3.100,0.476) -- (3.114,0.476) -- (3.127,0.714) -- (3.141,0.238) -- (3.155,0.476) -- (3.169,0.476) -- (3.183,0.714) -- (3.197,0.238) -- (3.211,0.476) -- (3.225,0.476) -- (3.239,0.714) -- (3.253,0.238) -- (3.266,0.476) -- (3.280,0.476) -- (3.294,0.714) -- (3.308,0.238) -- (3.322,0.476) -- (3.336,0.476) -- (3.350,0.714) -- (3.364,0.238) -- (3.378,0.476) -- (3.392,0.476) -- (3.405,0.714) -- (3.419,0.238) -- (3.433,0.476) -- (3.447,0.476) -- (3.461,0.714) -- (3.475,0.238) -- (3.489,0.476) -- (3.503,0.476) -- (3.517,0.714) -- (3.531,0.238) -- (3.544,0.476) -- (3.558,0.476) -- (3.572,0.714) -- (3.586,0.238) -- (3.600,0.476);
  \draw[scblack,solid,line width=0.7pt] (0,0.238) -- (3.600,0.238);
  \draw[scblue,dotted,line width=0.7pt] (0.000,0.238) -- (0.014,0.000) -- (0.028,0.238) -- (0.042,0.238) -- (0.056,0.238) -- (0.069,0.238) -- (0.083,0.238) -- (0.097,0.238) -- (0.111,0.238) -- (0.125,0.238) -- (0.139,0.238) -- (0.153,0.238) -- (0.167,0.238) -- (0.181,0.238) -- (0.195,0.238) -- (0.208,0.238) -- (0.222,0.238) -- (0.236,0.238) -- (0.250,0.238) -- (0.264,0.238) -- (0.278,0.238) -- (0.292,0.238) -- (0.306,0.238) -- (0.320,0.238) -- (0.334,0.238) -- (0.347,0.238) -- (0.361,0.238) -- (0.375,0.238) -- (0.389,0.238) -- (0.403,0.238) -- (0.417,0.238) -- (0.431,0.238) -- (0.445,0.238) -- (0.459,0.238) -- (0.473,0.238) -- (0.486,0.238) -- (0.500,0.238) -- (0.514,0.238) -- (0.528,0.238) -- (0.542,0.238) -- (0.556,0.238) -- (0.570,0.238) -- (0.584,0.238) -- (0.598,0.238) -- (0.612,0.238) -- (0.625,0.238) -- (0.639,0.238) -- (0.653,0.238) -- (0.667,0.238) -- (0.681,0.238) -- (0.695,0.238) -- (0.709,0.238) -- (0.723,0.238) -- (0.737,0.238) -- (0.751,0.238) -- (0.764,0.238) -- (0.778,0.238) -- (0.792,0.238) -- (0.806,0.238) -- (0.820,0.238) -- (0.834,0.238) -- (0.848,0.476) -- (0.862,0.714) -- (0.876,0.714) -- (0.890,0.714) -- (0.903,0.714) -- (0.917,0.714) -- (0.931,0.714) -- (0.945,0.714) -- (0.959,0.714) -- (0.973,0.714) -- (0.987,0.714) -- (1.001,0.714) -- (1.015,0.714) -- (1.029,0.714) -- (1.042,0.714) -- (1.056,0.714) -- (1.070,0.714) -- (1.084,0.714) -- (1.098,0.714) -- (1.112,0.714) -- (1.126,0.714) -- (1.140,0.714) -- (1.154,0.714) -- (1.168,0.714) -- (1.181,0.714) -- (1.195,0.714) -- (1.209,0.714) -- (1.223,0.714) -- (1.237,0.714) -- (1.251,0.714) -- (1.265,0.714) -- (1.279,0.714) -- (1.293,0.714) -- (1.307,0.714) -- (1.320,0.714) -- (1.334,0.714) -- (1.348,0.714) -- (1.362,0.714) -- (1.376,0.714) -- (1.390,0.714) -- (1.404,0.714) -- (1.418,0.714) -- (1.432,0.714) -- (1.446,0.714) -- (1.459,0.714) -- (1.473,0.714) -- (1.487,0.714) -- (1.501,0.714) -- (1.515,0.714) -- (1.529,0.714) -- (1.543,0.714) -- (1.557,0.714) -- (1.571,0.714) -- (1.585,0.714) -- (1.598,0.714) -- (1.612,0.714) -- (1.626,0.714) -- (1.640,0.714) -- (1.654,0.714) -- (1.668,0.714) -- (1.682,0.714) -- (1.696,0.714) -- (1.710,0.714) -- (1.724,0.714) -- (1.737,0.714) -- (1.751,0.714) -- (1.765,0.714) -- (1.779,0.714) -- (1.793,0.714) -- (1.807,0.714) -- (1.821,0.714) -- (1.835,0.714) -- (1.849,0.714) -- (1.863,0.714) -- (1.876,0.714) -- (1.890,0.714) -- (1.904,0.714) -- (1.918,0.714) -- (1.932,0.714) -- (1.946,0.714) -- (1.960,0.714) -- (1.974,0.714) -- (1.988,0.714) -- (2.002,0.714) -- (2.015,0.714) -- (2.029,0.714) -- (2.043,0.714) -- (2.057,0.714) -- (2.071,0.714) -- (2.085,0.714) -- (2.099,0.714) -- (2.113,0.714) -- (2.127,0.714) -- (2.141,0.714) -- (2.154,0.714) -- (2.168,0.714) -- (2.182,0.714) -- (2.196,0.714) -- (2.210,0.714) -- (2.224,0.714) -- (2.238,1.905) -- (2.252,1.905) -- (2.266,1.905) -- (2.280,1.905) -- (2.293,1.905) -- (2.307,1.905) -- (2.321,1.905) -- (2.335,1.905) -- (2.349,1.905) -- (2.363,1.905) -- (2.377,1.905) -- (2.391,1.905) -- (2.405,1.905) -- (2.419,1.905) -- (2.432,1.905) -- (2.446,1.905) -- (2.460,1.905) -- (2.474,1.905) -- (2.488,1.905) -- (2.502,1.905) -- (2.516,1.905) -- (2.530,1.905) -- (2.544,1.905) -- (2.558,1.905) -- (2.571,1.905) -- (2.585,1.905) -- (2.599,1.905) -- (2.613,1.905) -- (2.627,1.905) -- (2.641,1.905) -- (2.655,1.905) -- (2.669,1.905) -- (2.683,1.905) -- (2.697,1.905) -- (2.710,0.714) -- (2.724,0.714) -- (2.738,0.714) -- (2.752,0.714) -- (2.766,0.714) -- (2.780,0.714) -- (2.794,0.714) -- (2.808,0.714) -- (2.822,0.714) -- (2.836,0.714) -- (2.849,0.714) -- (2.863,0.714) -- (2.877,0.714) -- (2.891,0.714) -- (2.905,0.714) -- (2.919,0.714) -- (2.933,0.714) -- (2.947,0.714) -- (2.961,0.714) -- (2.975,0.714) -- (2.988,0.714) -- (3.002,0.714) -- (3.016,0.714) -- (3.030,0.714) -- (3.044,0.714) -- (3.058,0.714) -- (3.072,0.714) -- (3.086,0.714) -- (3.100,0.714) -- (3.114,0.714) -- (3.127,0.714) -- (3.141,0.714) -- (3.155,0.714) -- (3.169,0.714) -- (3.183,0.714) -- (3.197,0.714) -- (3.211,0.714) -- (3.225,0.714) -- (3.239,0.714) -- (3.253,0.714) -- (3.266,0.714) -- (3.280,0.714) -- (3.294,0.714) -- (3.308,0.714) -- (3.322,0.714) -- (3.336,0.714) -- (3.350,0.714) -- (3.364,0.714) -- (3.378,0.714) -- (3.392,0.714) -- (3.405,0.714) -- (3.419,0.714) -- (3.433,0.714) -- (3.447,0.714) -- (3.461,0.714) -- (3.475,0.714) -- (3.489,0.714) -- (3.503,0.714) -- (3.517,0.714) -- (3.531,0.714) -- (3.544,0.714) -- (3.558,0.714) -- (3.572,0.714) -- (3.586,0.714) -- (3.600,0.714);
  \draw[scgreen,densely dashed,line width=1.0pt] (0.000,0.238) -- (0.014,0.238) -- (0.028,0.238) -- (0.042,0.238) -- (0.056,0.238) -- (0.069,0.238) -- (0.083,0.238) -- (0.097,0.238) -- (0.111,0.238) -- (0.125,0.238) -- (0.139,0.238) -- (0.153,0.238) -- (0.167,0.238) -- (0.181,0.238) -- (0.195,0.238) -- (0.208,0.238) -- (0.222,0.238) -- (0.236,0.238) -- (0.250,0.238) -- (0.264,0.238) -- (0.278,0.238) -- (0.292,0.238) -- (0.306,0.238) -- (0.320,0.238) -- (0.334,0.238) -- (0.347,0.238) -- (0.361,0.238) -- (0.375,0.238) -- (0.389,0.238) -- (0.403,0.238) -- (0.417,0.238) -- (0.431,0.238) -- (0.445,0.238) -- (0.459,0.238) -- (0.473,0.238) -- (0.486,0.238) -- (0.500,0.238) -- (0.514,0.238) -- (0.528,0.238) -- (0.542,0.238) -- (0.556,0.238) -- (0.570,0.238) -- (0.584,0.238) -- (0.598,0.238) -- (0.612,0.238) -- (0.625,0.238) -- (0.639,0.238) -- (0.653,0.238) -- (0.667,0.238) -- (0.681,0.238) -- (0.695,0.238) -- (0.709,0.238) -- (0.723,0.238) -- (0.737,0.238) -- (0.751,0.238) -- (0.764,0.238) -- (0.778,0.238) -- (0.792,0.238) -- (0.806,0.238) -- (0.820,0.238) -- (0.834,0.238) -- (0.848,0.238) -- (0.862,0.238) -- (0.876,0.238) -- (0.890,0.238) -- (0.903,0.714) -- (0.917,0.714) -- (0.931,0.714) -- (0.945,0.714) -- (0.959,0.714) -- (0.973,0.714) -- (0.987,0.714) -- (1.001,0.714) -- (1.015,0.714) -- (1.029,0.714) -- (1.042,0.714) -- (1.056,0.714) -- (1.070,0.714) -- (1.084,0.714) -- (1.098,0.714) -- (1.112,0.714) -- (1.126,0.714) -- (1.140,0.714) -- (1.154,0.714) -- (1.168,0.714) -- (1.181,0.714) -- (1.195,0.714) -- (1.209,0.714) -- (1.223,0.714) -- (1.237,0.714) -- (1.251,0.714) -- (1.265,0.714) -- (1.279,0.714) -- (1.293,0.714) -- (1.307,0.714) -- (1.320,0.714) -- (1.334,0.714) -- (1.348,0.714) -- (1.362,0.714) -- (1.376,0.714) -- (1.390,0.714) -- (1.404,0.714) -- (1.418,0.714) -- (1.432,0.714) -- (1.446,0.714) -- (1.459,0.714) -- (1.473,0.714) -- (1.487,0.714) -- (1.501,0.714) -- (1.515,0.714) -- (1.529,0.714) -- (1.543,0.714) -- (1.557,0.714) -- (1.571,0.714) -- (1.585,0.714) -- (1.598,0.714) -- (1.612,0.714) -- (1.626,0.714) -- (1.640,0.714) -- (1.654,0.714) -- (1.668,0.714) -- (1.682,0.714) -- (1.696,0.714) -- (1.710,0.714) -- (1.724,0.714) -- (1.737,0.714) -- (1.751,0.714) -- (1.765,0.714) -- (1.779,0.714) -- (1.793,0.714) -- (1.807,0.714) -- (1.821,0.714) -- (1.835,0.714) -- (1.849,0.714) -- (1.863,0.714) -- (1.876,0.714) -- (1.890,0.714) -- (1.904,0.714) -- (1.918,0.714) -- (1.932,0.714) -- (1.946,0.714) -- (1.960,0.714) -- (1.974,0.714) -- (1.988,0.714) -- (2.002,0.714) -- (2.015,0.714) -- (2.029,0.714) -- (2.043,0.714) -- (2.057,0.714) -- (2.071,0.714) -- (2.085,0.714) -- (2.099,0.714) -- (2.113,0.714) -- (2.127,0.714) -- (2.141,0.714) -- (2.154,0.714) -- (2.168,0.714) -- (2.182,0.714) -- (2.196,0.714) -- (2.210,0.714) -- (2.224,0.714) -- (2.238,0.714) -- (2.252,0.714) -- (2.266,0.714) -- (2.280,0.714) -- (2.293,0.714) -- (2.307,0.714) -- (2.321,0.714) -- (2.335,0.714) -- (2.349,0.714) -- (2.363,0.714) -- (2.377,0.714) -- (2.391,0.714) -- (2.405,0.714) -- (2.419,0.714) -- (2.432,0.714) -- (2.446,0.714) -- (2.460,0.714) -- (2.474,0.714) -- (2.488,0.714) -- (2.502,0.714) -- (2.516,0.714) -- (2.530,0.714) -- (2.544,0.714) -- (2.558,0.714) -- (2.571,0.714) -- (2.585,0.714) -- (2.599,0.714) -- (2.613,0.714) -- (2.627,0.714) -- (2.641,0.714) -- (2.655,0.714) -- (2.669,0.714) -- (2.683,0.714) -- (2.697,0.714) -- (2.710,0.714) -- (2.724,0.714) -- (2.738,0.714) -- (2.752,0.714) -- (2.766,0.714) -- (2.780,0.714) -- (2.794,0.714) -- (2.808,0.714) -- (2.822,0.714) -- (2.836,0.714) -- (2.849,0.714) -- (2.863,0.714) -- (2.877,0.714) -- (2.891,0.714) -- (2.905,0.714) -- (2.919,0.714) -- (2.933,0.714) -- (2.947,0.714) -- (2.961,0.714) -- (2.975,0.714) -- (2.988,0.714) -- (3.002,0.714) -- (3.016,0.714) -- (3.030,0.714) -- (3.044,0.714) -- (3.058,0.714) -- (3.072,0.714) -- (3.086,0.714) -- (3.100,0.714) -- (3.114,0.714) -- (3.127,0.714) -- (3.141,0.714) -- (3.155,0.714) -- (3.169,0.714) -- (3.183,0.714) -- (3.197,0.714) -- (3.211,0.714) -- (3.225,0.714) -- (3.239,0.714) -- (3.253,0.714) -- (3.266,0.714) -- (3.280,0.714) -- (3.294,0.714) -- (3.308,0.714) -- (3.322,0.714) -- (3.336,0.714) -- (3.350,0.714) -- (3.364,0.714) -- (3.378,0.714) -- (3.392,0.714) -- (3.405,0.714) -- (3.419,0.714) -- (3.433,0.714) -- (3.447,0.714) -- (3.461,0.714) -- (3.475,0.714) -- (3.489,0.714) -- (3.503,0.714) -- (3.517,0.714) -- (3.531,0.714) -- (3.544,0.714) -- (3.558,0.714) -- (3.572,0.714) -- (3.586,0.714) -- (3.600,0.714);
  \node[scgreen,font=\scriptsize] at (0.903,2.020) {$\bullet$};
  \draw[scred,line width=0.5pt] (2.335,0) -- (2.335,0.16);
  \draw[scred,line width=0.5pt] (2.349,0) -- (2.349,0.16);
  \draw[scred,line width=0.5pt] (2.363,0) -- (2.363,0.16);
  \draw[scred,line width=0.5pt] (2.377,0) -- (2.377,0.16);
  \draw[scred,line width=0.5pt] (2.391,0) -- (2.391,0.16);
  \draw[scred,line width=0.5pt] (2.405,0) -- (2.405,0.16);
  \draw[scred,line width=0.5pt] (2.419,0) -- (2.419,0.16);
  \draw[scred,line width=0.5pt] (2.432,0) -- (2.432,0.16);
  \draw[scred,line width=0.5pt] (2.446,0) -- (2.446,0.16);
  \draw[scred,line width=0.5pt] (2.460,0) -- (2.460,0.16);
  \draw[scblack] (0,0) rectangle (3.600,2.000);
  \draw[scblack] (0.000,0) -- (0.000,-0.08);
  \node[font=\scriptsize,below] at (0.000,-0.08) {563};
  \draw[scblack] (3.600,0) -- (3.600,-0.08);
  \node[font=\scriptsize,below] at (3.600,-0.08) {1599};
  \draw[scblack] (0,0.000) -- (-0.08,0.000);
  \node[font=\scriptsize,left] at (-0.10,0.000) {0};
  \draw[scblack] (0,0.238) -- (-0.08,0.238);
  \node[font=\scriptsize,left] at (-0.10,0.238) {$\theta_0$};
  \draw[scblack] (0,2.004) -- (-0.08,2.004);
  \node[font=\scriptsize,left] at (-0.10,2.004) {9.0};
  \node[font={\scriptsize\bfseries},above] at (1.800,2.040) {Y3};
  \end{scope}
  \begin{scope}[xshift=0.00cm,yshift=-3.15cm]
  \fill[scevent,opacity=0.14] (0.824,0) rectangle (3.600,2.000);
  \draw[scgray,solid,line width=0.4pt] (0.000,0.000) -- (0.014,0.635) -- (0.028,0.635) -- (0.042,0.000) -- (0.056,0.000) -- (0.069,0.635) -- (0.083,0.635) -- (0.097,0.000) -- (0.111,0.000) -- (0.125,0.635) -- (0.139,0.635) -- (0.153,0.000) -- (0.167,0.000) -- (0.181,0.635) -- (0.195,0.635) -- (0.208,0.000) -- (0.222,0.000) -- (0.236,0.635) -- (0.250,0.635) -- (0.264,0.000) -- (0.278,0.000) -- (0.292,0.635) -- (0.306,0.635) -- (0.320,0.000) -- (0.334,0.000) -- (0.347,0.635) -- (0.361,0.635) -- (0.375,0.000) -- (0.389,0.000) -- (0.403,0.635) -- (0.417,0.635) -- (0.431,0.000) -- (0.445,0.000) -- (0.459,0.635) -- (0.473,0.635) -- (0.486,0.000) -- (0.500,0.000) -- (0.514,0.635) -- (0.528,0.635) -- (0.542,0.000) -- (0.556,0.000) -- (0.570,0.635) -- (0.584,0.635) -- (0.598,0.000) -- (0.612,0.000) -- (0.625,0.635) -- (0.639,0.635) -- (0.653,0.000) -- (0.667,0.000) -- (0.681,0.635) -- (0.695,0.635) -- (0.709,0.000) -- (0.723,0.000) -- (0.737,0.635) -- (0.751,0.635) -- (0.764,0.000) -- (0.778,0.000) -- (0.792,0.635) -- (0.806,0.635) -- (0.820,0.000) -- (0.834,1.270) -- (0.848,1.905) -- (0.862,0.635) -- (0.876,1.270) -- (0.890,1.270) -- (0.903,1.905) -- (0.917,0.635) -- (0.931,1.270) -- (0.945,1.270) -- (0.959,1.905) -- (0.973,0.635) -- (0.987,1.270) -- (1.001,1.270) -- (1.015,1.905) -- (1.029,0.635) -- (1.042,1.270) -- (1.056,1.270) -- (1.070,1.905) -- (1.084,0.635) -- (1.098,1.270) -- (1.112,1.270) -- (1.126,1.905) -- (1.140,0.635) -- (1.154,1.270) -- (1.168,1.270) -- (1.181,1.905) -- (1.195,0.635) -- (1.209,1.270) -- (1.223,1.270) -- (1.237,1.905) -- (1.251,0.635) -- (1.265,1.270) -- (1.279,1.270) -- (1.293,1.905) -- (1.307,0.635) -- (1.320,1.270) -- (1.334,1.270) -- (1.348,1.905) -- (1.362,0.635) -- (1.376,1.270) -- (1.390,1.270) -- (1.404,1.905) -- (1.418,0.635) -- (1.432,1.270) -- (1.446,1.270) -- (1.459,1.905) -- (1.473,0.635) -- (1.487,1.270) -- (1.501,1.270) -- (1.515,1.905) -- (1.529,0.635) -- (1.543,1.270) -- (1.557,1.270) -- (1.571,1.905) -- (1.585,0.635) -- (1.598,1.270) -- (1.612,1.270) -- (1.626,1.905) -- (1.640,0.635) -- (1.654,1.270) -- (1.668,1.270) -- (1.682,1.905) -- (1.696,0.635) -- (1.710,1.270) -- (1.724,1.270) -- (1.737,1.905) -- (1.751,0.635) -- (1.765,1.270) -- (1.779,1.270) -- (1.793,1.905) -- (1.807,0.635) -- (1.821,1.270) -- (1.835,1.270) -- (1.849,1.905) -- (1.863,0.635) -- (1.876,1.270) -- (1.890,1.270) -- (1.904,1.905) -- (1.918,0.635) -- (1.932,1.270) -- (1.946,1.270) -- (1.960,1.905) -- (1.974,0.635) -- (1.988,1.270) -- (2.002,1.270) -- (2.015,1.905) -- (2.029,0.635) -- (2.043,1.270) -- (2.057,1.270) -- (2.071,1.905) -- (2.085,0.635) -- (2.099,1.270) -- (2.113,1.270) -- (2.127,1.905) -- (2.141,0.635) -- (2.154,1.270) -- (2.168,1.270) -- (2.182,1.905) -- (2.196,0.635) -- (2.210,1.270) -- (2.224,1.270) -- (2.238,1.905) -- (2.252,0.635) -- (2.266,1.270) -- (2.280,1.270) -- (2.293,1.905) -- (2.307,0.635) -- (2.321,1.270) -- (2.335,1.270) -- (2.349,1.905) -- (2.363,0.635) -- (2.377,1.270) -- (2.391,1.270) -- (2.405,1.905) -- (2.419,0.635) -- (2.432,1.270) -- (2.446,1.270) -- (2.460,1.905) -- (2.474,0.635) -- (2.488,1.270) -- (2.502,1.270) -- (2.516,1.905) -- (2.530,0.635) -- (2.544,1.270) -- (2.558,1.270) -- (2.571,1.905) -- (2.585,0.635) -- (2.599,1.270) -- (2.613,1.270) -- (2.627,1.905) -- (2.641,0.635) -- (2.655,1.270) -- (2.669,1.270) -- (2.683,1.905) -- (2.697,0.635) -- (2.710,1.270) -- (2.724,1.270) -- (2.738,1.905) -- (2.752,0.635) -- (2.766,1.270) -- (2.780,1.270) -- (2.794,1.905) -- (2.808,0.635) -- (2.822,1.270) -- (2.836,1.270) -- (2.849,1.905) -- (2.863,0.635) -- (2.877,1.270) -- (2.891,1.270) -- (2.905,1.905) -- (2.919,0.635) -- (2.933,1.270) -- (2.947,1.270) -- (2.961,1.905) -- (2.975,0.635) -- (2.988,1.270) -- (3.002,1.270) -- (3.016,1.905) -- (3.030,0.635) -- (3.044,1.270) -- (3.058,1.270) -- (3.072,1.905) -- (3.086,0.635) -- (3.100,1.270) -- (3.114,1.270) -- (3.127,1.905) -- (3.141,0.635) -- (3.155,1.270) -- (3.169,1.270) -- (3.183,1.905) -- (3.197,0.635) -- (3.211,1.270) -- (3.225,1.270) -- (3.239,1.905) -- (3.253,0.635) -- (3.266,1.270) -- (3.280,1.270) -- (3.294,1.905) -- (3.308,0.635) -- (3.322,1.270) -- (3.336,1.270) -- (3.350,1.905) -- (3.364,0.635) -- (3.378,1.270) -- (3.392,1.270) -- (3.405,1.905) -- (3.419,0.635) -- (3.433,1.270) -- (3.447,1.270) -- (3.461,1.905) -- (3.475,0.635) -- (3.489,1.270) -- (3.503,1.270) -- (3.517,1.905) -- (3.531,0.635) -- (3.544,1.270) -- (3.558,1.270) -- (3.572,1.905) -- (3.586,0.635) -- (3.600,1.270);
  \draw[scblack,solid,line width=0.7pt] (0,0.635) -- (3.600,0.635);
  \draw[scblue,dotted,line width=0.7pt] (0.000,0.635) -- (0.014,0.000) -- (0.028,0.635) -- (0.042,0.635) -- (0.056,0.635) -- (0.069,0.635) -- (0.083,0.635) -- (0.097,0.635) -- (0.111,0.635) -- (0.125,0.635) -- (0.139,0.635) -- (0.153,0.635) -- (0.167,0.635) -- (0.181,0.635) -- (0.195,0.635) -- (0.208,0.635) -- (0.222,0.635) -- (0.236,0.635) -- (0.250,0.635) -- (0.264,0.635) -- (0.278,0.635) -- (0.292,0.635) -- (0.306,0.635) -- (0.320,0.635) -- (0.334,0.635) -- (0.347,0.635) -- (0.361,0.635) -- (0.375,0.635) -- (0.389,0.635) -- (0.403,0.635) -- (0.417,0.635) -- (0.431,0.635) -- (0.445,0.635) -- (0.459,0.635) -- (0.473,0.635) -- (0.486,0.635) -- (0.500,0.635) -- (0.514,0.635) -- (0.528,0.635) -- (0.542,0.635) -- (0.556,0.635) -- (0.570,0.635) -- (0.584,0.635) -- (0.598,0.635) -- (0.612,0.635) -- (0.625,0.635) -- (0.639,0.635) -- (0.653,0.635) -- (0.667,0.635) -- (0.681,0.635) -- (0.695,0.635) -- (0.709,0.635) -- (0.723,0.635) -- (0.737,0.635) -- (0.751,0.635) -- (0.764,0.635) -- (0.778,0.635) -- (0.792,0.635) -- (0.806,0.635) -- (0.820,0.635) -- (0.834,0.635) -- (0.848,1.270) -- (0.862,1.905) -- (0.876,1.905) -- (0.890,1.905) -- (0.903,1.905) -- (0.917,1.905) -- (0.931,1.905) -- (0.945,1.905) -- (0.959,1.905) -- (0.973,1.905) -- (0.987,1.905) -- (1.001,1.905) -- (1.015,1.905) -- (1.029,1.905) -- (1.042,1.905) -- (1.056,1.905) -- (1.070,1.905) -- (1.084,1.905) -- (1.098,1.905) -- (1.112,1.905) -- (1.126,1.905) -- (1.140,1.905) -- (1.154,1.905) -- (1.168,1.905) -- (1.181,1.905) -- (1.195,1.905) -- (1.209,1.905) -- (1.223,1.905) -- (1.237,1.905) -- (1.251,1.905) -- (1.265,1.905) -- (1.279,1.905) -- (1.293,1.905) -- (1.307,1.905) -- (1.320,1.905) -- (1.334,1.905) -- (1.348,1.905) -- (1.362,1.905) -- (1.376,1.905) -- (1.390,1.905) -- (1.404,1.905) -- (1.418,1.905) -- (1.432,1.905) -- (1.446,1.905) -- (1.459,1.905) -- (1.473,1.905) -- (1.487,1.905) -- (1.501,1.905) -- (1.515,1.905) -- (1.529,1.905) -- (1.543,1.905) -- (1.557,1.905) -- (1.571,1.905) -- (1.585,1.905) -- (1.598,1.905) -- (1.612,1.905) -- (1.626,1.905) -- (1.640,1.905) -- (1.654,1.905) -- (1.668,1.905) -- (1.682,1.905) -- (1.696,1.905) -- (1.710,1.905) -- (1.724,1.905) -- (1.737,1.905) -- (1.751,1.905) -- (1.765,1.905) -- (1.779,1.905) -- (1.793,1.905) -- (1.807,1.905) -- (1.821,1.905) -- (1.835,1.905) -- (1.849,1.905) -- (1.863,1.905) -- (1.876,1.905) -- (1.890,1.905) -- (1.904,1.905) -- (1.918,1.905) -- (1.932,1.905) -- (1.946,1.905) -- (1.960,1.905) -- (1.974,1.905) -- (1.988,1.905) -- (2.002,1.905) -- (2.015,1.905) -- (2.029,1.905) -- (2.043,1.905) -- (2.057,1.905) -- (2.071,1.905) -- (2.085,1.905) -- (2.099,1.905) -- (2.113,1.905) -- (2.127,1.905) -- (2.141,1.905) -- (2.154,1.905) -- (2.168,1.905) -- (2.182,1.905) -- (2.196,1.905) -- (2.210,1.905) -- (2.224,1.905) -- (2.238,1.905) -- (2.252,1.905) -- (2.266,1.905) -- (2.280,1.905) -- (2.293,1.905) -- (2.307,1.905) -- (2.321,1.905) -- (2.335,1.905) -- (2.349,1.905) -- (2.363,1.905) -- (2.377,1.905) -- (2.391,1.905) -- (2.405,1.905) -- (2.419,1.905) -- (2.432,1.905) -- (2.446,1.905) -- (2.460,1.905) -- (2.474,1.905) -- (2.488,1.905) -- (2.502,1.905) -- (2.516,1.905) -- (2.530,1.905) -- (2.544,1.905) -- (2.558,1.905) -- (2.571,1.905) -- (2.585,1.905) -- (2.599,1.905) -- (2.613,1.905) -- (2.627,1.905) -- (2.641,1.905) -- (2.655,1.905) -- (2.669,1.905) -- (2.683,1.905) -- (2.697,1.905) -- (2.710,1.905) -- (2.724,1.905) -- (2.738,1.905) -- (2.752,1.905) -- (2.766,1.905) -- (2.780,1.905) -- (2.794,1.905) -- (2.808,1.905) -- (2.822,1.905) -- (2.836,1.905) -- (2.849,1.905) -- (2.863,1.905) -- (2.877,1.905) -- (2.891,1.905) -- (2.905,1.905) -- (2.919,1.905) -- (2.933,1.905) -- (2.947,1.905) -- (2.961,1.905) -- (2.975,1.905) -- (2.988,1.905) -- (3.002,1.905) -- (3.016,1.905) -- (3.030,1.905) -- (3.044,1.905) -- (3.058,1.905) -- (3.072,1.905) -- (3.086,1.905) -- (3.100,1.905) -- (3.114,1.905) -- (3.127,1.905) -- (3.141,1.905) -- (3.155,1.905) -- (3.169,1.905) -- (3.183,1.905) -- (3.197,1.905) -- (3.211,1.905) -- (3.225,1.905) -- (3.239,1.905) -- (3.253,1.905) -- (3.266,1.905) -- (3.280,1.905) -- (3.294,1.905) -- (3.308,1.905) -- (3.322,1.905) -- (3.336,1.905) -- (3.350,1.905) -- (3.364,1.905) -- (3.378,1.905) -- (3.392,1.905) -- (3.405,1.905) -- (3.419,1.905) -- (3.433,1.905) -- (3.447,1.905) -- (3.461,1.905) -- (3.475,1.905) -- (3.489,1.905) -- (3.503,1.905) -- (3.517,1.905) -- (3.531,1.905) -- (3.544,1.905) -- (3.558,1.905) -- (3.572,1.905) -- (3.586,1.905) -- (3.600,1.905);
  \draw[scgreen,densely dashed,line width=1.0pt] (0.000,0.635) -- (0.014,0.635) -- (0.028,0.635) -- (0.042,0.635) -- (0.056,0.635) -- (0.069,0.635) -- (0.083,0.635) -- (0.097,0.635) -- (0.111,0.635) -- (0.125,0.635) -- (0.139,0.635) -- (0.153,0.635) -- (0.167,0.635) -- (0.181,0.635) -- (0.195,0.635) -- (0.208,0.635) -- (0.222,0.635) -- (0.236,0.635) -- (0.250,0.635) -- (0.264,0.635) -- (0.278,0.635) -- (0.292,0.635) -- (0.306,0.635) -- (0.320,0.635) -- (0.334,0.635) -- (0.347,0.635) -- (0.361,0.635) -- (0.375,0.635) -- (0.389,0.635) -- (0.403,0.635) -- (0.417,0.635) -- (0.431,0.635) -- (0.445,0.635) -- (0.459,0.635) -- (0.473,0.635) -- (0.486,0.635) -- (0.500,0.635) -- (0.514,0.635) -- (0.528,0.635) -- (0.542,0.635) -- (0.556,0.635) -- (0.570,0.635) -- (0.584,0.635) -- (0.598,0.635) -- (0.612,0.635) -- (0.625,0.635) -- (0.639,0.635) -- (0.653,0.635) -- (0.667,0.635) -- (0.681,0.635) -- (0.695,0.635) -- (0.709,0.635) -- (0.723,0.635) -- (0.737,0.635) -- (0.751,0.635) -- (0.764,0.635) -- (0.778,0.635) -- (0.792,0.635) -- (0.806,0.635) -- (0.820,0.635) -- (0.834,0.635) -- (0.848,0.635) -- (0.862,0.635) -- (0.876,0.635) -- (0.890,0.635) -- (0.903,1.905) -- (0.917,1.905) -- (0.931,1.905) -- (0.945,1.905) -- (0.959,1.905) -- (0.973,1.905) -- (0.987,1.905) -- (1.001,1.905) -- (1.015,1.905) -- (1.029,1.905) -- (1.042,1.905) -- (1.056,1.905) -- (1.070,1.905) -- (1.084,1.905) -- (1.098,1.905) -- (1.112,1.905) -- (1.126,1.905) -- (1.140,1.905) -- (1.154,1.905) -- (1.168,1.905) -- (1.181,1.905) -- (1.195,1.905) -- (1.209,1.905) -- (1.223,1.905) -- (1.237,1.905) -- (1.251,1.905) -- (1.265,1.905) -- (1.279,1.905) -- (1.293,1.905) -- (1.307,1.905) -- (1.320,1.905) -- (1.334,1.905) -- (1.348,1.905) -- (1.362,1.905) -- (1.376,1.905) -- (1.390,1.905) -- (1.404,1.905) -- (1.418,1.905) -- (1.432,1.905) -- (1.446,1.905) -- (1.459,1.905) -- (1.473,1.905) -- (1.487,1.905) -- (1.501,1.905) -- (1.515,1.905) -- (1.529,1.905) -- (1.543,1.905) -- (1.557,1.905) -- (1.571,1.905) -- (1.585,1.905) -- (1.598,1.905) -- (1.612,1.905) -- (1.626,1.905) -- (1.640,1.905) -- (1.654,1.905) -- (1.668,1.905) -- (1.682,1.905) -- (1.696,1.905) -- (1.710,1.905) -- (1.724,1.905) -- (1.737,1.905) -- (1.751,1.905) -- (1.765,1.905) -- (1.779,1.905) -- (1.793,1.905) -- (1.807,1.905) -- (1.821,1.905) -- (1.835,1.905) -- (1.849,1.905) -- (1.863,1.905) -- (1.876,1.905) -- (1.890,1.905) -- (1.904,1.905) -- (1.918,1.905) -- (1.932,1.905) -- (1.946,1.905) -- (1.960,1.905) -- (1.974,1.905) -- (1.988,1.905) -- (2.002,1.905) -- (2.015,1.905) -- (2.029,1.905) -- (2.043,1.905) -- (2.057,1.905) -- (2.071,1.905) -- (2.085,1.905) -- (2.099,1.905) -- (2.113,1.905) -- (2.127,1.905) -- (2.141,1.905) -- (2.154,1.905) -- (2.168,1.905) -- (2.182,1.905) -- (2.196,1.905) -- (2.210,1.905) -- (2.224,1.905) -- (2.238,1.905) -- (2.252,1.905) -- (2.266,1.905) -- (2.280,1.905) -- (2.293,1.905) -- (2.307,1.905) -- (2.321,1.905) -- (2.335,1.905) -- (2.349,1.905) -- (2.363,1.905) -- (2.377,1.905) -- (2.391,1.905) -- (2.405,1.905) -- (2.419,1.905) -- (2.432,1.905) -- (2.446,1.905) -- (2.460,1.905) -- (2.474,1.905) -- (2.488,1.905) -- (2.502,1.905) -- (2.516,1.905) -- (2.530,1.905) -- (2.544,1.905) -- (2.558,1.905) -- (2.571,1.905) -- (2.585,1.905) -- (2.599,1.905) -- (2.613,1.905) -- (2.627,1.905) -- (2.641,1.905) -- (2.655,1.905) -- (2.669,1.905) -- (2.683,1.905) -- (2.697,1.905) -- (2.710,1.905) -- (2.724,1.905) -- (2.738,1.905) -- (2.752,1.905) -- (2.766,1.905) -- (2.780,1.905) -- (2.794,1.905) -- (2.808,1.905) -- (2.822,1.905) -- (2.836,1.905) -- (2.849,1.905) -- (2.863,1.905) -- (2.877,1.905) -- (2.891,1.905) -- (2.905,1.905) -- (2.919,1.905) -- (2.933,1.905) -- (2.947,1.905) -- (2.961,1.905) -- (2.975,1.905) -- (2.988,1.905) -- (3.002,1.905) -- (3.016,1.905) -- (3.030,1.905) -- (3.044,1.905) -- (3.058,1.905) -- (3.072,1.905) -- (3.086,1.905) -- (3.100,1.905) -- (3.114,1.905) -- (3.127,1.905) -- (3.141,1.905) -- (3.155,1.905) -- (3.169,1.905) -- (3.183,1.905) -- (3.197,1.905) -- (3.211,1.905) -- (3.225,1.905) -- (3.239,1.905) -- (3.253,1.905) -- (3.266,1.905) -- (3.280,1.905) -- (3.294,1.905) -- (3.308,1.905) -- (3.322,1.905) -- (3.336,1.905) -- (3.350,1.905) -- (3.364,1.905) -- (3.378,1.905) -- (3.392,1.905) -- (3.405,1.905) -- (3.419,1.905) -- (3.433,1.905) -- (3.447,1.905) -- (3.461,1.905) -- (3.475,1.905) -- (3.489,1.905) -- (3.503,1.905) -- (3.517,1.905) -- (3.531,1.905) -- (3.544,1.905) -- (3.558,1.905) -- (3.572,1.905) -- (3.586,1.905) -- (3.600,1.905);
  \node[scgreen,font=\scriptsize] at (0.903,2.020) {$\bullet$};
  \draw[scblack] (0,0) rectangle (3.600,2.000);
  \draw[scblack] (0.000,0) -- (0.000,-0.08);
  \node[font=\scriptsize,below] at (0.000,-0.08) {563};
  \draw[scblack] (3.600,0) -- (3.600,-0.08);
  \node[font=\scriptsize,below] at (3.600,-0.08) {1599};
  \draw[scblack] (0,0.000) -- (-0.08,0.000);
  \node[font=\scriptsize,left] at (-0.10,0.000) {0};
  \draw[scblack] (0,0.635) -- (-0.08,0.635);
  \node[font=\scriptsize,left] at (-0.10,0.635) {$\theta_0$};
  \draw[scblack] (0,2.019) -- (-0.08,2.019);
  \node[font=\scriptsize,left] at (-0.10,2.019) {3.4};
  \node[font=\scriptsize,rotate=90] at (-0.720,1.000) {score};
  \node[font={\scriptsize\bfseries},above] at (1.800,2.040) {Y4};
  \end{scope}
  \begin{scope}[xshift=4.70cm,yshift=-3.15cm]
  \fill[scevent,opacity=0.14] (0.824,0) rectangle (1.171,2.000);
  \draw[scgray,solid,line width=0.4pt] (0.000,0.000) -- (0.014,0.317) -- (0.028,0.317) -- (0.042,0.000) -- (0.056,0.000) -- (0.069,0.317) -- (0.083,0.317) -- (0.097,0.000) -- (0.111,0.000) -- (0.125,0.317) -- (0.139,0.317) -- (0.153,0.000) -- (0.167,0.000) -- (0.181,0.317) -- (0.195,0.317) -- (0.208,0.000) -- (0.222,0.000) -- (0.236,0.317) -- (0.250,0.317) -- (0.264,0.000) -- (0.278,0.000) -- (0.292,0.317) -- (0.306,0.317) -- (0.320,0.000) -- (0.334,0.000) -- (0.347,0.317) -- (0.361,0.317) -- (0.375,0.000) -- (0.389,0.000) -- (0.403,0.317) -- (0.417,0.317) -- (0.431,0.000) -- (0.445,0.000) -- (0.459,0.317) -- (0.473,0.317) -- (0.486,0.000) -- (0.500,0.000) -- (0.514,0.317) -- (0.528,0.317) -- (0.542,0.000) -- (0.556,0.000) -- (0.570,0.317) -- (0.584,0.317) -- (0.598,0.000) -- (0.612,0.000) -- (0.625,0.317) -- (0.639,0.317) -- (0.653,0.000) -- (0.667,0.000) -- (0.681,0.317) -- (0.695,0.317) -- (0.709,0.000) -- (0.723,0.000) -- (0.737,0.317) -- (0.751,0.317) -- (0.764,0.000) -- (0.778,0.000) -- (0.792,0.317) -- (0.806,0.317) -- (0.820,0.000) -- (0.834,1.905) -- (0.848,0.317) -- (0.862,1.587) -- (0.876,0.000) -- (0.890,1.905) -- (0.903,0.317) -- (0.917,1.587) -- (0.931,0.000) -- (0.945,1.905) -- (0.959,0.317) -- (0.973,1.587) -- (0.987,0.000) -- (1.001,1.905) -- (1.015,0.317) -- (1.029,1.587) -- (1.042,0.000) -- (1.056,1.905) -- (1.070,0.317) -- (1.084,1.587) -- (1.098,0.000) -- (1.112,1.905) -- (1.126,0.317) -- (1.140,1.587) -- (1.154,0.000) -- (1.168,1.905) -- (1.181,0.317) -- (1.195,0.317) -- (1.209,0.000) -- (1.223,0.000) -- (1.237,0.317) -- (1.251,0.317) -- (1.265,0.000) -- (1.279,0.000) -- (1.293,0.317) -- (1.307,0.317) -- (1.320,0.000) -- (1.334,0.000) -- (1.348,0.317) -- (1.362,0.317) -- (1.376,0.000) -- (1.390,0.000) -- (1.404,0.317) -- (1.418,0.317) -- (1.432,0.000) -- (1.446,0.000) -- (1.459,0.317) -- (1.473,0.317) -- (1.487,0.000) -- (1.501,0.000) -- (1.515,0.317) -- (1.529,0.317) -- (1.543,0.000) -- (1.557,0.000) -- (1.571,0.317) -- (1.585,0.317) -- (1.598,0.000) -- (1.612,0.000) -- (1.626,0.317) -- (1.640,0.317) -- (1.654,0.000) -- (1.668,0.000) -- (1.682,0.317) -- (1.696,0.317) -- (1.710,0.000) -- (1.724,0.000) -- (1.737,0.317) -- (1.751,0.317) -- (1.765,0.000) -- (1.779,0.000) -- (1.793,0.317) -- (1.807,0.317) -- (1.821,0.000) -- (1.835,0.000) -- (1.849,0.317) -- (1.863,0.317) -- (1.876,0.000) -- (1.890,0.000) -- (1.904,0.317) -- (1.918,0.317) -- (1.932,0.000) -- (1.946,0.000) -- (1.960,0.317) -- (1.974,0.317) -- (1.988,0.000) -- (2.002,0.000) -- (2.015,0.317) -- (2.029,0.317) -- (2.043,0.000) -- (2.057,0.000) -- (2.071,0.317) -- (2.085,0.317) -- (2.099,0.000) -- (2.113,0.000) -- (2.127,0.317) -- (2.141,0.317) -- (2.154,0.000) -- (2.168,0.000) -- (2.182,0.317) -- (2.196,0.317) -- (2.210,0.000) -- (2.224,0.000) -- (2.238,0.317) -- (2.252,0.317) -- (2.266,0.000) -- (2.280,0.000) -- (2.293,0.317) -- (2.307,0.317) -- (2.321,0.000) -- (2.335,0.000) -- (2.349,0.317) -- (2.363,0.317) -- (2.377,0.000) -- (2.391,0.000) -- (2.405,0.317) -- (2.419,0.317) -- (2.432,0.000) -- (2.446,0.000) -- (2.460,0.317) -- (2.474,0.317) -- (2.488,0.000) -- (2.502,0.000) -- (2.516,0.317) -- (2.530,0.317) -- (2.544,0.000) -- (2.558,0.000) -- (2.571,0.317) -- (2.585,0.317) -- (2.599,0.000) -- (2.613,0.000) -- (2.627,0.317) -- (2.641,0.317) -- (2.655,0.000) -- (2.669,0.000) -- (2.683,0.317) -- (2.697,0.317) -- (2.710,0.000) -- (2.724,0.000) -- (2.738,0.317) -- (2.752,0.317) -- (2.766,0.000) -- (2.780,0.000) -- (2.794,0.317) -- (2.808,0.317) -- (2.822,0.000) -- (2.836,0.000) -- (2.849,0.317) -- (2.863,0.317) -- (2.877,0.000) -- (2.891,0.000) -- (2.905,0.317) -- (2.919,0.317) -- (2.933,0.000) -- (2.947,0.000) -- (2.961,0.317) -- (2.975,0.317) -- (2.988,0.000) -- (3.002,0.000) -- (3.016,0.317) -- (3.030,0.317) -- (3.044,0.000) -- (3.058,0.000) -- (3.072,0.317) -- (3.086,0.317) -- (3.100,0.000) -- (3.114,0.000) -- (3.127,0.317) -- (3.141,0.317) -- (3.155,0.000) -- (3.169,0.000) -- (3.183,0.317) -- (3.197,0.317) -- (3.211,0.000) -- (3.225,0.000) -- (3.239,0.317) -- (3.253,0.317) -- (3.266,0.000) -- (3.280,0.000) -- (3.294,0.317) -- (3.308,0.317) -- (3.322,0.000) -- (3.336,0.000) -- (3.350,0.317) -- (3.364,0.317) -- (3.378,0.000) -- (3.392,0.000) -- (3.405,0.317) -- (3.419,0.317) -- (3.433,0.000) -- (3.447,0.000) -- (3.461,0.317) -- (3.475,0.317) -- (3.489,0.000) -- (3.503,0.000) -- (3.517,0.317) -- (3.531,0.317) -- (3.544,0.000) -- (3.558,0.000) -- (3.572,0.317) -- (3.586,0.317) -- (3.600,0.000);
  \draw[scblack,solid,line width=0.7pt] (0,0.317) -- (3.600,0.317);
  \draw[scblue,dotted,line width=0.7pt] (0.000,0.317) -- (0.014,0.000) -- (0.028,0.317) -- (0.042,0.317) -- (0.056,0.317) -- (0.069,0.317) -- (0.083,0.317) -- (0.097,0.317) -- (0.111,0.317) -- (0.125,0.317) -- (0.139,0.317) -- (0.153,0.317) -- (0.167,0.317) -- (0.181,0.317) -- (0.195,0.317) -- (0.208,0.317) -- (0.222,0.317) -- (0.236,0.317) -- (0.250,0.317) -- (0.264,0.317) -- (0.278,0.317) -- (0.292,0.317) -- (0.306,0.317) -- (0.320,0.317) -- (0.334,0.317) -- (0.347,0.317) -- (0.361,0.317) -- (0.375,0.317) -- (0.389,0.317) -- (0.403,0.317) -- (0.417,0.317) -- (0.431,0.317) -- (0.445,0.317) -- (0.459,0.317) -- (0.473,0.317) -- (0.486,0.317) -- (0.500,0.317) -- (0.514,0.317) -- (0.528,0.317) -- (0.542,0.317) -- (0.556,0.317) -- (0.570,0.317) -- (0.584,0.317) -- (0.598,0.317) -- (0.612,0.317) -- (0.625,0.317) -- (0.639,0.317) -- (0.653,0.317) -- (0.667,0.317) -- (0.681,0.317) -- (0.695,0.317) -- (0.709,0.317) -- (0.723,0.317) -- (0.737,0.317) -- (0.751,0.317) -- (0.764,0.317) -- (0.778,0.317) -- (0.792,0.317) -- (0.806,0.317) -- (0.820,0.317) -- (0.834,0.317) -- (0.848,1.905) -- (0.862,1.905) -- (0.876,1.905) -- (0.890,1.905) -- (0.903,1.905) -- (0.917,1.905) -- (0.931,1.905) -- (0.945,1.905) -- (0.959,1.905) -- (0.973,1.905) -- (0.987,1.905) -- (1.001,1.905) -- (1.015,1.905) -- (1.029,1.905) -- (1.042,1.905) -- (1.056,1.905) -- (1.070,1.905) -- (1.084,1.905) -- (1.098,1.905) -- (1.112,1.905) -- (1.126,1.905) -- (1.140,1.905) -- (1.154,1.905) -- (1.168,1.905) -- (1.181,1.905) -- (1.195,1.905) -- (1.209,1.905) -- (1.223,1.905) -- (1.237,1.905) -- (1.251,1.905) -- (1.265,1.905) -- (1.279,1.905) -- (1.293,1.905) -- (1.307,1.905) -- (1.320,0.317) -- (1.334,0.317) -- (1.348,0.317) -- (1.362,0.317) -- (1.376,0.317) -- (1.390,0.317) -- (1.404,0.317) -- (1.418,0.317) -- (1.432,0.317) -- (1.446,0.317) -- (1.459,0.317) -- (1.473,0.317) -- (1.487,0.317) -- (1.501,0.317) -- (1.515,0.317) -- (1.529,0.317) -- (1.543,0.317) -- (1.557,0.317) -- (1.571,0.317) -- (1.585,0.317) -- (1.598,0.317) -- (1.612,0.317) -- (1.626,0.317) -- (1.640,0.317) -- (1.654,0.317) -- (1.668,0.317) -- (1.682,0.317) -- (1.696,0.317) -- (1.710,0.317) -- (1.724,0.317) -- (1.737,0.317) -- (1.751,0.317) -- (1.765,0.317) -- (1.779,0.317) -- (1.793,0.317) -- (1.807,0.317) -- (1.821,0.317) -- (1.835,0.317) -- (1.849,0.317) -- (1.863,0.317) -- (1.876,0.317) -- (1.890,0.317) -- (1.904,0.317) -- (1.918,0.317) -- (1.932,0.317) -- (1.946,0.317) -- (1.960,0.317) -- (1.974,0.317) -- (1.988,0.317) -- (2.002,0.317) -- (2.015,0.317) -- (2.029,0.317) -- (2.043,0.317) -- (2.057,0.317) -- (2.071,0.317) -- (2.085,0.317) -- (2.099,0.317) -- (2.113,0.317) -- (2.127,0.317) -- (2.141,0.317) -- (2.154,0.317) -- (2.168,0.317) -- (2.182,0.317) -- (2.196,0.317) -- (2.210,0.317) -- (2.224,0.317) -- (2.238,0.317) -- (2.252,0.317) -- (2.266,0.317) -- (2.280,0.317) -- (2.293,0.317) -- (2.307,0.317) -- (2.321,0.317) -- (2.335,0.317) -- (2.349,0.317) -- (2.363,0.317) -- (2.377,0.317) -- (2.391,0.317) -- (2.405,0.317) -- (2.419,0.317) -- (2.432,0.317) -- (2.446,0.317) -- (2.460,0.317) -- (2.474,0.317) -- (2.488,0.317) -- (2.502,0.317) -- (2.516,0.317) -- (2.530,0.317) -- (2.544,0.317) -- (2.558,0.317) -- (2.571,0.317) -- (2.585,0.317) -- (2.599,0.317) -- (2.613,0.317) -- (2.627,0.317) -- (2.641,0.317) -- (2.655,0.317) -- (2.669,0.317) -- (2.683,0.317) -- (2.697,0.317) -- (2.710,0.317) -- (2.724,0.317) -- (2.738,0.317) -- (2.752,0.317) -- (2.766,0.317) -- (2.780,0.317) -- (2.794,0.317) -- (2.808,0.317) -- (2.822,0.317) -- (2.836,0.317) -- (2.849,0.317) -- (2.863,0.317) -- (2.877,0.317) -- (2.891,0.317) -- (2.905,0.317) -- (2.919,0.317) -- (2.933,0.317) -- (2.947,0.317) -- (2.961,0.317) -- (2.975,0.317) -- (2.988,0.317) -- (3.002,0.317) -- (3.016,0.317) -- (3.030,0.317) -- (3.044,0.317) -- (3.058,0.317) -- (3.072,0.317) -- (3.086,0.317) -- (3.100,0.317) -- (3.114,0.317) -- (3.127,0.317) -- (3.141,0.317) -- (3.155,0.317) -- (3.169,0.317) -- (3.183,0.317) -- (3.197,0.317) -- (3.211,0.317) -- (3.225,0.317) -- (3.239,0.317) -- (3.253,0.317) -- (3.266,0.317) -- (3.280,0.317) -- (3.294,0.317) -- (3.308,0.317) -- (3.322,0.317) -- (3.336,0.317) -- (3.350,0.317) -- (3.364,0.317) -- (3.378,0.317) -- (3.392,0.317) -- (3.405,0.317) -- (3.419,0.317) -- (3.433,0.317) -- (3.447,0.317) -- (3.461,0.317) -- (3.475,0.317) -- (3.489,0.317) -- (3.503,0.317) -- (3.517,0.317) -- (3.531,0.317) -- (3.544,0.317) -- (3.558,0.317) -- (3.572,0.317) -- (3.586,0.317) -- (3.600,0.317);
  \draw[scgreen,densely dashed,line width=1.0pt] (0.000,0.317) -- (0.014,0.317) -- (0.028,0.317) -- (0.042,0.317) -- (0.056,0.317) -- (0.069,0.317) -- (0.083,0.317) -- (0.097,0.317) -- (0.111,0.317) -- (0.125,0.317) -- (0.139,0.317) -- (0.153,0.317) -- (0.167,0.317) -- (0.181,0.317) -- (0.195,0.317) -- (0.208,0.317) -- (0.222,0.317) -- (0.236,0.317) -- (0.250,0.317) -- (0.264,0.317) -- (0.278,0.317) -- (0.292,0.317) -- (0.306,0.317) -- (0.320,0.317) -- (0.334,0.317) -- (0.347,0.317) -- (0.361,0.317) -- (0.375,0.317) -- (0.389,0.317) -- (0.403,0.317) -- (0.417,0.317) -- (0.431,0.317) -- (0.445,0.317) -- (0.459,0.317) -- (0.473,0.317) -- (0.486,0.317) -- (0.500,0.317) -- (0.514,0.317) -- (0.528,0.317) -- (0.542,0.317) -- (0.556,0.317) -- (0.570,0.317) -- (0.584,0.317) -- (0.598,0.317) -- (0.612,0.317) -- (0.625,0.317) -- (0.639,0.317) -- (0.653,0.317) -- (0.667,0.317) -- (0.681,0.317) -- (0.695,0.317) -- (0.709,0.317) -- (0.723,0.317) -- (0.737,0.317) -- (0.751,0.317) -- (0.764,0.317) -- (0.778,0.317) -- (0.792,0.317) -- (0.806,0.317) -- (0.820,0.317) -- (0.834,0.317) -- (0.848,0.317) -- (0.862,0.317) -- (0.876,0.317) -- (0.890,0.317) -- (0.903,0.317) -- (0.917,0.317) -- (0.931,0.317) -- (0.945,0.317) -- (0.959,0.317) -- (0.973,0.317) -- (0.987,0.317) -- (1.001,0.317) -- (1.015,0.317) -- (1.029,0.317) -- (1.042,0.317) -- (1.056,0.317) -- (1.070,0.317) -- (1.084,0.317) -- (1.098,0.317) -- (1.112,0.317) -- (1.126,0.317) -- (1.140,0.317) -- (1.154,0.317) -- (1.168,0.317) -- (1.181,0.317) -- (1.195,0.317) -- (1.209,0.317) -- (1.223,0.317) -- (1.237,0.317) -- (1.251,0.317) -- (1.265,0.317) -- (1.279,0.317) -- (1.293,0.317) -- (1.307,0.317) -- (1.320,0.317) -- (1.334,0.317) -- (1.348,0.317) -- (1.362,0.317) -- (1.376,0.317) -- (1.390,0.317) -- (1.404,0.317) -- (1.418,0.317) -- (1.432,0.317) -- (1.446,0.317) -- (1.459,0.317) -- (1.473,0.317) -- (1.487,0.317) -- (1.501,0.317) -- (1.515,0.317) -- (1.529,0.317) -- (1.543,0.317) -- (1.557,0.317) -- (1.571,0.317) -- (1.585,0.317) -- (1.598,0.317) -- (1.612,0.317) -- (1.626,0.317) -- (1.640,0.317) -- (1.654,0.317) -- (1.668,0.317) -- (1.682,0.317) -- (1.696,0.317) -- (1.710,0.317) -- (1.724,0.317) -- (1.737,0.317) -- (1.751,0.317) -- (1.765,0.317) -- (1.779,0.317) -- (1.793,0.317) -- (1.807,0.317) -- (1.821,0.317) -- (1.835,0.317) -- (1.849,0.317) -- (1.863,0.317) -- (1.876,0.317) -- (1.890,0.317) -- (1.904,0.317) -- (1.918,0.317) -- (1.932,0.317) -- (1.946,0.317) -- (1.960,0.317) -- (1.974,0.317) -- (1.988,0.317) -- (2.002,0.317) -- (2.015,0.317) -- (2.029,0.317) -- (2.043,0.317) -- (2.057,0.317) -- (2.071,0.317) -- (2.085,0.317) -- (2.099,0.317) -- (2.113,0.317) -- (2.127,0.317) -- (2.141,0.317) -- (2.154,0.317) -- (2.168,0.317) -- (2.182,0.317) -- (2.196,0.317) -- (2.210,0.317) -- (2.224,0.317) -- (2.238,0.317) -- (2.252,0.317) -- (2.266,0.317) -- (2.280,0.317) -- (2.293,0.317) -- (2.307,0.317) -- (2.321,0.317) -- (2.335,0.317) -- (2.349,0.317) -- (2.363,0.317) -- (2.377,0.317) -- (2.391,0.317) -- (2.405,0.317) -- (2.419,0.317) -- (2.432,0.317) -- (2.446,0.317) -- (2.460,0.317) -- (2.474,0.317) -- (2.488,0.317) -- (2.502,0.317) -- (2.516,0.317) -- (2.530,0.317) -- (2.544,0.317) -- (2.558,0.317) -- (2.571,0.317) -- (2.585,0.317) -- (2.599,0.317) -- (2.613,0.317) -- (2.627,0.317) -- (2.641,0.317) -- (2.655,0.317) -- (2.669,0.317) -- (2.683,0.317) -- (2.697,0.317) -- (2.710,0.317) -- (2.724,0.317) -- (2.738,0.317) -- (2.752,0.317) -- (2.766,0.317) -- (2.780,0.317) -- (2.794,0.317) -- (2.808,0.317) -- (2.822,0.317) -- (2.836,0.317) -- (2.849,0.317) -- (2.863,0.317) -- (2.877,0.317) -- (2.891,0.317) -- (2.905,0.317) -- (2.919,0.317) -- (2.933,0.317) -- (2.947,0.317) -- (2.961,0.317) -- (2.975,0.317) -- (2.988,0.317) -- (3.002,0.317) -- (3.016,0.317) -- (3.030,0.317) -- (3.044,0.317) -- (3.058,0.317) -- (3.072,0.317) -- (3.086,0.317) -- (3.100,0.317) -- (3.114,0.317) -- (3.127,0.317) -- (3.141,0.317) -- (3.155,0.317) -- (3.169,0.317) -- (3.183,0.317) -- (3.197,0.317) -- (3.211,0.317) -- (3.225,0.317) -- (3.239,0.317) -- (3.253,0.317) -- (3.266,0.317) -- (3.280,0.317) -- (3.294,0.317) -- (3.308,0.317) -- (3.322,0.317) -- (3.336,0.317) -- (3.350,0.317) -- (3.364,0.317) -- (3.378,0.317) -- (3.392,0.317) -- (3.405,0.317) -- (3.419,0.317) -- (3.433,0.317) -- (3.447,0.317) -- (3.461,0.317) -- (3.475,0.317) -- (3.489,0.317) -- (3.503,0.317) -- (3.517,0.317) -- (3.531,0.317) -- (3.544,0.317) -- (3.558,0.317) -- (3.572,0.317) -- (3.586,0.317) -- (3.600,0.317);
  \draw[scred,line width=0.5pt] (0.945,0) -- (0.945,0.16);
  \draw[scred,line width=0.5pt] (0.959,0) -- (0.959,0.16);
  \draw[scred,line width=0.5pt] (0.973,0) -- (0.973,0.16);
  \draw[scred,line width=0.5pt] (0.987,0) -- (0.987,0.16);
  \draw[scred,line width=0.5pt] (1.001,0) -- (1.001,0.16);
  \draw[scred,line width=0.5pt] (1.015,0) -- (1.015,0.16);
  \draw[scred,line width=0.5pt] (1.029,0) -- (1.029,0.16);
  \draw[scred,line width=0.5pt] (1.042,0) -- (1.042,0.16);
  \draw[scred,line width=0.5pt] (1.056,0) -- (1.056,0.16);
  \draw[scred,line width=0.5pt] (1.070,0) -- (1.070,0.16);
  \draw[scblack] (0,0) rectangle (3.600,2.000);
  \draw[scblack] (0.000,0) -- (0.000,-0.08);
  \node[font=\scriptsize,below] at (0.000,-0.08) {563};
  \draw[scblack] (3.600,0) -- (3.600,-0.08);
  \node[font=\scriptsize,below] at (3.600,-0.08) {1599};
  \draw[scblack] (0,0.000) -- (-0.08,0.000);
  \node[font=\scriptsize,left] at (-0.10,0.000) {0};
  \draw[scblack] (0,0.317) -- (-0.08,0.317);
  \node[font=\scriptsize,left] at (-0.10,0.317) {$\theta_0$};
  \draw[scblack] (0,1.990) -- (-0.08,1.990);
  \node[font=\scriptsize,left] at (-0.10,1.990) {6.7};
  \node[font=\scriptsize,below=0.42cm] at (1.800,0) {sample index};
  \node[font={\scriptsize\bfseries},above] at (1.800,2.040) {Y5};
  \end{scope}
  \begin{scope}[xshift=9.40cm,yshift=-3.15cm]
  \draw[scgray,solid,line width=0.4pt] (0.000,0.000) -- (0.014,0.317) -- (0.028,0.317) -- (0.042,0.000) -- (0.056,0.000) -- (0.069,0.317) -- (0.083,0.317) -- (0.097,0.000) -- (0.111,0.000) -- (0.125,0.317) -- (0.139,0.317) -- (0.153,0.000) -- (0.167,0.000) -- (0.181,0.317) -- (0.195,0.317) -- (0.208,0.000) -- (0.222,0.000) -- (0.236,0.317) -- (0.250,0.317) -- (0.264,0.000) -- (0.278,0.000) -- (0.292,0.317) -- (0.306,0.317) -- (0.320,0.000) -- (0.334,0.000) -- (0.347,0.317) -- (0.361,0.317) -- (0.375,0.000) -- (0.389,0.000) -- (0.403,0.317) -- (0.417,0.317) -- (0.431,0.000) -- (0.445,0.000) -- (0.459,0.317) -- (0.473,0.317) -- (0.486,0.000) -- (0.500,0.000) -- (0.514,0.317) -- (0.528,0.317) -- (0.542,0.000) -- (0.556,0.000) -- (0.570,0.317) -- (0.584,0.317) -- (0.598,0.000) -- (0.612,0.000) -- (0.625,0.317) -- (0.639,0.317) -- (0.653,0.000) -- (0.667,0.000) -- (0.681,0.317) -- (0.695,0.317) -- (0.709,0.000) -- (0.723,0.000) -- (0.737,0.317) -- (0.751,0.317) -- (0.764,0.000) -- (0.778,0.000) -- (0.792,0.317) -- (0.806,0.317) -- (0.820,0.000) -- (0.834,1.587) -- (0.848,1.905) -- (0.862,1.270) -- (0.876,1.587) -- (0.890,1.587) -- (0.903,1.905) -- (0.917,1.270) -- (0.931,1.587) -- (0.945,1.587) -- (0.959,1.905) -- (0.973,1.270) -- (0.987,1.587) -- (1.001,1.587) -- (1.015,1.905) -- (1.029,1.270) -- (1.042,1.587) -- (1.056,1.587) -- (1.070,1.905) -- (1.084,1.270) -- (1.098,1.587) -- (1.112,1.587) -- (1.126,1.905) -- (1.140,1.270) -- (1.154,1.587) -- (1.168,1.587) -- (1.181,1.905) -- (1.195,1.270) -- (1.209,1.587) -- (1.223,1.587) -- (1.237,1.905) -- (1.251,1.270) -- (1.265,1.587) -- (1.279,1.587) -- (1.293,1.905) -- (1.307,1.270) -- (1.320,1.587) -- (1.334,1.587) -- (1.348,1.905) -- (1.362,1.270) -- (1.376,1.587) -- (1.390,1.587) -- (1.404,1.905) -- (1.418,1.270) -- (1.432,1.587) -- (1.446,1.587) -- (1.459,1.905) -- (1.473,1.270) -- (1.487,1.587) -- (1.501,1.587) -- (1.515,1.905) -- (1.529,1.270) -- (1.543,1.587) -- (1.557,1.587) -- (1.571,1.905) -- (1.585,1.270) -- (1.598,1.587) -- (1.612,1.587) -- (1.626,1.905) -- (1.640,1.270) -- (1.654,1.587) -- (1.668,1.587) -- (1.682,1.905) -- (1.696,1.270) -- (1.710,1.587) -- (1.724,1.587) -- (1.737,1.905) -- (1.751,1.270) -- (1.765,1.587) -- (1.779,1.587) -- (1.793,1.905) -- (1.807,1.270) -- (1.821,1.587) -- (1.835,1.587) -- (1.849,1.905) -- (1.863,1.270) -- (1.876,1.587) -- (1.890,1.587) -- (1.904,1.905) -- (1.918,1.270) -- (1.932,1.587) -- (1.946,1.587) -- (1.960,1.905) -- (1.974,1.270) -- (1.988,1.587) -- (2.002,1.587) -- (2.015,1.905) -- (2.029,1.270) -- (2.043,1.587) -- (2.057,1.587) -- (2.071,1.905) -- (2.085,1.270) -- (2.099,1.587) -- (2.113,1.587) -- (2.127,1.905) -- (2.141,1.270) -- (2.154,1.587) -- (2.168,1.587) -- (2.182,1.905) -- (2.196,1.270) -- (2.210,1.587) -- (2.224,1.587) -- (2.238,1.905) -- (2.252,1.270) -- (2.266,1.587) -- (2.280,1.587) -- (2.293,1.905) -- (2.307,1.270) -- (2.321,1.587) -- (2.335,1.587) -- (2.349,1.905) -- (2.363,1.270) -- (2.377,1.587) -- (2.391,1.587) -- (2.405,1.905) -- (2.419,1.270) -- (2.432,1.587) -- (2.446,1.587) -- (2.460,1.905) -- (2.474,1.270) -- (2.488,1.587) -- (2.502,1.587) -- (2.516,1.905) -- (2.530,1.270) -- (2.544,1.587) -- (2.558,1.587) -- (2.571,1.905) -- (2.585,1.270) -- (2.599,1.587) -- (2.613,1.587) -- (2.627,1.905) -- (2.641,1.270) -- (2.655,1.587) -- (2.669,1.587) -- (2.683,1.905) -- (2.697,1.270) -- (2.710,1.587) -- (2.724,1.587) -- (2.738,1.905) -- (2.752,1.270) -- (2.766,1.587) -- (2.780,1.587) -- (2.794,1.905) -- (2.808,1.270) -- (2.822,1.587) -- (2.836,1.587) -- (2.849,1.905) -- (2.863,1.270) -- (2.877,1.587) -- (2.891,1.587) -- (2.905,1.905) -- (2.919,1.270) -- (2.933,1.587) -- (2.947,1.587) -- (2.961,1.905) -- (2.975,1.270) -- (2.988,1.587) -- (3.002,1.587) -- (3.016,1.905) -- (3.030,1.270) -- (3.044,1.587) -- (3.058,1.587) -- (3.072,1.905) -- (3.086,1.270) -- (3.100,1.587) -- (3.114,1.587) -- (3.127,1.905) -- (3.141,1.270) -- (3.155,1.587) -- (3.169,1.587) -- (3.183,1.905) -- (3.197,1.270) -- (3.211,1.587) -- (3.225,1.587) -- (3.239,1.905) -- (3.253,1.270) -- (3.266,1.587) -- (3.280,1.587) -- (3.294,1.905) -- (3.308,1.270) -- (3.322,1.587) -- (3.336,1.587) -- (3.350,1.905) -- (3.364,1.270) -- (3.378,1.587) -- (3.392,1.587) -- (3.405,1.905) -- (3.419,1.270) -- (3.433,1.587) -- (3.447,1.587) -- (3.461,1.905) -- (3.475,1.270) -- (3.489,1.587) -- (3.503,1.587) -- (3.517,1.905) -- (3.531,1.270) -- (3.544,1.587) -- (3.558,1.587) -- (3.572,1.905) -- (3.586,1.270) -- (3.600,1.587);
  \draw[scblack,solid,line width=0.7pt] (0,0.317) -- (3.600,0.317);
  \draw[scblue,dotted,line width=0.7pt] (0.000,0.317) -- (0.014,0.000) -- (0.028,0.317) -- (0.042,0.317) -- (0.056,0.317) -- (0.069,0.317) -- (0.083,0.317) -- (0.097,0.317) -- (0.111,0.317) -- (0.125,0.317) -- (0.139,0.317) -- (0.153,0.317) -- (0.167,0.317) -- (0.181,0.317) -- (0.195,0.317) -- (0.208,0.317) -- (0.222,0.317) -- (0.236,0.317) -- (0.250,0.317) -- (0.264,0.317) -- (0.278,0.317) -- (0.292,0.317) -- (0.306,0.317) -- (0.320,0.317) -- (0.334,0.317) -- (0.347,0.317) -- (0.361,0.317) -- (0.375,0.317) -- (0.389,0.317) -- (0.403,0.317) -- (0.417,0.317) -- (0.431,0.317) -- (0.445,0.317) -- (0.459,0.317) -- (0.473,0.317) -- (0.486,0.317) -- (0.500,0.317) -- (0.514,0.317) -- (0.528,0.317) -- (0.542,0.317) -- (0.556,0.317) -- (0.570,0.317) -- (0.584,0.317) -- (0.598,0.317) -- (0.612,0.317) -- (0.625,0.317) -- (0.639,0.317) -- (0.653,0.317) -- (0.667,0.317) -- (0.681,0.317) -- (0.695,0.317) -- (0.709,0.317) -- (0.723,0.317) -- (0.737,0.317) -- (0.751,0.317) -- (0.764,0.317) -- (0.778,0.317) -- (0.792,0.317) -- (0.806,0.317) -- (0.820,0.317) -- (0.834,0.317) -- (0.848,1.587) -- (0.862,1.905) -- (0.876,1.905) -- (0.890,1.905) -- (0.903,1.905) -- (0.917,1.905) -- (0.931,1.905) -- (0.945,1.905) -- (0.959,1.905) -- (0.973,1.905) -- (0.987,1.905) -- (1.001,1.905) -- (1.015,1.905) -- (1.029,1.905) -- (1.042,1.905) -- (1.056,1.905) -- (1.070,1.905) -- (1.084,1.905) -- (1.098,1.905) -- (1.112,1.905) -- (1.126,1.905) -- (1.140,1.905) -- (1.154,1.905) -- (1.168,1.905) -- (1.181,1.905) -- (1.195,1.905) -- (1.209,1.905) -- (1.223,1.905) -- (1.237,1.905) -- (1.251,1.905) -- (1.265,1.905) -- (1.279,1.905) -- (1.293,1.905) -- (1.307,1.905) -- (1.320,1.905) -- (1.334,1.905) -- (1.348,1.905) -- (1.362,1.905) -- (1.376,1.905) -- (1.390,1.905) -- (1.404,1.905) -- (1.418,1.905) -- (1.432,1.905) -- (1.446,1.905) -- (1.459,1.905) -- (1.473,1.905) -- (1.487,1.905) -- (1.501,1.905) -- (1.515,1.905) -- (1.529,1.905) -- (1.543,1.905) -- (1.557,1.905) -- (1.571,1.905) -- (1.585,1.905) -- (1.598,1.905) -- (1.612,1.905) -- (1.626,1.905) -- (1.640,1.905) -- (1.654,1.905) -- (1.668,1.905) -- (1.682,1.905) -- (1.696,1.905) -- (1.710,1.905) -- (1.724,1.905) -- (1.737,1.905) -- (1.751,1.905) -- (1.765,1.905) -- (1.779,1.905) -- (1.793,1.905) -- (1.807,1.905) -- (1.821,1.905) -- (1.835,1.905) -- (1.849,1.905) -- (1.863,1.905) -- (1.876,1.905) -- (1.890,1.905) -- (1.904,1.905) -- (1.918,1.905) -- (1.932,1.905) -- (1.946,1.905) -- (1.960,1.905) -- (1.974,1.905) -- (1.988,1.905) -- (2.002,1.905) -- (2.015,1.905) -- (2.029,1.905) -- (2.043,1.905) -- (2.057,1.905) -- (2.071,1.905) -- (2.085,1.905) -- (2.099,1.905) -- (2.113,1.905) -- (2.127,1.905) -- (2.141,1.905) -- (2.154,1.905) -- (2.168,1.905) -- (2.182,1.905) -- (2.196,1.905) -- (2.210,1.905) -- (2.224,1.905) -- (2.238,1.905) -- (2.252,1.905) -- (2.266,1.905) -- (2.280,1.905) -- (2.293,1.905) -- (2.307,1.905) -- (2.321,1.905) -- (2.335,1.905) -- (2.349,1.905) -- (2.363,1.905) -- (2.377,1.905) -- (2.391,1.905) -- (2.405,1.905) -- (2.419,1.905) -- (2.432,1.905) -- (2.446,1.905) -- (2.460,1.905) -- (2.474,1.905) -- (2.488,1.905) -- (2.502,1.905) -- (2.516,1.905) -- (2.530,1.905) -- (2.544,1.905) -- (2.558,1.905) -- (2.571,1.905) -- (2.585,1.905) -- (2.599,1.905) -- (2.613,1.905) -- (2.627,1.905) -- (2.641,1.905) -- (2.655,1.905) -- (2.669,1.905) -- (2.683,1.905) -- (2.697,1.905) -- (2.710,1.905) -- (2.724,1.905) -- (2.738,1.905) -- (2.752,1.905) -- (2.766,1.905) -- (2.780,1.905) -- (2.794,1.905) -- (2.808,1.905) -- (2.822,1.905) -- (2.836,1.905) -- (2.849,1.905) -- (2.863,1.905) -- (2.877,1.905) -- (2.891,1.905) -- (2.905,1.905) -- (2.919,1.905) -- (2.933,1.905) -- (2.947,1.905) -- (2.961,1.905) -- (2.975,1.905) -- (2.988,1.905) -- (3.002,1.905) -- (3.016,1.905) -- (3.030,1.905) -- (3.044,1.905) -- (3.058,1.905) -- (3.072,1.905) -- (3.086,1.905) -- (3.100,1.905) -- (3.114,1.905) -- (3.127,1.905) -- (3.141,1.905) -- (3.155,1.905) -- (3.169,1.905) -- (3.183,1.905) -- (3.197,1.905) -- (3.211,1.905) -- (3.225,1.905) -- (3.239,1.905) -- (3.253,1.905) -- (3.266,1.905) -- (3.280,1.905) -- (3.294,1.905) -- (3.308,1.905) -- (3.322,1.905) -- (3.336,1.905) -- (3.350,1.905) -- (3.364,1.905) -- (3.378,1.905) -- (3.392,1.905) -- (3.405,1.905) -- (3.419,1.905) -- (3.433,1.905) -- (3.447,1.905) -- (3.461,1.905) -- (3.475,1.905) -- (3.489,1.905) -- (3.503,1.905) -- (3.517,1.905) -- (3.531,1.905) -- (3.544,1.905) -- (3.558,1.905) -- (3.572,1.905) -- (3.586,1.905) -- (3.600,1.905);
  \draw[scgreen,densely dashed,line width=1.0pt] (0.000,0.317) -- (0.014,0.317) -- (0.028,0.317) -- (0.042,0.317) -- (0.056,0.317) -- (0.069,0.317) -- (0.083,0.317) -- (0.097,0.317) -- (0.111,0.317) -- (0.125,0.317) -- (0.139,0.317) -- (0.153,0.317) -- (0.167,0.317) -- (0.181,0.317) -- (0.195,0.317) -- (0.208,0.317) -- (0.222,0.317) -- (0.236,0.317) -- (0.250,0.317) -- (0.264,0.317) -- (0.278,0.317) -- (0.292,0.317) -- (0.306,0.317) -- (0.320,0.317) -- (0.334,0.317) -- (0.347,0.317) -- (0.361,0.317) -- (0.375,0.317) -- (0.389,0.317) -- (0.403,0.317) -- (0.417,0.317) -- (0.431,0.317) -- (0.445,0.317) -- (0.459,0.317) -- (0.473,0.317) -- (0.486,0.317) -- (0.500,0.317) -- (0.514,0.317) -- (0.528,0.317) -- (0.542,0.317) -- (0.556,0.317) -- (0.570,0.317) -- (0.584,0.317) -- (0.598,0.317) -- (0.612,0.317) -- (0.625,0.317) -- (0.639,0.317) -- (0.653,0.317) -- (0.667,0.317) -- (0.681,0.317) -- (0.695,0.317) -- (0.709,0.317) -- (0.723,0.317) -- (0.737,0.317) -- (0.751,0.317) -- (0.764,0.317) -- (0.778,0.317) -- (0.792,0.317) -- (0.806,0.317) -- (0.820,0.317) -- (0.834,0.317) -- (0.848,0.317) -- (0.862,0.317) -- (0.876,0.317) -- (0.890,0.317) -- (0.903,0.317) -- (0.917,0.317) -- (0.931,0.317) -- (0.945,0.317) -- (0.959,0.317) -- (0.973,0.317) -- (0.987,0.317) -- (1.001,0.317) -- (1.015,0.317) -- (1.029,0.317) -- (1.042,0.317) -- (1.056,0.317) -- (1.070,0.317) -- (1.084,0.317) -- (1.098,0.317) -- (1.112,0.317) -- (1.126,0.317) -- (1.140,0.317) -- (1.154,0.317) -- (1.168,0.317) -- (1.181,0.317) -- (1.195,0.317) -- (1.209,0.317) -- (1.223,0.317) -- (1.237,0.317) -- (1.251,0.317) -- (1.265,0.317) -- (1.279,0.317) -- (1.293,0.317) -- (1.307,0.317) -- (1.320,0.317) -- (1.334,0.317) -- (1.348,0.317) -- (1.362,0.317) -- (1.376,0.317) -- (1.390,0.317) -- (1.404,0.317) -- (1.418,0.317) -- (1.432,0.317) -- (1.446,0.317) -- (1.459,0.317) -- (1.473,0.317) -- (1.487,0.317) -- (1.501,0.317) -- (1.515,0.317) -- (1.529,0.317) -- (1.543,0.317) -- (1.557,0.317) -- (1.571,0.317) -- (1.585,0.317) -- (1.598,0.317) -- (1.612,0.317) -- (1.626,0.317) -- (1.640,0.317) -- (1.654,0.317) -- (1.668,0.317) -- (1.682,0.317) -- (1.696,0.317) -- (1.710,0.317) -- (1.724,0.317) -- (1.737,0.317) -- (1.751,0.317) -- (1.765,0.317) -- (1.779,0.317) -- (1.793,0.317) -- (1.807,0.317) -- (1.821,0.317) -- (1.835,0.317) -- (1.849,0.317) -- (1.863,0.317) -- (1.876,0.317) -- (1.890,0.317) -- (1.904,0.317) -- (1.918,0.317) -- (1.932,0.317) -- (1.946,0.317) -- (1.960,0.317) -- (1.974,0.317) -- (1.988,0.317) -- (2.002,0.317) -- (2.015,0.317) -- (2.029,0.317) -- (2.043,0.317) -- (2.057,0.317) -- (2.071,0.317) -- (2.085,0.317) -- (2.099,0.317) -- (2.113,0.317) -- (2.127,0.317) -- (2.141,0.317) -- (2.154,0.317) -- (2.168,0.317) -- (2.182,0.317) -- (2.196,0.317) -- (2.210,0.317) -- (2.224,0.317) -- (2.238,0.317) -- (2.252,0.317) -- (2.266,0.317) -- (2.280,0.317) -- (2.293,0.317) -- (2.307,0.317) -- (2.321,0.317) -- (2.335,0.317) -- (2.349,0.317) -- (2.363,0.317) -- (2.377,0.317) -- (2.391,0.317) -- (2.405,0.317) -- (2.419,0.317) -- (2.432,0.317) -- (2.446,0.317) -- (2.460,0.317) -- (2.474,0.317) -- (2.488,0.317) -- (2.502,0.317) -- (2.516,0.317) -- (2.530,0.317) -- (2.544,0.317) -- (2.558,0.317) -- (2.571,0.317) -- (2.585,0.317) -- (2.599,0.317) -- (2.613,0.317) -- (2.627,0.317) -- (2.641,0.317) -- (2.655,0.317) -- (2.669,0.317) -- (2.683,0.317) -- (2.697,0.317) -- (2.710,0.317) -- (2.724,0.317) -- (2.738,0.317) -- (2.752,0.317) -- (2.766,0.317) -- (2.780,0.317) -- (2.794,0.317) -- (2.808,0.317) -- (2.822,0.317) -- (2.836,0.317) -- (2.849,0.317) -- (2.863,0.317) -- (2.877,0.317) -- (2.891,0.317) -- (2.905,0.317) -- (2.919,0.317) -- (2.933,0.317) -- (2.947,0.317) -- (2.961,0.317) -- (2.975,0.317) -- (2.988,0.317) -- (3.002,0.317) -- (3.016,0.317) -- (3.030,0.317) -- (3.044,0.317) -- (3.058,0.317) -- (3.072,0.317) -- (3.086,0.317) -- (3.100,0.317) -- (3.114,0.317) -- (3.127,0.317) -- (3.141,0.317) -- (3.155,0.317) -- (3.169,0.317) -- (3.183,0.317) -- (3.197,0.317) -- (3.211,0.317) -- (3.225,0.317) -- (3.239,0.317) -- (3.253,0.317) -- (3.266,0.317) -- (3.280,0.317) -- (3.294,0.317) -- (3.308,0.317) -- (3.322,0.317) -- (3.336,0.317) -- (3.350,0.317) -- (3.364,0.317) -- (3.378,0.317) -- (3.392,0.317) -- (3.405,0.317) -- (3.419,0.317) -- (3.433,0.317) -- (3.447,0.317) -- (3.461,0.317) -- (3.475,0.317) -- (3.489,0.317) -- (3.503,0.317) -- (3.517,0.317) -- (3.531,0.317) -- (3.544,0.317) -- (3.558,0.317) -- (3.572,0.317) -- (3.586,0.317) -- (3.600,0.317);
  \draw[scred,line width=0.5pt] (0.890,0) -- (0.890,0.16);
  \draw[scred,line width=0.5pt] (0.903,0) -- (0.903,0.16);
  \draw[scred,line width=0.5pt] (0.917,0) -- (0.917,0.16);
  \draw[scred,line width=0.5pt] (0.931,0) -- (0.931,0.16);
  \draw[scred,line width=0.5pt] (0.945,0) -- (0.945,0.16);
  \draw[scred,line width=0.5pt] (0.959,0) -- (0.959,0.16);
  \draw[scred,line width=0.5pt] (0.973,0) -- (0.973,0.16);
  \draw[scred,line width=0.5pt] (0.987,0) -- (0.987,0.16);
  \draw[scred,line width=0.5pt] (1.001,0) -- (1.001,0.16);
  \draw[scred,line width=0.5pt] (1.015,0) -- (1.015,0.16);
  \draw[scblack] (0,0) rectangle (3.600,2.000);
  \draw[scblack] (0.000,0) -- (0.000,-0.08);
  \node[font=\scriptsize,below] at (0.000,-0.08) {563};
  \draw[scblack] (3.600,0) -- (3.600,-0.08);
  \node[font=\scriptsize,below] at (3.600,-0.08) {1599};
  \draw[scblack] (0,0.000) -- (-0.08,0.000);
  \node[font=\scriptsize,left] at (-0.10,0.000) {0};
  \draw[scblack] (0,0.317) -- (-0.08,0.317);
  \node[font=\scriptsize,left] at (-0.10,0.317) {$\theta_0$};
  \draw[scblack] (0,1.990) -- (-0.08,1.990);
  \node[font=\scriptsize,left] at (-0.10,1.990) {6.7};
  \node[font=\scriptsize,below=0.42cm] at (1.800,0) {sample index};
  \node[font={\scriptsize\bfseries},above] at (1.800,2.040) {Y6};
  \end{scope}
  \begin{scope}[xshift=0.00cm,yshift=-6.30cm]
  \fill[scevent,opacity=0.14] (2.561,0) rectangle (2.908,2.000);
  \draw[scgray,solid,line width=0.4pt] (0.000,0.000) -- (0.014,0.272) -- (0.028,0.272) -- (0.042,0.000) -- (0.056,0.000) -- (0.069,0.272) -- (0.083,0.272) -- (0.097,0.000) -- (0.111,0.000) -- (0.125,0.272) -- (0.139,0.272) -- (0.153,0.000) -- (0.167,0.000) -- (0.181,0.272) -- (0.195,0.272) -- (0.208,0.000) -- (0.222,0.000) -- (0.236,0.272) -- (0.250,0.272) -- (0.264,0.000) -- (0.278,0.000) -- (0.292,0.272) -- (0.306,0.272) -- (0.320,0.000) -- (0.334,0.000) -- (0.347,0.272) -- (0.361,0.272) -- (0.375,0.000) -- (0.389,0.000) -- (0.403,0.272) -- (0.417,0.272) -- (0.431,0.000) -- (0.445,0.000) -- (0.459,0.272) -- (0.473,0.272) -- (0.486,0.000) -- (0.500,0.000) -- (0.514,0.272) -- (0.528,0.272) -- (0.542,0.000) -- (0.556,0.000) -- (0.570,0.272) -- (0.584,0.272) -- (0.598,0.000) -- (0.612,0.000) -- (0.625,0.272) -- (0.639,0.272) -- (0.653,0.000) -- (0.667,0.000) -- (0.681,0.272) -- (0.695,0.272) -- (0.709,0.000) -- (0.723,0.000) -- (0.737,0.272) -- (0.751,0.272) -- (0.764,0.000) -- (0.778,0.000) -- (0.792,0.272) -- (0.806,0.272) -- (0.820,0.000) -- (0.834,0.544) -- (0.848,0.816) -- (0.862,0.272) -- (0.876,0.544) -- (0.890,0.544) -- (0.903,0.816) -- (0.917,0.272) -- (0.931,0.544) -- (0.945,0.544) -- (0.959,0.816) -- (0.973,0.272) -- (0.987,0.544) -- (1.001,0.544) -- (1.015,0.816) -- (1.029,0.272) -- (1.042,0.544) -- (1.056,0.544) -- (1.070,0.816) -- (1.084,0.272) -- (1.098,0.544) -- (1.112,0.544) -- (1.126,0.816) -- (1.140,0.272) -- (1.154,0.544) -- (1.168,0.544) -- (1.181,0.816) -- (1.195,0.272) -- (1.209,0.544) -- (1.223,0.544) -- (1.237,0.816) -- (1.251,0.272) -- (1.265,0.544) -- (1.279,0.544) -- (1.293,0.816) -- (1.307,0.272) -- (1.320,0.544) -- (1.334,0.544) -- (1.348,0.816) -- (1.362,0.272) -- (1.376,0.544) -- (1.390,0.544) -- (1.404,0.816) -- (1.418,0.272) -- (1.432,0.544) -- (1.446,0.544) -- (1.459,0.816) -- (1.473,0.272) -- (1.487,0.544) -- (1.501,0.544) -- (1.515,0.816) -- (1.529,0.272) -- (1.543,0.544) -- (1.557,0.544) -- (1.571,0.816) -- (1.585,0.272) -- (1.598,0.544) -- (1.612,0.544) -- (1.626,0.816) -- (1.640,0.272) -- (1.654,0.544) -- (1.668,0.544) -- (1.682,0.816) -- (1.696,0.272) -- (1.710,0.544) -- (1.724,0.544) -- (1.737,0.816) -- (1.751,0.272) -- (1.765,0.544) -- (1.779,0.544) -- (1.793,0.816) -- (1.807,0.272) -- (1.821,0.544) -- (1.835,0.544) -- (1.849,0.816) -- (1.863,0.272) -- (1.876,0.000) -- (1.890,0.000) -- (1.904,0.272) -- (1.918,0.272) -- (1.932,0.000) -- (1.946,0.000) -- (1.960,0.272) -- (1.974,0.272) -- (1.988,0.000) -- (2.002,0.000) -- (2.015,0.272) -- (2.029,0.272) -- (2.043,0.000) -- (2.057,0.000) -- (2.071,0.272) -- (2.085,0.272) -- (2.099,0.000) -- (2.113,0.000) -- (2.127,0.272) -- (2.141,0.272) -- (2.154,0.000) -- (2.168,0.000) -- (2.182,0.272) -- (2.196,0.272) -- (2.210,0.000) -- (2.224,0.000) -- (2.238,0.272) -- (2.252,0.272) -- (2.266,0.000) -- (2.280,0.000) -- (2.293,0.272) -- (2.307,0.272) -- (2.321,0.000) -- (2.335,0.000) -- (2.349,0.272) -- (2.363,0.272) -- (2.377,0.000) -- (2.391,0.000) -- (2.405,0.272) -- (2.419,0.272) -- (2.432,0.000) -- (2.446,0.000) -- (2.460,0.272) -- (2.474,0.272) -- (2.488,0.000) -- (2.502,0.000) -- (2.516,0.272) -- (2.530,0.272) -- (2.544,0.000) -- (2.558,0.000) -- (2.571,1.905) -- (2.585,0.272) -- (2.599,1.633) -- (2.613,0.000) -- (2.627,1.905) -- (2.641,0.272) -- (2.655,1.633) -- (2.669,0.000) -- (2.683,1.905) -- (2.697,0.272) -- (2.710,1.633) -- (2.724,0.000) -- (2.738,1.905) -- (2.752,0.272) -- (2.766,1.633) -- (2.780,0.000) -- (2.794,1.905) -- (2.808,0.272) -- (2.822,1.633) -- (2.836,0.000) -- (2.849,1.905) -- (2.863,0.272) -- (2.877,1.633) -- (2.891,0.000) -- (2.905,1.905) -- (2.919,0.272) -- (2.933,0.000) -- (2.947,0.000) -- (2.961,0.272) -- (2.975,0.272) -- (2.988,0.000) -- (3.002,0.000) -- (3.016,0.272) -- (3.030,0.272) -- (3.044,0.000) -- (3.058,0.000) -- (3.072,0.272) -- (3.086,0.272) -- (3.100,0.000) -- (3.114,0.000) -- (3.127,0.272) -- (3.141,0.272) -- (3.155,0.000) -- (3.169,0.000) -- (3.183,0.272) -- (3.197,0.272) -- (3.211,0.000) -- (3.225,0.000) -- (3.239,0.272) -- (3.253,0.272) -- (3.266,0.000) -- (3.280,0.000) -- (3.294,0.272) -- (3.308,0.272) -- (3.322,0.000) -- (3.336,0.000) -- (3.350,0.272) -- (3.364,0.272) -- (3.378,0.000) -- (3.392,0.000) -- (3.405,0.272) -- (3.419,0.272) -- (3.433,0.000) -- (3.447,0.000) -- (3.461,0.272) -- (3.475,0.272) -- (3.489,0.000) -- (3.503,0.000) -- (3.517,0.272) -- (3.531,0.272) -- (3.544,0.000) -- (3.558,0.000) -- (3.572,0.272) -- (3.586,0.272) -- (3.600,0.000);
  \draw[scblack,solid,line width=0.7pt] (0,0.272) -- (3.600,0.272);
  \draw[scblue,dotted,line width=0.7pt] (0.000,0.272) -- (0.014,0.000) -- (0.028,0.272) -- (0.042,0.272) -- (0.056,0.272) -- (0.069,0.272) -- (0.083,0.272) -- (0.097,0.272) -- (0.111,0.272) -- (0.125,0.272) -- (0.139,0.272) -- (0.153,0.272) -- (0.167,0.272) -- (0.181,0.272) -- (0.195,0.272) -- (0.208,0.272) -- (0.222,0.272) -- (0.236,0.272) -- (0.250,0.272) -- (0.264,0.272) -- (0.278,0.272) -- (0.292,0.272) -- (0.306,0.272) -- (0.320,0.272) -- (0.334,0.272) -- (0.347,0.272) -- (0.361,0.272) -- (0.375,0.272) -- (0.389,0.272) -- (0.403,0.272) -- (0.417,0.272) -- (0.431,0.272) -- (0.445,0.272) -- (0.459,0.272) -- (0.473,0.272) -- (0.486,0.272) -- (0.500,0.272) -- (0.514,0.272) -- (0.528,0.272) -- (0.542,0.272) -- (0.556,0.272) -- (0.570,0.272) -- (0.584,0.272) -- (0.598,0.272) -- (0.612,0.272) -- (0.625,0.272) -- (0.639,0.272) -- (0.653,0.272) -- (0.667,0.272) -- (0.681,0.272) -- (0.695,0.272) -- (0.709,0.272) -- (0.723,0.272) -- (0.737,0.272) -- (0.751,0.272) -- (0.764,0.272) -- (0.778,0.272) -- (0.792,0.272) -- (0.806,0.272) -- (0.820,0.272) -- (0.834,0.272) -- (0.848,0.544) -- (0.862,0.816) -- (0.876,0.816) -- (0.890,0.816) -- (0.903,0.816) -- (0.917,0.816) -- (0.931,0.816) -- (0.945,0.816) -- (0.959,0.816) -- (0.973,0.816) -- (0.987,0.816) -- (1.001,0.816) -- (1.015,0.816) -- (1.029,0.816) -- (1.042,0.816) -- (1.056,0.816) -- (1.070,0.816) -- (1.084,0.816) -- (1.098,0.816) -- (1.112,0.816) -- (1.126,0.816) -- (1.140,0.816) -- (1.154,0.816) -- (1.168,0.816) -- (1.181,0.816) -- (1.195,0.816) -- (1.209,0.816) -- (1.223,0.816) -- (1.237,0.816) -- (1.251,0.816) -- (1.265,0.816) -- (1.279,0.816) -- (1.293,0.816) -- (1.307,0.816) -- (1.320,0.816) -- (1.334,0.816) -- (1.348,0.816) -- (1.362,0.816) -- (1.376,0.816) -- (1.390,0.816) -- (1.404,0.816) -- (1.418,0.816) -- (1.432,0.816) -- (1.446,0.816) -- (1.459,0.816) -- (1.473,0.816) -- (1.487,0.816) -- (1.501,0.816) -- (1.515,0.816) -- (1.529,0.816) -- (1.543,0.816) -- (1.557,0.816) -- (1.571,0.816) -- (1.585,0.816) -- (1.598,0.816) -- (1.612,0.816) -- (1.626,0.816) -- (1.640,0.816) -- (1.654,0.816) -- (1.668,0.816) -- (1.682,0.816) -- (1.696,0.816) -- (1.710,0.816) -- (1.724,0.816) -- (1.737,0.816) -- (1.751,0.816) -- (1.765,0.816) -- (1.779,0.816) -- (1.793,0.816) -- (1.807,0.816) -- (1.821,0.816) -- (1.835,0.816) -- (1.849,0.816) -- (1.863,0.816) -- (1.876,0.816) -- (1.890,0.816) -- (1.904,0.816) -- (1.918,0.816) -- (1.932,0.816) -- (1.946,0.816) -- (1.960,0.816) -- (1.974,0.816) -- (1.988,0.816) -- (2.002,0.272) -- (2.015,0.272) -- (2.029,0.272) -- (2.043,0.272) -- (2.057,0.272) -- (2.071,0.272) -- (2.085,0.272) -- (2.099,0.272) -- (2.113,0.272) -- (2.127,0.272) -- (2.141,0.272) -- (2.154,0.272) -- (2.168,0.272) -- (2.182,0.272) -- (2.196,0.272) -- (2.210,0.272) -- (2.224,0.272) -- (2.238,0.272) -- (2.252,0.272) -- (2.266,0.272) -- (2.280,0.272) -- (2.293,0.272) -- (2.307,0.272) -- (2.321,0.272) -- (2.335,0.272) -- (2.349,0.272) -- (2.363,0.272) -- (2.377,0.272) -- (2.391,0.272) -- (2.405,0.272) -- (2.419,0.272) -- (2.432,0.272) -- (2.446,0.272) -- (2.460,0.272) -- (2.474,0.272) -- (2.488,0.272) -- (2.502,0.272) -- (2.516,0.272) -- (2.530,0.272) -- (2.544,0.272) -- (2.558,0.272) -- (2.571,0.272) -- (2.585,1.905) -- (2.599,1.905) -- (2.613,1.905) -- (2.627,1.905) -- (2.641,1.905) -- (2.655,1.905) -- (2.669,1.905) -- (2.683,1.905) -- (2.697,1.905) -- (2.710,1.905) -- (2.724,1.905) -- (2.738,1.905) -- (2.752,1.905) -- (2.766,1.905) -- (2.780,1.905) -- (2.794,1.905) -- (2.808,1.905) -- (2.822,1.905) -- (2.836,1.905) -- (2.849,1.905) -- (2.863,1.905) -- (2.877,1.905) -- (2.891,1.905) -- (2.905,1.905) -- (2.919,1.905) -- (2.933,1.905) -- (2.947,1.905) -- (2.961,1.905) -- (2.975,1.905) -- (2.988,1.905) -- (3.002,1.905) -- (3.016,1.905) -- (3.030,1.905) -- (3.044,1.905) -- (3.058,0.272) -- (3.072,0.272) -- (3.086,0.272) -- (3.100,0.272) -- (3.114,0.272) -- (3.127,0.272) -- (3.141,0.272) -- (3.155,0.272) -- (3.169,0.272) -- (3.183,0.272) -- (3.197,0.272) -- (3.211,0.272) -- (3.225,0.272) -- (3.239,0.272) -- (3.253,0.272) -- (3.266,0.272) -- (3.280,0.272) -- (3.294,0.272) -- (3.308,0.272) -- (3.322,0.272) -- (3.336,0.272) -- (3.350,0.272) -- (3.364,0.272) -- (3.378,0.272) -- (3.392,0.272) -- (3.405,0.272) -- (3.419,0.272) -- (3.433,0.272) -- (3.447,0.272) -- (3.461,0.272) -- (3.475,0.272) -- (3.489,0.272) -- (3.503,0.272) -- (3.517,0.272) -- (3.531,0.272) -- (3.544,0.272) -- (3.558,0.272) -- (3.572,0.272) -- (3.586,0.272) -- (3.600,0.272);
  \draw[scgreen,densely dashed,line width=1.0pt] (0.000,0.272) -- (0.014,0.272) -- (0.028,0.272) -- (0.042,0.272) -- (0.056,0.272) -- (0.069,0.272) -- (0.083,0.272) -- (0.097,0.272) -- (0.111,0.272) -- (0.125,0.272) -- (0.139,0.272) -- (0.153,0.272) -- (0.167,0.272) -- (0.181,0.272) -- (0.195,0.272) -- (0.208,0.272) -- (0.222,0.272) -- (0.236,0.272) -- (0.250,0.272) -- (0.264,0.272) -- (0.278,0.272) -- (0.292,0.272) -- (0.306,0.272) -- (0.320,0.272) -- (0.334,0.272) -- (0.347,0.272) -- (0.361,0.272) -- (0.375,0.272) -- (0.389,0.272) -- (0.403,0.272) -- (0.417,0.272) -- (0.431,0.272) -- (0.445,0.272) -- (0.459,0.272) -- (0.473,0.272) -- (0.486,0.272) -- (0.500,0.272) -- (0.514,0.272) -- (0.528,0.272) -- (0.542,0.272) -- (0.556,0.272) -- (0.570,0.272) -- (0.584,0.272) -- (0.598,0.272) -- (0.612,0.272) -- (0.625,0.272) -- (0.639,0.272) -- (0.653,0.272) -- (0.667,0.272) -- (0.681,0.272) -- (0.695,0.272) -- (0.709,0.272) -- (0.723,0.272) -- (0.737,0.272) -- (0.751,0.272) -- (0.764,0.272) -- (0.778,0.272) -- (0.792,0.272) -- (0.806,0.272) -- (0.820,0.272) -- (0.834,0.272) -- (0.848,0.272) -- (0.862,0.272) -- (0.876,0.272) -- (0.890,0.272) -- (0.903,0.816) -- (0.917,0.816) -- (0.931,0.816) -- (0.945,0.816) -- (0.959,0.816) -- (0.973,0.816) -- (0.987,0.816) -- (1.001,0.816) -- (1.015,0.816) -- (1.029,0.816) -- (1.042,0.816) -- (1.056,0.816) -- (1.070,0.816) -- (1.084,0.816) -- (1.098,0.816) -- (1.112,0.816) -- (1.126,0.816) -- (1.140,0.816) -- (1.154,0.816) -- (1.168,0.816) -- (1.181,0.816) -- (1.195,0.816) -- (1.209,0.816) -- (1.223,0.816) -- (1.237,0.816) -- (1.251,0.816) -- (1.265,0.816) -- (1.279,0.816) -- (1.293,0.816) -- (1.307,0.816) -- (1.320,0.816) -- (1.334,0.816) -- (1.348,0.816) -- (1.362,0.816) -- (1.376,0.816) -- (1.390,0.816) -- (1.404,0.816) -- (1.418,0.816) -- (1.432,0.816) -- (1.446,0.816) -- (1.459,0.816) -- (1.473,0.816) -- (1.487,0.816) -- (1.501,0.816) -- (1.515,0.816) -- (1.529,0.816) -- (1.543,0.816) -- (1.557,0.816) -- (1.571,0.816) -- (1.585,0.816) -- (1.598,0.816) -- (1.612,0.816) -- (1.626,0.816) -- (1.640,0.816) -- (1.654,0.816) -- (1.668,0.816) -- (1.682,0.816) -- (1.696,0.816) -- (1.710,0.816) -- (1.724,0.816) -- (1.737,0.816) -- (1.751,0.816) -- (1.765,0.816) -- (1.779,0.816) -- (1.793,0.816) -- (1.807,0.816) -- (1.821,0.816) -- (1.835,0.816) -- (1.849,0.816) -- (1.863,0.816) -- (1.876,0.816) -- (1.890,0.816) -- (1.904,0.816) -- (1.918,0.816) -- (1.932,0.816) -- (1.946,0.816) -- (1.960,0.816) -- (1.974,0.816) -- (1.988,0.272) -- (2.002,0.272) -- (2.015,0.272) -- (2.029,0.272) -- (2.043,0.272) -- (2.057,0.272) -- (2.071,0.272) -- (2.085,0.272) -- (2.099,0.272) -- (2.113,0.272) -- (2.127,0.272) -- (2.141,0.272) -- (2.154,0.272) -- (2.168,0.272) -- (2.182,0.272) -- (2.196,0.272) -- (2.210,0.272) -- (2.224,0.272) -- (2.238,0.272) -- (2.252,0.272) -- (2.266,0.272) -- (2.280,0.272) -- (2.293,0.272) -- (2.307,0.272) -- (2.321,0.272) -- (2.335,0.272) -- (2.349,0.272) -- (2.363,0.272) -- (2.377,0.272) -- (2.391,0.272) -- (2.405,0.272) -- (2.419,0.272) -- (2.432,0.272) -- (2.446,0.272) -- (2.460,0.272) -- (2.474,0.272) -- (2.488,0.272) -- (2.502,0.272) -- (2.516,0.272) -- (2.530,0.272) -- (2.544,0.272) -- (2.558,0.272) -- (2.571,0.272) -- (2.585,0.272) -- (2.599,0.272) -- (2.613,0.272) -- (2.627,0.272) -- (2.641,0.272) -- (2.655,0.272) -- (2.669,0.272) -- (2.683,0.272) -- (2.697,0.272) -- (2.710,0.272) -- (2.724,0.272) -- (2.738,0.272) -- (2.752,0.272) -- (2.766,0.272) -- (2.780,0.272) -- (2.794,0.272) -- (2.808,0.272) -- (2.822,0.272) -- (2.836,0.272) -- (2.849,0.272) -- (2.863,0.272) -- (2.877,0.272) -- (2.891,0.272) -- (2.905,0.272) -- (2.919,0.272) -- (2.933,0.272) -- (2.947,0.272) -- (2.961,0.272) -- (2.975,0.272) -- (2.988,0.272) -- (3.002,0.272) -- (3.016,0.272) -- (3.030,0.272) -- (3.044,0.272) -- (3.058,0.272) -- (3.072,0.272) -- (3.086,0.272) -- (3.100,0.272) -- (3.114,0.272) -- (3.127,0.272) -- (3.141,0.272) -- (3.155,0.272) -- (3.169,0.272) -- (3.183,0.272) -- (3.197,0.272) -- (3.211,0.272) -- (3.225,0.272) -- (3.239,0.272) -- (3.253,0.272) -- (3.266,0.272) -- (3.280,0.272) -- (3.294,0.272) -- (3.308,0.272) -- (3.322,0.272) -- (3.336,0.272) -- (3.350,0.272) -- (3.364,0.272) -- (3.378,0.272) -- (3.392,0.272) -- (3.405,0.272) -- (3.419,0.272) -- (3.433,0.272) -- (3.447,0.272) -- (3.461,0.272) -- (3.475,0.272) -- (3.489,0.272) -- (3.503,0.272) -- (3.517,0.272) -- (3.531,0.272) -- (3.544,0.272) -- (3.558,0.272) -- (3.572,0.272) -- (3.586,0.272) -- (3.600,0.272);
  \node[scgreen,font=\scriptsize] at (0.903,2.020) {$\bullet$};
  \node[scgreen,font=\scriptsize] at (1.988,2.020) {$\bullet$};
  \draw[scred,line width=0.5pt] (2.683,0) -- (2.683,0.16);
  \draw[scred,line width=0.5pt] (2.697,0) -- (2.697,0.16);
  \draw[scred,line width=0.5pt] (2.710,0) -- (2.710,0.16);
  \draw[scred,line width=0.5pt] (2.724,0) -- (2.724,0.16);
  \draw[scred,line width=0.5pt] (2.738,0) -- (2.738,0.16);
  \draw[scred,line width=0.5pt] (2.752,0) -- (2.752,0.16);
  \draw[scred,line width=0.5pt] (2.766,0) -- (2.766,0.16);
  \draw[scred,line width=0.5pt] (2.780,0) -- (2.780,0.16);
  \draw[scred,line width=0.5pt] (2.794,0) -- (2.794,0.16);
  \draw[scred,line width=0.5pt] (2.808,0) -- (2.808,0.16);
  \draw[scblack] (0,0) rectangle (3.600,2.000);
  \draw[scblack] (0.000,0) -- (0.000,-0.08);
  \node[font=\scriptsize,below] at (0.000,-0.08) {563};
  \draw[scblack] (3.600,0) -- (3.600,-0.08);
  \node[font=\scriptsize,below] at (3.600,-0.08) {1599};
  \draw[scblack] (0,0.000) -- (-0.08,0.000);
  \node[font=\scriptsize,left] at (-0.10,0.000) {0};
  \draw[scblack] (0,0.272) -- (-0.08,0.272);
  \node[font=\scriptsize,left] at (-0.10,0.272) {$\theta_0$};
  \draw[scblack] (0,2.011) -- (-0.08,2.011);
  \node[font=\scriptsize,left] at (-0.10,2.011) {7.9};
  \node[font=\scriptsize,below=0.42cm] at (1.800,0) {sample index};
  \node[font=\scriptsize,rotate=90] at (-0.720,1.000) {score};
  \node[font={\scriptsize\bfseries},above] at (1.800,2.040) {Y7};
  \end{scope}
  \begin{scope}[xshift=4.70cm,yshift=-6.30cm]
  \draw[scgray,solid,line width=0.9pt] (0.100,1.900) -- (0.600,1.900);
  \node[font=\scriptsize,right] at (0.650,1.900) {score};
  \draw[scblack,solid,line width=0.9pt] (0.100,1.560) -- (0.600,1.560);
  \node[font=\scriptsize,right] at (0.650,1.560) {fixed $\theta_0$};
  \draw[scblue,dotted,line width=0.9pt] (0.100,1.220) -- (0.600,1.220);
  \node[font=\scriptsize,right] at (0.650,1.220) {rolling (all)};
  \draw[scgreen,densely dashed,line width=0.9pt] (0.100,0.880) -- (0.600,0.880);
  \node[font=\scriptsize,right] at (0.650,0.880) {anchored $\theta$};
  \draw[scred,line width=0.6pt] (0.1,0.5) -- (0.6,0.5);
  \node[font=\scriptsize,right] at (0.65,0.5) {defer window};
  \node[scgreen,font=\scriptsize] at (0.35,0.14) {$\bullet$};
  \node[font=\scriptsize,right] at (0.65,0.14) {recalibration};
  \end{scope}
\end{tikzpicture}%
}
  \caption{Synthetic cases Y1--Y7: replay score (grey), fixed $\theta_0$ (black),
  admit-all rolling threshold (blue dotted) and the controller's threshold (green
  dashed), with its deferrals (red ticks), recalibrations (green dots) and labelled
  faults (shaded).}
  \label{fig:synthetic-map}
\end{figure}

\paragraph{Which guard component matters.}
Table~\ref{tab:ablation} removes one guard component at a time, on the same trace. The
cap $\kappa$ matters most on Y6, a step larger than the cap: the full controller cannot
follow it, so it defers 10 times, escalates, and keeps alerting (non-event alert
duration $757$); without the cap it follows the step and the duration falls to $16$.
Removing deferral costs most on the Y2 ramp, where non-event alert duration rises by
$373$. These are measured effects of fixed components, not tuning.

\begin{table}[tbp]
  \centering
  \caption{Change in each synthetic outcome when one guard component (cap $\kappa$,
  stability $\rho$ or deferral $D$) is removed, relative to the full controller. Only
  non-zero changes are shown.}
  \label{tab:ablation}
  \small
  \begin{tabular}{llrrrr}
\toprule
Removed guard & Scen. & $\Delta$non-ev. & $\Delta$defer & $\Delta$recal. & $\Delta$in-ev. \\
\midrule
cap $\kappa$ & Y3 & -3 & -10 & +2 & -0.130 \\
 & Y5 & -3 & -10 & +2 & -0.130 \\
 & Y6 & -741 & -9 & +1 & -- \\
 & Y7 & -3 & -10 & +2 & -0.130 \\
\midrule
stability $\rho$ & Y2 & +0 & -3 & +0 & -- \\
\midrule
deferral $D$ & Y2 & +373 & -3 & -1 & -- \\
 & Y3 & +0 & -10 & +0 & +0.200 \\
 & Y5 & +0 & -10 & +0 & +0.200 \\
 & Y6 & +40 & -10 & +0 & -- \\
 & Y7 & +0 & -10 & +0 & +0.200 \\
\bottomrule
\end{tabular}%

\end{table}

\subsection{Workload against coverage}\label{sec:results-tradeoff}
\Smeasured{} Figure~\ref{fig:tradeoff} places every arm by decision coverage and
non-event alert duration. On SKAB nearly all arms sit at full coverage, so only the
vertical spread of workload separates them. On the benign cases the Sun arm is quiet
only because it decides on $22/260$ (Y1) and $23/260$ (Y2) windows. Section~\ref{sec:limitations}
reads this trade-off against us.

\begin{figure}[tbp]
  \centering
  \resizebox{0.78\linewidth}{!}{\providecolor{scblack}{RGB}{0,0,0}%
\providecolor{scblue}{RGB}{31,78,150}%
\providecolor{scred}{RGB}{170,40,30}%
\providecolor{scgreen}{RGB}{20,110,60}%
\providecolor{scgray}{RGB}{120,120,120}%
\providecolor{scevent}{RGB}{60,60,60}%
\begin{tikzpicture}[>=stealth,line join=round,line cap=round]
  \begin{scope}[xshift=0.00cm,yshift=0.00cm]
  \node[scblack] at (4.049,0.000) {$\bullet$};
  \node[scblack] at (4.049,2.859) {$\bullet$};
  \node[scblack] at (4.049,0.262) {$\bullet$};
  \node[scblack] at (4.049,0.000) {$\bullet$};
  \node[scblack] at (4.049,0.000) {$\bullet$};
  \node[scblack] at (4.049,0.000) {$\bullet$};
  \node[scblack] at (4.049,0.000) {$\bullet$};
  \node[scblack] at (4.049,0.000) {$\bullet$};
  \node[scblack] at (4.049,0.000) {$\bullet$};
  \node[scblack] at (4.049,0.000) {$\bullet$};
  \node[scblack] at (0.286,0.000) {$\bullet$};
  \draw[scblack,line width=0.2pt] (0.290,0.000) -- (0.450,0.200);
  \node[font=\scriptsize,scblack,anchor=west,fill=white,fill opacity=0.9,text opacity=1,inner sep=0.5pt] at (0.450,0.200) {11};
  \draw[scblack,line width=0.2pt] (4.050,0.000) -- (3.770,0.420);
  \node[font=\scriptsize,scblack,anchor=east,fill=white,fill opacity=0.9,text opacity=1,inner sep=0.5pt] at (3.770,0.420) {1,4,5,6,7,8,9,10};
  \draw[scblack,line width=0.2pt] (4.050,2.860) -- (3.770,3.280);
  \node[font=\scriptsize,scblack,anchor=east,fill=white,fill opacity=0.9,text opacity=1,inner sep=0.5pt] at (3.770,3.280) {2};
  \draw[scblack,line width=0.2pt] (4.050,0.260) -- (3.770,0.720);
  \node[font=\scriptsize,scblack,anchor=east,fill=white,fill opacity=0.9,text opacity=1,inner sep=0.5pt] at (3.770,0.720) {3};
  \draw[scblack] (0,0) rectangle (4.500,3.400);
  \draw[scblack] (0.301,0) -- (0.301,-0.08);
  \node[font=\scriptsize,below] at (0.301,-0.08) {0.38};
  \draw[scblack] (4.049,0) -- (4.049,-0.08);
  \node[font=\scriptsize,below] at (4.049,-0.08) {1.00};
  \draw[scblack] (0,0.000) -- (-0.08,0.000);
  \node[font=\scriptsize,left] at (-0.10,0.000) {0};
  \draw[scblack] (0,1.705) -- (-0.08,1.705);
  \node[font=\scriptsize,left] at (-0.10,1.705) {65};
  \draw[scblack] (0,3.410) -- (-0.08,3.410);
  \node[font=\scriptsize,left] at (-0.10,3.410) {130};
  \node[font=\scriptsize,below=0.42cm] at (2.250,0) {decision coverage};
  \node[font=\scriptsize,rotate=90] at (-0.920,1.700) {non-event alert dur.\ (samples)};
  \node[font={\scriptsize\bfseries},above] at (2.250,3.440) {valve1};
  \end{scope}
  \begin{scope}[xshift=5.70cm,yshift=0.00cm]
  \node[scblack] at (4.049,0.550) {$\bullet$};
  \node[scblack] at (4.049,2.858) {$\bullet$};
  \node[scblack] at (4.049,0.330) {$\bullet$};
  \node[scblack] at (4.049,0.000) {$\bullet$};
  \node[scblack] at (4.049,0.247) {$\bullet$};
  \node[scblack] at (4.049,0.550) {$\bullet$};
  \node[scblack] at (4.049,0.550) {$\bullet$};
  \node[scblack] at (4.049,0.550) {$\bullet$};
  \node[scblack] at (4.049,0.550) {$\bullet$};
  \node[scblack] at (4.049,0.550) {$\bullet$};
  \node[scblack] at (0.289,0.110) {$\bullet$};
  \draw[scblack,line width=0.2pt] (0.290,0.110) -- (0.450,0.270);
  \node[font=\scriptsize,scblack,anchor=west,fill=white,fill opacity=0.9,text opacity=1,inner sep=0.5pt] at (0.450,0.270) {11};
  \draw[scblack,line width=0.2pt] (4.050,0.550) -- (3.770,1.320);
  \node[font=\scriptsize,scblack,anchor=east,fill=white,fill opacity=0.9,text opacity=1,inner sep=0.5pt] at (3.770,1.320) {1,6,7,8,9,10};
  \draw[scblack,line width=0.2pt] (4.050,2.860) -- (3.770,3.280);
  \node[font=\scriptsize,scblack,anchor=east,fill=white,fill opacity=0.9,text opacity=1,inner sep=0.5pt] at (3.770,3.280) {2};
  \draw[scblack,line width=0.2pt] (4.050,0.330) -- (3.770,1.020);
  \node[font=\scriptsize,scblack,anchor=east,fill=white,fill opacity=0.9,text opacity=1,inner sep=0.5pt] at (3.770,1.020) {3};
  \draw[scblack,line width=0.2pt] (4.050,0.000) -- (3.770,0.420);
  \node[font=\scriptsize,scblack,anchor=east,fill=white,fill opacity=0.9,text opacity=1,inner sep=0.5pt] at (3.770,0.420) {4};
  \draw[scblack,line width=0.2pt] (4.050,0.250) -- (3.770,0.720);
  \node[font=\scriptsize,scblack,anchor=east,fill=white,fill opacity=0.9,text opacity=1,inner sep=0.5pt] at (3.770,0.720) {5};
  \draw[scblack] (0,0) rectangle (4.500,3.400);
  \draw[scblack] (0.307,0) -- (0.307,-0.08);
  \node[font=\scriptsize,below] at (0.307,-0.08) {0.41};
  \draw[scblack] (4.049,0) -- (4.049,-0.08);
  \node[font=\scriptsize,below] at (4.049,-0.08) {1.00};
  \draw[scblack] (0,0.000) -- (-0.08,0.000);
  \node[font=\scriptsize,left] at (-0.10,0.000) {0};
  \draw[scblack] (0,1.704) -- (-0.08,1.704);
  \node[font=\scriptsize,left] at (-0.10,1.704) {62};
  \draw[scblack] (0,3.408) -- (-0.08,3.408);
  \node[font=\scriptsize,left] at (-0.10,3.408) {124};
  \node[font=\scriptsize,below=0.42cm] at (2.250,0) {decision coverage};
  \node[font=\scriptsize,rotate=90] at (-0.920,1.700) {non-event alert dur.\ (samples)};
  \node[font={\scriptsize\bfseries},above] at (2.250,3.440) {valve2};
  \end{scope}
  \begin{scope}[xshift=0.00cm,yshift=-4.90cm]
  \node[scblack] at (4.047,1.874) {$\bullet$};
  \node[scblack] at (4.047,2.879) {$\bullet$};
  \node[scblack] at (4.047,0.038) {$\bullet$};
  \node[scblack] at (4.047,0.615) {$\bullet$};
  \node[scblack] at (4.047,2.464) {$\bullet$};
  \node[scblack] at (4.047,2.502) {$\bullet$};
  \node[scblack] at (4.047,0.050) {$\bullet$};
  \node[scblack] at (4.047,0.050) {$\bullet$};
  \node[scblack] at (4.047,0.050) {$\bullet$};
  \node[scblack] at (4.047,0.050) {$\bullet$};
  \node[scblack] at (0.268,0.025) {$\bullet$};
  \draw[scblack,line width=0.2pt] (0.270,0.030) -- (0.430,0.200);
  \node[font=\scriptsize,scblack,anchor=west,fill=white,fill opacity=0.9,text opacity=1,inner sep=0.5pt] at (0.430,0.200) {11};
  \draw[scblack,line width=0.2pt] (4.050,1.870) -- (3.770,2.290);
  \node[font=\scriptsize,scblack,anchor=east,fill=white,fill opacity=0.9,text opacity=1,inner sep=0.5pt] at (3.770,2.290) {1};
  \draw[scblack,line width=0.2pt] (4.050,2.880) -- (3.770,3.300);
  \node[font=\scriptsize,scblack,anchor=east,fill=white,fill opacity=0.9,text opacity=1,inner sep=0.5pt] at (3.770,3.300) {2};
  \draw[scblack,line width=0.2pt] (4.050,0.040) -- (3.770,0.460);
  \node[font=\scriptsize,scblack,anchor=east,fill=white,fill opacity=0.9,text opacity=1,inner sep=0.5pt] at (3.770,0.460) {3};
  \draw[scblack,line width=0.2pt] (4.050,0.620) -- (3.770,1.060);
  \node[font=\scriptsize,scblack,anchor=east,fill=white,fill opacity=0.9,text opacity=1,inner sep=0.5pt] at (3.770,1.060) {4};
  \draw[scblack,line width=0.2pt] (4.050,2.460) -- (3.770,2.700);
  \node[font=\scriptsize,scblack,anchor=east,fill=white,fill opacity=0.9,text opacity=1,inner sep=0.5pt] at (3.770,2.700) {5};
  \draw[scblack,line width=0.2pt] (4.050,2.500) -- (3.770,3.000);
  \node[font=\scriptsize,scblack,anchor=east,fill=white,fill opacity=0.9,text opacity=1,inner sep=0.5pt] at (3.770,3.000) {6};
  \draw[scblack,line width=0.2pt] (4.050,0.050) -- (3.770,0.760);
  \node[font=\scriptsize,scblack,anchor=east,fill=white,fill opacity=0.9,text opacity=1,inner sep=0.5pt] at (3.770,0.760) {7,8,9,10};
  \draw[scblack] (0,0) rectangle (4.500,3.400);
  \draw[scblack] (0.249,0) -- (0.249,-0.08);
  \node[font=\scriptsize,below] at (0.249,-0.08) {0.08};
  \draw[scblack] (4.047,0) -- (4.047,-0.08);
  \node[font=\scriptsize,below] at (4.047,-0.08) {1.00};
  \draw[scblack] (0,0.000) -- (-0.08,0.000);
  \node[font=\scriptsize,left] at (-0.10,0.000) {0};
  \draw[scblack] (0,1.698) -- (-0.08,1.698);
  \node[font=\scriptsize,left] at (-0.10,1.698) {541};
  \draw[scblack] (0,3.400) -- (-0.08,3.400);
  \node[font=\scriptsize,left] at (-0.10,3.400) {1083};
  \node[font=\scriptsize,below=0.42cm] at (2.250,0) {decision coverage};
  \node[font=\scriptsize,rotate=90] at (-0.920,1.700) {non-event alert dur.\ (samples)};
  \node[font={\scriptsize\bfseries},above] at (2.250,3.440) {Y1 (benign step)};
  \end{scope}
  \begin{scope}[xshift=5.70cm,yshift=-4.90cm]
  \node[scblack] at (4.047,1.811) {$\bullet$};
  \node[scblack] at (4.047,2.879) {$\bullet$};
  \node[scblack] at (4.047,0.088) {$\bullet$};
  \node[scblack] at (4.047,0.590) {$\bullet$};
  \node[scblack] at (4.047,2.351) {$\bullet$};
  \node[scblack] at (4.047,2.389) {$\bullet$};
  \node[scblack] at (3.999,0.640) {$\bullet$};
  \node[scblack] at (3.999,0.640) {$\bullet$};
  \node[scblack] at (4.047,0.640) {$\bullet$};
  \node[scblack] at (4.047,1.811) {$\bullet$};
  \node[scblack] at (0.268,0.075) {$\bullet$};
  \draw[scblack,line width=0.2pt] (0.270,0.080) -- (0.430,0.240);
  \node[font=\scriptsize,scblack,anchor=west,fill=white,fill opacity=0.9,text opacity=1,inner sep=0.5pt] at (0.430,0.240) {11};
  \draw[scblack,line width=0.2pt] (4.000,0.640) -- (3.720,1.310);
  \node[font=\scriptsize,scblack,anchor=east,fill=white,fill opacity=0.9,text opacity=1,inner sep=0.5pt] at (3.720,1.310) {7,8};
  \draw[scblack,line width=0.2pt] (4.050,1.810) -- (3.770,2.230);
  \node[font=\scriptsize,scblack,anchor=east,fill=white,fill opacity=0.9,text opacity=1,inner sep=0.5pt] at (3.770,2.230) {1,10};
  \draw[scblack,line width=0.2pt] (4.050,2.880) -- (3.770,3.300);
  \node[font=\scriptsize,scblack,anchor=east,fill=white,fill opacity=0.9,text opacity=1,inner sep=0.5pt] at (3.770,3.300) {2};
  \draw[scblack,line width=0.2pt] (4.050,0.090) -- (3.770,0.510);
  \node[font=\scriptsize,scblack,anchor=east,fill=white,fill opacity=0.9,text opacity=1,inner sep=0.5pt] at (3.770,0.510) {3};
  \draw[scblack,line width=0.2pt] (4.050,0.590) -- (3.770,1.010);
  \node[font=\scriptsize,scblack,anchor=east,fill=white,fill opacity=0.9,text opacity=1,inner sep=0.5pt] at (3.770,1.010) {4};
  \draw[scblack,line width=0.2pt] (4.050,2.350) -- (3.770,2.700);
  \node[font=\scriptsize,scblack,anchor=east,fill=white,fill opacity=0.9,text opacity=1,inner sep=0.5pt] at (3.770,2.700) {5};
  \draw[scblack,line width=0.2pt] (4.050,2.390) -- (3.770,3.000);
  \node[font=\scriptsize,scblack,anchor=east,fill=white,fill opacity=0.9,text opacity=1,inner sep=0.5pt] at (3.770,3.000) {6};
  \draw[scblack,line width=0.2pt] (4.050,0.640) -- (3.770,1.610);
  \node[font=\scriptsize,scblack,anchor=east,fill=white,fill opacity=0.9,text opacity=1,inner sep=0.5pt] at (3.770,1.610) {9};
  \draw[scblack] (0,0) rectangle (4.500,3.400);
  \draw[scblack] (0.275,0) -- (0.275,-0.08);
  \node[font=\scriptsize,below] at (0.275,-0.08) {0.09};
  \draw[scblack] (4.047,0) -- (4.047,-0.08);
  \node[font=\scriptsize,below] at (4.047,-0.08) {1.00};
  \draw[scblack] (0,0.000) -- (-0.08,0.000);
  \node[font=\scriptsize,left] at (-0.10,0.000) {0};
  \draw[scblack] (0,1.698) -- (-0.08,1.698);
  \node[font=\scriptsize,left] at (-0.10,1.698) {541};
  \draw[scblack] (0,3.400) -- (-0.08,3.400);
  \node[font=\scriptsize,left] at (-0.10,3.400) {1083};
  \node[font=\scriptsize,below=0.42cm] at (2.250,0) {decision coverage};
  \node[font=\scriptsize,rotate=90] at (-0.920,1.700) {non-event alert dur.\ (samples)};
  \node[font={\scriptsize\bfseries},above] at (2.250,3.440) {Y2 (benign ramp)};
  \end{scope}
  \node[font=\scriptsize,scblack,anchor=north west,align=left] at (11.200,0.200) {1~fixed\\2~rolling\\3~rolling all\\4~k consecutive\\5~m of n\\6~hysteresis\\7~a-recalibration\\8~a-no cap\\9~a-no stability\\10~a-no defer\\11~sun confidence sequence};
\end{tikzpicture}%
}
  \caption{Decision coverage ($x$) against non-event alert duration in samples ($y$),
  one numbered marker per arm (key at right; coincident arms share a label). Arms to the
  bottom right are quiet \emph{and} decide on every window. Y1 and Y2 contain no fault
  event.}
  \label{fig:tradeoff}
\end{figure}
